%==============================================================================
%  577deputes.fr — Note méthodologique
%  Formulation des indicateurs, analyse de sensibilité, rapprochement à la
%  littérature, limites.
%
%  Compilation : pdfLaTeX (Overleaf : "pdfLaTeX", aucune dépendance exotique).
%  Trilingue : Partie I français · Partie II English · Partie III español.
%  Les parties EN et ES sont, à ce stade, des squelettes (%==============================================================================
%  Part II — English version (full translation, mirrors Part I).
%==============================================================================
\clearpage
\selectlanguage{english}
\part{English version}
\setcounter{section}{0}   % each language part is numbered 1, 2, 3, … on its own

\section{Introduction and scope}

\subsection{Object}
\emph{577deputes.fr} is a \emph{local-first} web application: it downloads the
official open datasets, ingests them into a SQLite database, and exposes browsing
pages and derived indicators for the 17\textsuperscript{th} legislature of the
French National Assembly (ongoing since July~2024). The application produces no
primary data; everything it displays is either raw data taken as-is from a public
source, or the result of a computation documented here.

\subsection{Sources}
\begin{itemize}
  \item \textbf{National Assembly} (\href{https://data.assemblee-nationale.fr}{data.assemblee-nationale.fr},
        \emph{Licence Ouverte~2.0}): actors, mandates and bodies; legislative files,
        documents and procedural acts; amendments; recorded (nominal) public
        ballots and individual votes; written, oral and questions to the
        Government; sitting agenda and verbatim records.
  \item \textbf{INSEE}: legal population by legislative constituency.
  \item \textbf{Ministry of the Interior} (data.gouv.fr): registered voters,
        turnout and abstention by constituency (2024 legislative elections).
\end{itemize}
The Assembly datasets are refreshed daily via conditional requests
(ETag / Last-Modified); a download failure does not affect the other sources.

\subsection{Volume (snapshot used for this document)}
\begin{center}\small
\begin{tabular}{@{}lr@{}}
\toprule
Deputies indexed (of which 577 currently sitting) & $\approx 630$ \\
Parliamentary questions & $\approx 17\,000$ \\
\quad of which with a published answer $>200$ characters & $10\,766$ \\
Amendments (17\textsuperscript{th} legislature) & $108\,115$ \\
Recorded public ballots & $6\,490$ \\
Individual votes & $1\,043\,305$ \\
(group, ballot) pairs with a majority position & $77\,748$ \\
\bottomrule
\end{tabular}
\end{center}

\subsection{What is out of scope}
\begin{itemize}
  \item \textbf{Show-of-hands votes} and \textbf{committee votes}: the
        \og{}Scrutins\fg{} dataset only contains recorded nominal ballots held in
        the plenary chamber. All cohesion, discipline and absenteeism figures rest
        on that perimeter (cf.~§\ref{sec:limites-en}).
  \item \textbf{Files opened under the 16\textsuperscript{th} legislature} and
        continued under the 17\textsuperscript{th}: they are indexed but attached
        to their legislature of origin, hence small discrepancies between counters.
  \item \textbf{Intent inference}: no indicator in this document claims to
        establish a motive. \og{}Mass filing\fg{}, \og{}phantom amendment\fg{},
        \og{}strategic absenteeism\fg{} describe \emph{measurable regularities}, not
        designs.
\end{itemize}

\section{Notation and conventions}\label{sec:notations-en}
We write $J(A,B) = |A\cap B| / |A\cup B|$ for the Jaccard coefficient of two sets.
For a ballot $v$, $\mathrm{for}(v)$, $\mathrm{against}(v)$, $\mathrm{abst}(v)$ are
the official counts and $\text{cast}(v)=\mathrm{for}(v)+\mathrm{against}(v)+\mathrm{abst}(v)$.
For a parliamentary group $G$ and a ballot $v$, $p^\star(v,G)$ is the
\emph{majority position} of the group as published by the Assembly (field
\code{positionMajoritaire}), taking values in $\{\text{for},\text{against},\text{abstention}\}$.
On the snapshot used, this field is populated for all $77\,748$ pairs $(G,v)$ and
never takes any other value (47.9~\% \og{}for\fg{}, 46.8~\% \og{}against\fg{},
5.3~\% \og{}abstention\fg{}): there is therefore, at our level, no exact tie to
arbitrate.

Each indicator subsection follows the same plan: \textbf{(a)~definition},
\textbf{(b)~algorithm \& parameters}, \textbf{(c)~sensitivity}, \textbf{(d)~literature},
\textbf{(e)~limitations}. The sensitivity figures come from the script
\code{methodo\_sensitivity.py} run on the snapshot above; the full tables are in
Appendix~\ref{sec:annexes-en} and in the deposit.

\section{The indicators}

\subsection{Near-duplicate amendment detection}\label{sec:clusters-en}

\paragraph{(a) Definition.} We seek to group amendments whose operative text is
\og{}essentially the same\fg{}, whether or not they share an author. For an
amendment $a$ we build $S_a$, the set of word $k$-\emph{shingles} (sequences of
$k$ consecutive words) of the normalised text, each hashed to 64~bits. Two
amendments $a,a'$ are deemed near-identical if $J(S_a,S_{a'})\ge\theta$. A
\emph{cluster} is a connected component (under the union of these pairs) of size
$\ge 2$.

\paragraph{(b) Algorithm \& parameters.}
\begin{enumerate}[label=\arabic*.]
  \item \emph{Normalisation}: strip HTML tags; Unicode decomposition (NFKD) and
        removal of diacritics; lowercasing; removal of any character outside
        \texttt{[a-z0-9 ]}; whitespace collapse; truncation to $L_{\max}=12\,000$
        characters.
  \item \emph{Shingling}: $k = 5$; $S_a=\{h(w_i\cdots w_{i+k-1})\}$ with $h$ a
        64-bit BLAKE2b hash. Amendment discarded if $|S_a| < s_{\min}=8$.
  \item \emph{MinHash signature}: $K=64$ hash functions $g_j(x)=x\oplus\sigma_j$
        ($\sigma_j$ deterministic seeds), signature $\mathrm{sig}_j(a)=\min_{x\in S_a} g_j(x)$.
  \item \emph{LSH banding}: the signature is split into $b=16$ bands of $r=4$
        components; $a$ and $a'$ are \emph{candidates} if they coincide on at
        least one band. The probability that a pair of true similarity $s$ becomes
        a candidate is $P(s)=1-(1-s^{\,r})^{b}$ (S-curve, threshold
        $\approx (1/b)^{1/r}\approx 0.55$).
  \item \emph{Verification}: for each candidate pair, compute the exact Jaccard on
        $S_a$, $S_{a'}$; pair kept if $J\ge\theta=0.80$.
  \item \emph{Clustering}: union-find on the kept pairs; clusters of size $\ge 2$.
\end{enumerate}
On the snapshot: $62\,476$ amendments pass the $s_{\min}$ threshold; we obtain
$\approx 9\,460$ clusters grouping $28\,124$ amendments (the production database
records $28\,055$ — the $0.25~\%$ gap comes from a bucket-size guard introduced in
the verification script, and constitutes our reproducibility check).

\paragraph{(c) Sensitivity.}
\emph{Threshold $\theta$} (full set, production LSH):
\begin{center}\small
\begin{tabular}{rrrr}
\toprule
$\theta$ & pairs kept & clusters & clustered amendments (\%) \\
\midrule
0.70 & 180\,407 & 9\,638 & 32\,953 \;(52.8~\%) \\
0.75 & 147\,438 & 9\,334 & 30\,794 \;(49.3~\%) \\
0.78 & \phantom{0}88\,007 & 9\,649 & 29\,159 \;(46.7~\%) \\
\textbf{0.80} & \phantom{0}\textbf{60\,381} & \textbf{9\,461} & \textbf{28\,124 \;(45.0~\%)} \\
0.82 & \phantom{0}54\,380 & 9\,253 & 26\,834 \;(43.0~\%) \\
0.85 & \phantom{0}43\,701 & 9\,023 & 25\,212 \;(40.4~\%) \\
0.90 & \phantom{0}34\,480 & 8\,440 & 22\,782 \;(36.5~\%) \\
0.95 & \phantom{0}30\,679 & 7\,840 & 20\,934 \;(33.5~\%) \\
\bottomrule
\end{tabular}
\end{center}
Sensitivity is \emph{moderate}: $\pm0.05$ around $0.80$ shifts the number of
clustered amendments by $\approx\pm10$ to $\pm15~\%$. But the distribution of
candidate-pair similarities (Appendix~\ref{sec:annexes-en}) shows that $0.80$
falls in a \emph{transition region}: there is a large mass of pairs in
$[0.75;0.85)$ (amendments sharing much boilerplate without being copies), whereas
near-verbatim copies form a distinct peak above $0.95$. A stricter threshold
($0.90$) would keep only near-verbatim copies; a looser one ($0.70$--$0.75$) would
also absorb heavy rewrites. The choice of $0.80$ is a deliberate compromise, not a
\og{}natural\fg{} point of the cloud. For $\theta<0.80$, the recall of candidate
pairs by the 16$\times$4 LSH declines (see below): the displayed counts are then
lower bounds.

\emph{Shingle size $k$} (sample of $\approx 3\,500$ amendments, exact Jaccard on
all pairs, $\theta=0.80$): $k=3 \Rightarrow 12.4~\%$ clustered; $k=5 \Rightarrow
8.3~\%$; $k=7 \Rightarrow 6.5~\%$. Shorter shingles are more permissive; $k=5$ is
the median setting.

\emph{LSH configuration} ($b\times r$): $P(s)=1-(1-s^r)^b$. For the production
setting 16$\times$4: $P(0.80)=0.9998$, $P(0.70)=0.988$, $P(0.50)=0.64$. The LSH
only affects the \emph{recall} of candidate pairs; the final \emph{precision} is
guaranteed by the exact Jaccard verification that follows. At the production
threshold, recall is therefore near-perfect.

\emph{Truncation $L_{\max}$ and floor $s_{\min}$}: only $192$ amendments
($0.18~\%$) exceed $12\,000$ characters after normalisation; raising $s_{\min}$
from $8$ to $50$ would exclude $38~\%$ of the currently kept amendments (the value
$8$ is deliberately permissive, to keep short operative texts).

\paragraph{(d) Literature.} The shingle representation and the estimation of
Jaccard similarity by min-wise hashing are due to Broder~\cite{broder1997,broder2000};
band-based blocking (LSH) follows Gionis-Indyk-Motwani~\cite{gionis1999} and the
exposition of Leskovec-Rajaraman-Ullman~\cite{lru2014}. Application to detecting
reuse of legislative text is explored notably by the \emph{Legislative Influence
Detector}~\cite{burgess2016} and the work on the diffusion of \og{}model
bills\fg{} (Hertel-Fernandez~\cite{hf2014}, Jansa et al.~\cite{jansa2019},
Linder et al.~\cite{linder2020}). Our contribution is not methodological: it is
the \emph{application} of known building blocks to the Assembly's amendment stream,
with a descriptive typology.

\paragraph{(e) Limitations.} (i)~Normalisation \emph{truncates} beyond $12\,000$
characters: two very long amendments are compared only on their prefix.
(ii)~Word shingling ignores order beyond the window $k$: a reordering of
paragraphs can collapse the Jaccard. (iii)~The threshold $0.80$ is in the
ambiguous region of the cloud: the exact cluster boundary is conventional.
(iv)~A \og{}cluster\fg{} says nothing about intentional coordination: the same
template may circulate via an association federation, a union, an activist site,
without any concertation between deputies.

\subsection{Cluster typology}\label{sec:typo-en}

\paragraph{(a) Definition.} For a cluster $C$, let $n(C)=|C|$ be its size and
$g(C)$ the number of distinct parliamentary groups among its amendments. We
classify:
\[
\tau(C)=\begin{cases}
\text{mass filing} & \text{if } n(C)\ge \param{T_{masse}}=15,\\[2pt]
\text{inter-group convergence} & \text{otherwise if } g(C)\ge 2,\\[2pt]
\text{intra-group amplification} & \text{otherwise if } g(C)\le 1 \text{ and } n(C)\ge \param{T_{ampl}}=3,\\[2pt]
\text{simple reuse} & \text{otherwise.}
\end{cases}
\]
Size always takes precedence; no intent is inferred (the internal label
\og{}obstruction\fg{} was neutralised to \og{}mass filing\fg{}).

\paragraph{(b) Algorithm \& parameters.} Direct computation on the cluster table
($\param{T_{masse}}=15$, $\param{T_{ampl}}=3$). On the snapshot (clusters from
$\theta=0.80$): $91$ \og{}mass filing\fg{} clusters ($2\,578$ amendments),
$4\,314$ \og{}convergence\fg{} ($13\,834$), $773$ \og{}amplification\fg{}
($3\,146$), $4\,283$ \og{}reuse\fg{} ($8\,566$).

\paragraph{(c) Sensitivity.}
\begin{center}\small
\begin{tabular}{rr rr rr rr}
\toprule
$\param{T_{masse}}$ & $\param{T_{ampl}}$ & \multicolumn{2}{c}{mass filing} & \multicolumn{2}{c}{convergence} & \multicolumn{2}{c}{amplification} \\
 & & cl. & amdt & cl. & amdt & cl. & amdt \\
\midrule
10 & 3 & 203 & 3\,875 & 4\,226 & 12\,812 & 749 & 2\,871 \\
\textbf{15} & \textbf{3} & \textbf{91} & \textbf{2\,578} & \textbf{4\,314} & \textbf{13\,834} & \textbf{773} & \textbf{3\,146} \\
20 & 3 & 43 & 1\,776 & 4\,352 & 14\,472 & 783 & 3\,310 \\
30 & 3 & 19 & 1\,218 & 4\,370 & 14\,880 & 789 & 3\,460 \\
50 & 3 & \phantom{0}7 & \phantom{0,}784 & 4\,377 & 15\,121 & 794 & 3\,653 \\
\midrule
15 & 2 & 91 & 2\,578 & 4\,314 & 13\,834 & 5\,056 & 11\,712 \\
15 & 5 & 91 & 2\,578 & 4\,314 & 13\,834 & 160 & 1\,138 \\
\bottomrule
\end{tabular}
\end{center}
The \emph{number} of \og{}mass filing\fg{} clusters is sensitive to
$\param{T_{masse}}$ ($91$ down to $7$ between thresholds $15$ and $50$) but the
\emph{volume} of amendments affected is less so; the $\param{T_{ampl}}$ threshold
mostly shifts clusters between \og{}amplification\fg{} and \og{}reuse\fg{}. The
size distribution is in Appendix~\ref{sec:annexes-en}.

\paragraph{(d) Literature.} The typology is ad~hoc and descriptive; it belongs to
the field studying \emph{parliamentary obstruction} and mass amendment filing as a
debate tactic (cf.~Koger~\cite{koger2010} for the U.S. case; the Assembly's own
Rules of Procedure regulate the \og{}joint discussion\fg{} of identical
amendments).

\paragraph{(e) Limitations.} Four boxes, round thresholds: a cluster of $14$
identical amendments and one of $16$ fall on opposite sides of a boundary with no
theoretical basis. The \og{}convergence\fg{} box does not distinguish a technical
consensus from an opportunistic alliance.

\subsection{Vote discipline (individual cohesion)}\label{sec:discipline-en}

\paragraph{(a) Definition.} Let $G_d$ be the group of deputy $d$. Let $V_d$ be the
set of ballots where $d$ cast a vote $p_d(v)\in\{\text{for},\text{against},\text{abstention}\}$
\emph{and} where $p^\star(v,G_d)$ is defined. We set
\[
\mathrm{cast}(d)=|V_d|,\quad
\mathrm{aligned}(d)=\big|\{v\in V_d : p_d(v)=p^\star(v,G_d)\}\big|,\quad
\mathrm{discipline}(d)=\frac{\mathrm{aligned}(d)}{\mathrm{cast}(d)}.
\]
Shown in rankings only if $\mathrm{cast}(d)\ge 50$ (or $\ge 100$ for the
\og{}dissidents\fg{} list, sorted by ascending discipline); unaffiliated (NI)
deputies are excluded.

\paragraph{(b) Algorithm \& parameters.} Everything rests on the official field
$p^\star$: we do \emph{not} compute a group's majority position ourselves. On the
snapshot, $\mathrm{discipline}$ has a mean of $\approx 91.9~\%$, a median of
$\approx 92.5~\%$, a bottom decile of $\approx 85~\%$.

\paragraph{(c) Sensitivity.} Effect of the cast-vote floor:
\begin{center}\small
\begin{tabular}{rrrrrr}
\toprule
floor & eligible & mean & median & $p_{10}$ & $p_{90}$ \\
\midrule
0   & 639 & 91.92 & 92.56 & 84.92 & 98.30 \\
20  & 628 & 91.87 & 92.50 & 84.93 & 98.25 \\
\textbf{50}  & \textbf{623} & \textbf{91.85} & \textbf{92.50} & \textbf{84.99} & \textbf{98.21} \\
100 & 611 & 91.83 & 92.45 & 85.02 & 98.18 \\
200 & 600 & 91.87 & 92.46 & 85.06 & 98.18 \\
\bottomrule
\end{tabular}
\end{center}
The indicator is \emph{very robust} to the floor: the mean and quantiles move by
only hundredths of a point between floors $0$ and $200$. The choice $\ge 50$ is
therefore not a sensitive parameter — it only avoids displaying a \og{}100~\%\fg{}
based on three votes.

\paragraph{(d) The \og{}circular bias\fg{}.} The position $p^\star(v,G_d)$ is the
plurality of the votes of $G_d$'s members, and $d$ is one of them: $d$ is thus
compared, in part, to a position $d$ helped define. This is a \emph{real} bias but
(i)~it is not of our making — it is the Assembly's official convention;
(ii)~for a group of $50$ to $120$ members, the weight of one individual vote in
the plurality is bounded by $\le 2~\%$; (iii)~the case of an exact tie, which would
make the plurality ambiguous, is arbitrated upstream by the Assembly and never
appears in the data (the field $p^\star$ is always populated and always in
$\{\text{for},\text{against},\text{abstention}\}$, cf.~§\ref{sec:notations-en}).

\paragraph{(e) Literature.} The discipline so defined is the \emph{individual}
(per-deputy) version of party cohesion indices: Rice index~\cite{rice1925},
the Agreement Index of Hix-Noury-Roland~\cite{hnr2005,hnr2007}, Carey's
discussion~\cite{carey2007} of \og{}competing principals\fg{}. It does \emph{not}
estimate ideal points: for that one would need NOMINATE-type or \emph{Optimal
Classification} methods (Poole-Rosenthal~\cite{pr1997}, Poole~\cite{poole2005}),
out of scope.

\paragraph{(e') Limitations.} (i)~Perimeter restricted to recorded nominal
ballots: the \og{}discipline\fg{} measured is that of \emph{cast and recorded}
votes, not group unity in general. (ii)~An abstention consistent with the group's
position counts as \og{}aligned\fg{}. (iii)~An often-absent deputy has a low
$\mathrm{cast}$: their discipline is measured on few cases (hence the floor).

\subsection{Inter-group cohesion matrix}\label{sec:matrice-en}

\paragraph{(a) Definition.} For two groups $G_1,G_2$, let $V_{12}$ be the set of
ballots where both have a defined majority position. The cohesion is
\[
\mathrm{coh}(G_1,G_2)=\frac{\big|\{v\in V_{12}:p^\star(v,G_1)=p^\star(v,G_2)\}\big|}{|V_{12}|},
\qquad \mathrm{coh}(G,G)=1.
\]

\paragraph{(b) Algorithm \& parameters.} No free parameter; a self-join on the
majority-position table.

\paragraph{(c) Sensitivity.} A debatable choice: \og{}abstention $=$ abstention\fg{}
counts as agreement. Recomputing the matrix without the ballots where both groups
abstain, the mean absolute difference is $0.21$~point and the maximum $0.83$~point.
The effect is therefore \emph{negligible}.

\paragraph{(d) Literature.} It is an \emph{agreement} matrix between blocs in the
sense of Hix-Noury-Roland~\cite{hnr2007}; it provides no one-dimensional ordering
(for that, see the spatial methods cited in §\ref{sec:discipline-en}).

\paragraph{(e) Limitations.} Cohesion computed on \emph{group positions}, not
individual votes: a group split $51/49$ and a unanimous group count the same. Two
\og{}against\fg{} votes for opposite reasons are indistinguishable from genuine
agreement.

\subsection{Ternary diagram of blocs}\label{sec:ternaire-en}

\paragraph{(a) Definition.} Given three \emph{anchor} groups $A_1,A_2,A_3$ placed
at the vertices of an equilateral triangle at positions $\mathbf{x}_{A_1},\mathbf{x}_{A_2},\mathbf{x}_{A_3}$,
a group $G$ with cohesions $c_i=\mathrm{coh}(G,A_i)$ is placed at
\[
(\alpha,\beta,\gamma)=\frac{(c_1,c_2,c_3)}{c_1+c_2+c_3},\qquad
\mathbf{x}_G=\alpha\,\mathbf{x}_{A_1}+\beta\,\mathbf{x}_{A_2}+\gamma\,\mathbf{x}_{A_3}.
\]
The default anchors are LFI, EPR, RN; the user can change them.

\paragraph{(b)--(c) Algorithm \& sensitivity.} The only \og{}parameter\fg{} is the
choice of the three anchors: the diagram is merely a \emph{re-projection} of the
cohesion matrix §\ref{sec:matrice-en} onto three chosen axes, and changes with
them.

\paragraph{(d)--(e) Literature \& limitations.} A visualisation tool, with no claim
of measurement: a group's position \emph{depends on the frame of reference}. Do
not read the distances as ideological distances.

\subsection{Coalitions by ballot (by topic)}\label{sec:coalitions-en}

\paragraph{(a) Definition.} We select the \og{}highly attended and divisive\fg{}
ballots: (a)~$\text{cast}(v)\ge \param{T_v}=300$ ($\approx 52~\%$ of the 577 seats);
(b)~$|\mathrm{for}(v)-\mathrm{against}(v)|\le \param{T_e}=50$. We sort by (b)
ascending then date descending and keep the top $N=10$; if fewer than $N/2$
ballots pass, we relax criterion (b). For each retained ballot, we display each
group's majority position.

\paragraph{(b)--(c) Parameters \& sensitivity.} Number of eligible ballots by
$(\param{T_v},\param{T_e})$:
\begin{center}\small
\begin{tabular}{r|rrrrr}
\toprule
$\param{T_v}\backslash\param{T_e}$ & 10 & 25 & 50 & 75 & 100 \\
\midrule
200 & \phantom{0}95 & 252 & 565 & 913 & 1\,200 \\
250 & \phantom{0}25 & \phantom{0}81 & 193 & 328 & \phantom{0,}440 \\
\textbf{300} & \phantom{0}\textbf{11} & \phantom{0}\textbf{36} & \phantom{00}\textbf{79} & 118 & \phantom{0,}155 \\
350 & \phantom{00}3 & \phantom{0}14 & \phantom{00}24 & \phantom{0}47 & \phantom{0,0}55 \\
400 & \phantom{00}2 & \phantom{00}7 & \phantom{00}11 & \phantom{0}19 & \phantom{0,0}20 \\
\bottomrule
\end{tabular}
\end{center}
At the production setting $(300,50)$, $79$ ballots pass: the fallback (b) is not
triggered. The selection is sensitive to both thresholds, as expected (it is a
filter, not an estimator).

\paragraph{(d)--(e) Literature \& limitations.} A purely descriptive approach: we
show \emph{who votes like whom} on close ballots, with no model. Major limitation:
a ballot \og{}close in votes\fg{} may be so because of asymmetric low attendance as
much as because of a real cleavage.

\subsection{\og{}Template\fg{} ministerial answers}\label{sec:templates-en}

\paragraph{(a) Definition.} For each written answer $q$ whose raw text exceeds
$200$ characters, we normalise (strip HTML, collapse whitespace), optionally remove
the leading salutation, lowercase: $t_q$. The answer is discarded if $|t_q|<100$.
The \emph{prefix key} is $\pi_L(q)=t_q[1:L]$ with $L=250$. A \og{}template
answer\fg{} is a prefix shared by at least $m=3$ distinct answers. The
\emph{template rate} is
\[
\rho \;=\; \frac{1}{|Q_{\text{elig}}|}\sum_{\pi \,:\, |\{q:\pi_L(q)=\pi\}|\ge m} \big|\{q:\pi_L(q)=\pi\}\big|.
\]
On the snapshot, $|Q_{\text{elig}}|=10\,766$ and $\rho=27.9~\%$.

\paragraph{(b)--(c) Parameters \& sensitivity.} Grid of the rate $\rho$ (in \%) by
prefix length $L$ and occurrence threshold $m$:
\begin{center}\small
\begin{tabular}{r|rrrr}
\toprule
$L\backslash m$ & 2 & \textbf{3} & 5 & 10 \\
\midrule
100  & 42.1 & 31.6 & 23.1 & 13.4 \\
\textbf{250}  & 37.0 & \textbf{27.9} & 20.3 & 11.6 \\
500  & 33.7 & 24.6 & 17.2 & \phantom{0}9.9 \\
1000 & 30.2 & 22.1 & 15.3 & \phantom{0}9.0 \\
\bottomrule
\end{tabular}
\end{center}
The rate is \emph{strongly sensitive to the threshold $m$} (ratio $\approx 3$--$4$
between $m=2$ and $m=10$) and \emph{moderately sensitive to $L$} (from $\approx
32~\%$ to $\approx 22~\%$ at fixed $m$). On this corpus, removing the salutation has
\emph{no} effect (the \og{}with\fg{} and \og{}without\fg{} removal variants
coincide exactly, the stored text not starting with such formulas). The value
shown on the site should therefore be read as \emph{one} point of a family of
measures, not as a constant: the honest phrasing is \og{}$\approx 28~\%$ of answers
share a $250$-character prefix with at least two others\fg{}.

\paragraph{(d) Literature.} Parliamentary questions as an oversight instrument, and
their answers as an object of study, are treated by Wiberg~\cite{wiberg1995},
Russo-Wiberg~\cite{rw2010}, Martin~\cite{martin2011}. Spotting standardised answers
by prefix grouping is a \emph{home-grown} heuristic, akin to boilerplate detection;
it is not a validated measure of \og{}officialese\fg{}.

\paragraph{(e) Limitations.} (i)~A common prefix is a \emph{necessary} but not
\emph{sufficient} condition for an \og{}empty\fg{} answer: two answers may share a
framing header then diverge entirely. (ii)~The threshold $m=3$ and the length
$L=250$ are arbitrary (cf.~grid). (iii)~The corpus is limited to published answers
longer than $200$ characters; unanswered questions do not enter the denominator of
$\rho$.

\subsection{Strategic absenteeism}\label{sec:absence-en}

\paragraph{(a) Definition.} Among ballots with $\text{cast}(v)\ge \param{T_v}=200$,
$v$ is said \emph{divisive} if $|\mathrm{for}(v)-\mathrm{against}(v)|/\text{cast}(v)<\param{c}=0.10$,
\emph{consensual} if $>\param{C}=0.50$, otherwise unclassified. For a deputy $d$,
let $\mathrm{pres}_{cli}(d)$ (resp.\ $\mathrm{pres}_{con}(d)$) be the \emph{share}
of divisive (resp.\ consensual) ballots where $d$ cast a vote. We set
$\mathrm{gap}(d)=\mathrm{pres}_{con}(d)-\mathrm{pres}_{cli}(d)$ (in points) and
$\mathrm{ratio}(d)=\mathrm{pres}_{con}(d)/\mathrm{pres}_{cli}(d)$. Shown if $d$ is
currently sitting and $\mathrm{tot}_{cli}(d)\ge \param{T_{cli}}=30$.

\paragraph{(b)--(c) Parameters \& sensitivity.} Four thresholds
($\param{T_v},\param{c},\param{C},\param{T_{cli}}$). The exact values for the
$4^4$ explored combinations are in Appendix~\ref{sec:annexes-en}; one observes that
the median of the \emph{gap} is close to $0$ — most deputies have a similar
attendance on divisive and consensual ballots — and that only outliers have a
marked \emph{gap}, whose amplitude varies from $\approx 13$ to $\approx 31$ points
depending on the thresholds.

\paragraph{(d)--(e) Literature \& limitations.} The idea relates to the literature
on legislative \emph{shirking} and selective participation in votes. Strong
limitations: (i)~not voting $\ne$ being absent — a present deputy may decline to
take part; (ii)~\og{}divisive\fg{} is defined by the \emph{outcome} of the vote,
not by its political salience \emph{a priori}; (iii)~a positive \emph{gap} has many
non-strategic causes (agenda, illness, distant constituency). A signal to explore,
not a verdict.

\subsection{\og{}Phantom\fg{} amendments}\label{sec:fantomes-en}

\paragraph{(a) Definition.} For first signatory $d$, over all their amendments:
\[
\begin{aligned}
\mathrm{defended}(d) &= \big|\{a : \mathrm{sort}(a)\in\{\text{Adopted, Rejected, Withdrawn, Discussed}\}\}\big|,\\[2pt]
\mathrm{phantom}(d) &= \big|\{a : \mathrm{sort}(a)=\text{Not defended, or empty}\}\big|,\\[2pt]
\mathrm{pct\_phantom}(d) &= \mathrm{phantom}(d)\,/\,\mathrm{total}(d).
\end{aligned}
\]
Shown if $\mathrm{total}(d)\ge \param{T_a}=50$; sorted by descending
$\mathrm{pct\_phantom}$.

\paragraph{(b)--(c) Parameters \& sensitivity.} Census of \og{}sort\fg{} codes
(share of all amendments): Rejected 22.9~\%; Inadmissible 16.3~\%; Pending 14.1~\%;
Adopted 13.0~\%; Withdrawn 10.2~\%; Discussed 9.0~\%; Lapsed 8.7~\%; Not defended
5.9~\%; Erased $\approx$0~\%. (No amendment has an empty/null sort on the snapshot.)
So $\mathrm{phantom}$ reduces in practice to \og{}Not defended\fg{} ($5.9~\%$),
$\mathrm{defended}$ covers $\approx 55~\%$, and the rest (Inadmissible, Pending,
Lapsed) is counted \emph{neither} as defended \emph{nor} as phantom. Effect of the
floor $\param{T_a}$:
\begin{center}\small
\begin{tabular}{rrrrr}
\toprule
$\param{T_a}$ & eligible authors & of which sitting & $\overline{\mathrm{pct\_phant.}}$ & $\mathrm{pct\_phant.}^{\max}$ \\
\midrule
10  & 558 & 513 & 7.4 & 61.5 \\
20  & 530 & 491 & 7.2 & 52.4 \\
\textbf{50}  & \textbf{439} & \textbf{412} & \textbf{6.5} & \textbf{50.0} \\
100 & 313 & 300 & 6.6 & 50.0 \\
200 & 167 & 163 & 5.8 & 34.7 \\
\bottomrule
\end{tabular}
\end{center}
The mean rate is stable ($\approx 6$--$7.5~\%$); the \emph{maximum} decreases with
the floor (extreme values are carried by low-volume authors). A \emph{continuous
view} (scatter plot $\mathrm{total}\times\mathrm{pct\_phantom}$, without a floor) is
provided in the deposit so as not to hide this dependence.

\paragraph{(d)--(e) Literature \& limitations.} No direct literature; we quantify an
asymmetry between \emph{filing} and \emph{defending} amendments. Limitations:
(i)~the \og{}Not defended\fg{} category also covers time pressure in the chamber,
independent of the author's will; (ii)~the choice to include \og{}Discussed\fg{}
in \og{}defended\fg{} and to exclude \og{}Lapsed\fg{} is debatable (the sensitivity
grid on the definition of \og{}defended\fg{} is in the deposit); (iii)~\og{}first
signatory\fg{} $\ne$ \og{}actual author\fg{} for collective amendments.

\subsection{Ministerial response delays and rankings}\label{sec:delais-en}

\paragraph{(a) Definition.} For a question $q$, $\mathrm{delay}(q)$ is the number of
days between filing and answer (undefined if unanswered). For a ministry (in the
sense of the short \og{}questioned ministry\fg{}) $\mu$:
\[
\overline{\mathrm{delay}}(\mu)=\operatorname*{mean}_{q:\,\mathrm{delay}(q)\ \text{defined}}\mathrm{delay}(q),\qquad
\mathrm{response\_rate}(\mu)=\frac{|\{q:\text{status}=\text{answered}\}|}{|\{q\ \text{addressed to }\mu\}|}.
\]
Shown if $|\{q\}|\ge \param{T_q}=30$ \emph{and} $|\{q:\text{status}=\text{answered}\}|\ge \param{T_r}=5$.
The response rate is shown alongside the delay: a short delay does not mean the same
thing if only $1$ question in $10$ received an answer.

\paragraph{(b)--(c) Parameters \& sensitivity.} Number of ministries retained by
$(\param{T_q},\param{T_r})$: $(10,1)\to 73$; $(20,*)\to 62$; $(30,*)\to 57$;
$(50,*)\to 52$; $(100,*)\to 40$. The $\param{T_r}$ threshold is almost inoperative
(a ministry receiving $\ge 30$ questions has almost always answered $\ge 5$).

\paragraph{(d)--(e) Literature \& limitations.} Cf.~§\ref{sec:templates-en} for
questions as an object of study. Limitations: the mean is sensitive to extreme
values (one answer at $700$~days drags the mean up); the \og{}questioned
ministry\fg{} may differ from the \og{}assigned ministry\fg{} that actually
answers.

\subsection{Simple derived indicators}\label{sec:simples-en}

\paragraph{(a) Cumulative activity.} $\mathrm{act}(d)=n_Q(d)+n_A(d)+n_V(d)$
(questions asked $+$ amendments filed $+$ votes cast). \emph{Structural warning}:
$n_V$ is an order of magnitude larger than $n_Q$ and $n_A$ (a diligent deputy casts
thousands of votes but files dozens of amendments and asks dozens of questions); the
cumulative index is therefore \emph{mechanically dominated by attendance at
ballots}. The site systematically accompanies it with the isolated sub-rankings
(questions only, amendments only, attendance only).

\paragraph{(b) Amendment adoption rate.} $\mathrm{adopted}(d)/\mathrm{total}(d)$
(per deputy or per group).

\paragraph{(c) Question response rate.} cf.~§\ref{sec:delais-en}.

\paragraph{(d) Age, seniority.} Derived from the date of birth and the mandate
start date. Trivial.

\paragraph{(e) Parametrisable ranking builder.} The site offers some thirty
\og{}metrics\fg{} combinable with optional filters (by group, period, text type,
etc.). They are all \emph{re-expositions} of the primitives above (counts, ratios,
means); there is no additional method.

\paragraph{(f) Constituency statistics.} $\mathrm{turnout}=\mathrm{voters}/\mathrm{registered}$,
population (INSEE), registered voters (Ministry of the Interior), joined by
$(\text{department code},\ \text{constituency number})$. Coverage on the snapshot:
$566$ constituencies with INSEE population (the $11$ constituencies for French
citizens abroad have none), $577$ with registered voters.

\section{Reproducibility}\label{sec:repro-en}
\begin{itemize}
  \item \textbf{Application code}: public, MIT-licensed
        (\href{https://github.com/RogericBot/577deputes}{github.com/RogericBot/577deputes}).
        The database and the photographs are not redistributed; they rebuild with
        the command \code{anqp bootstrap} (download and ingest the sources).
  \item \textbf{Derived caches}: the tables \code{deputy\_discipline\_cache},
        \code{groupe\_discipline\_cache}, \code{deputy\_amd\_cache},
        \code{groupe\_amd\_cache}, \code{scrutin\_groupes}, \code{amendement\_clusters},
        \code{circo\_stats} are recomputed at ingestion time; they contain nothing
        that does not derive from the sources.
  \item \textbf{Sensitivity analyses}: the script \code{scripts/methodo\_sensitivity.py}
        (standard library only, read-only) reproduces all the tables in this
        document; it is attached to the Zenodo deposit with its outputs (CSV and
        Markdown). All random draws use a fixed seed.
  \item \textbf{Snapshot}: the figures in this document correspond to a copy of the
        database taken for its writing (volume in §1.3); they may differ slightly
        from the current state of the site, which updates daily.
\end{itemize}

\section{General limitations and known biases}\label{sec:limites-en}
\begin{enumerate}[label=L\arabic*.]
  \item \textbf{Ballot perimeter.} Only recorded nominal ballots held in the
        plenary chamber exist in the data. Show-of-hands votes and committee votes
        are not there. Every vote indicator (discipline, cohesion, absenteeism,
        coalitions) inherits this limitation. Moreover, the Assembly very often
        votes with only a fraction of the chamber present (median $\approx 150$
        voters out of $577$): a low count reflects that day's attendance, not a
        truncated vote, and there is no quorum required for most votes.
  \item \textbf{Amendment \og{}author\fg{}.} The field used is the \emph{first
        signatory}; for a collective amendment it does not necessarily reflect the
        drafting.
  \item \textbf{Straddling legislatures.} Some files opened under the
        16\textsuperscript{th} legislature and continued under the
        17\textsuperscript{th} are attached to their legislature of origin; hence
        small discrepancies between counters depending on the perimeter chosen.
  \item \textbf{Heuristic nature of the typologies.} \og{}Mass filing\fg{},
        \og{}convergence\fg{}, \og{}amplification\fg{}, \og{}reuse\fg{},
        \og{}template answer\fg{} are labels defined by round thresholds, not
        theoretical categories.
  \item \textbf{No intent inference.} The indicators describe measurable
        regularities. They do not say \emph{why}: a mass filing may enrich the
        debate or slow it down; an attendance \emph{gap} often has prosaic causes;
        an answer with a common prefix may be substantive after that prefix.
  \item \textbf{No ideal-point estimation.} No indicator produces a left-right scale
        or a spatial map; the cohesion matrix and the ternary diagram are only
        re-projected pairwise agreements.
  \item \textbf{Dependence on the data producers.} A gap, a delay or a format change
        at the Assembly, INSEE or the Ministry of the Interior propagates directly;
        the site explicitly signals when a source is unavailable rather than
        inventing.
  \item \textbf{Snapshot.} The site updates daily; any cited figure is dated by the
        day it was computed.
\end{enumerate}

\section{Appendices}\label{sec:annexes-en}
The full tables below are also provided in CSV/Markdown in the deposit (one file per
table, plus \code{SUMMARY.md}).
\begin{itemize}\small
  \item \textbf{A1} Jaccard threshold on the full set; \textbf{A1bis} histogram of
        candidate similarities; \textbf{A2} $k$ and $\theta$ on a sample (exact
        Jaccard); \textbf{A3} LSH curve $P(s)=1-(1-s^r)^b$; \textbf{A4} exposure to
        $L_{\max}$; \textbf{A4bis} floor $s_{\min}$.
  \item \textbf{B} typology thresholds; \textbf{B2} cluster size distribution.
  \item \textbf{C0} values of \code{positionMajoritaire}; \textbf{C1} discipline
        floor.
  \item \textbf{D} cohesion matrix, \og{}shared abstention\fg{} effect.
  \item \textbf{F} \og{}coalitions by ballot\fg{} thresholds.
  \item \textbf{G} \og{}template answers\fg{} grid ($L\times m\times$ salutation
        removal).
  \item \textbf{H} \og{}strategic absenteeism\fg{} grid ($\param{T_v}\times\param{c}\times\param{C}\times\param{T_{cli}}$).
  \item \textbf{I0} census of sorts; \textbf{I1} phantom floor; \textbf{I2}
        continuous view (scatter plot).
  \item \textbf{J} ministry-filter thresholds; \textbf{K} constituency-statistics
        coverage.
\end{itemize}
, %==============================================================================
%  Parte III — Versión en español (traducción completa, refleja la Parte I).
%==============================================================================
\clearpage
\selectlanguage{spanish}
\part{Versión en español}
\setcounter{section}{0}   % cada parte lingüística se numera 1, 2, 3, … por su cuenta

\section{Introducción y alcance}

\subsection{Objeto}
\emph{577deputes.fr} es una aplicación web \emph{local-first}: descarga los
conjuntos de datos abiertos oficiales, los ingiere en una base de datos SQLite y
expone páginas de consulta e indicadores derivados para la XVII\textsuperscript{a}
legislatura de la Asamblea Nacional francesa (en curso desde julio de 2024). La
aplicación no produce ningún dato primario; todo lo que muestra es, o bien un dato
en bruto tomado tal cual de una fuente pública, o bien el resultado de un cálculo
documentado aquí.

\subsection{Fuentes}
\begin{itemize}
  \item \textbf{Asamblea Nacional} (\href{https://data.assemblee-nationale.fr}{data.assemblee-nationale.fr},
        \emph{Licence Ouverte~2.0}): actores, mandatos y órganos; expedientes
        legislativos, documentos y actos de procedimiento; enmiendas; votaciones
        públicas nominales y votos individuales; preguntas escritas, orales y al
        Gobierno; orden del día de las sesiones y actas literales.
  \item \textbf{INSEE}: población legal por circunscripción legislativa.
  \item \textbf{Ministerio del Interior} (data.gouv.fr): inscritos, votantes y
        abstención por circunscripción (elecciones legislativas de 2024).
\end{itemize}
Los conjuntos de datos de la Asamblea se actualizan a diario mediante peticiones
condicionales (ETag / Last-Modified); un fallo de descarga no altera las demás
fuentes.

\subsection{Volumen (instantánea utilizada para este documento)}
\begin{center}\small
\begin{tabular}{@{}lr@{}}
\toprule
Diputados indexados (de los cuales 577 en ejercicio) & $\approx 630$ \\
Preguntas parlamentarias & $\approx 17\,000$ \\
\quad de las cuales con respuesta publicada $>200$ caracteres & $10\,766$ \\
Enmiendas (XVII\textsuperscript{a} legislatura) & $108\,115$ \\
Votaciones públicas nominales & $6\,490$ \\
Votos individuales & $1\,043\,305$ \\
Pares (grupo, votación) con posición mayoritaria & $77\,748$ \\
\bottomrule
\end{tabular}
\end{center}

\subsection{Lo que queda fuera del alcance}
\begin{itemize}
  \item \textbf{Las votaciones a mano alzada} y \textbf{las votaciones en
        comisión}: el conjunto de datos \og{}Scrutins\fg{} solo contiene las
        votaciones públicas nominales celebradas en el pleno. Todas las cifras de
        cohesión, disciplina y absentismo se apoyan en ese perímetro
        (cf.~§\ref{sec:limites-es}).
  \item \textbf{Los expedientes abiertos bajo la XVI\textsuperscript{a}
        legislatura} y continuados bajo la XVII\textsuperscript{a}: están indexados
        pero asignados a su legislatura de origen, de ahí pequeñas discrepancias
        entre contadores.
  \item \textbf{La inferencia de intención}: ningún indicador de este documento
        pretende establecer un motivo. \og{}Presentación masiva\fg{}, \og{}enmienda
        fantasma\fg{}, \og{}absentismo estratégico\fg{} describen \emph{regularidades
        medibles}, no designios.
\end{itemize}

\section{Notación y convenciones}\label{sec:notations-es}
Escribimos $J(A,B) = |A\cap B| / |A\cup B|$ para el coeficiente de Jaccard de dos
conjuntos. Para una votación $v$, $\mathrm{a\_favor}(v)$, $\mathrm{en\_contra}(v)$,
$\mathrm{abst}(v)$ son los recuentos oficiales y
$\text{emitidos}(v)=\mathrm{a\_favor}(v)+\mathrm{en\_contra}(v)+\mathrm{abst}(v)$.
Para un grupo parlamentario $G$ y una votación $v$, $p^\star(v,G)$ es la
\emph{posición mayoritaria} del grupo tal como la publica la Asamblea (campo
\code{positionMajoritaire}), con valores en $\{\text{a favor},\text{en contra},\text{abstención}\}$.
En la instantánea utilizada, ese campo está cumplimentado para los $77\,748$ pares
$(G,v)$ y nunca toma otro valor (47,9~\% \og{}a favor\fg{}, 46,8~\% \og{}en
contra\fg{}, 5,3~\% \og{}abstención\fg{}): por lo tanto no existe, a nuestro nivel,
ningún empate exacto que arbitrar.

Cada subsección de indicador sigue el mismo esquema: \textbf{(a)~definición},
\textbf{(b)~algoritmo y parámetros}, \textbf{(c)~sensibilidad}, \textbf{(d)~literatura},
\textbf{(e)~límites}. Las cifras de sensibilidad provienen del script
\code{methodo\_sensitivity.py} ejecutado sobre la instantánea anterior; las tablas
completas están en el anexo~\ref{sec:annexes-es} y en el depósito.

\section{Los indicadores}

\subsection{Detección de enmiendas casi idénticas}\label{sec:clusters-es}

\paragraph{(a) Definición.} Se busca agrupar las enmiendas cuyo texto dispositivo
es \og{}esencialmente el mismo\fg{}, compartan o no autor. Para una enmienda $a$ se
construye $S_a$, el conjunto de los $k$-\emph{shingles} de palabras (secuencias de
$k$ palabras consecutivas) del texto normalizado, cada uno troceado a 64~bits. Dos
enmiendas $a,a'$ se consideran casi idénticas si $J(S_a,S_{a'})\ge\theta$. Un
\emph{clúster} es una componente conexa (en el sentido de la unión de esos pares)
de tamaño $\ge 2$.

\paragraph{(b) Algoritmo y parámetros.}
\begin{enumerate}[label=\arabic*.]
  \item \emph{Normalización}: supresión de etiquetas HTML; descomposición Unicode
        (NFKD) y supresión de diacríticos; paso a minúsculas; supresión de todo
        carácter fuera de \texttt{[a-z0-9 ]}; normalización de espacios; truncado a
        $L_{\max}=12\,000$ caracteres.
  \item \emph{Shingling}: $k = 5$; $S_a=\{h(w_i\cdots w_{i+k-1})\}$ con $h$ un
        hash BLAKE2b de 64~bits. Enmienda descartada si $|S_a| < s_{\min}=8$.
  \item \emph{Firma MinHash}: $K=64$ funciones hash $g_j(x)=x\oplus\sigma_j$
        ($\sigma_j$ semillas deterministas), firma $\mathrm{sig}_j(a)=\min_{x\in S_a} g_j(x)$.
  \item \emph{Bloqueo LSH}: la firma se divide en $b=16$ bandas de $r=4$
        componentes; $a$ y $a'$ son \emph{candidatos} si coinciden en al menos una
        banda. La probabilidad de que un par de similitud real $s$ sea candidato es
        $P(s)=1-(1-s^{\,r})^{b}$ (curva en S, umbral $\approx (1/b)^{1/r}\approx 0{,}55$).
  \item \emph{Verificación}: para cada par candidato, cálculo del Jaccard exacto
        sobre $S_a$, $S_{a'}$; par retenido si $J\ge\theta=0{,}80$.
  \item \emph{Agrupamiento}: union-find sobre los pares retenidos; clústeres de
        tamaño $\ge 2$.
\end{enumerate}
En la instantánea: $62\,476$ enmiendas pasan el umbral $s_{\min}$; se obtienen
$\approx 9\,460$ clústeres que agrupan $28\,124$ enmiendas (la base de producción
registra $28\,055$ — la diferencia de $0{,}25~\%$ proviene de una salvaguarda de
tamaño de \emph{bucket} introducida en el script de verificación, y constituye
nuestro control de reproducibilidad).

\paragraph{(c) Sensibilidad.}
\emph{Umbral $\theta$} (conjunto completo, LSH de producción):
\begin{center}\small
\begin{tabular}{rrrr}
\toprule
$\theta$ & pares retenidos & clústeres & enmiendas agrupadas (\%) \\
\midrule
0,70 & 180\,407 & 9\,638 & 32\,953 \;(52,8~\%) \\
0,75 & 147\,438 & 9\,334 & 30\,794 \;(49,3~\%) \\
0,78 & \phantom{0}88\,007 & 9\,649 & 29\,159 \;(46,7~\%) \\
\textbf{0,80} & \phantom{0}\textbf{60\,381} & \textbf{9\,461} & \textbf{28\,124 \;(45,0~\%)} \\
0,82 & \phantom{0}54\,380 & 9\,253 & 26\,834 \;(43,0~\%) \\
0,85 & \phantom{0}43\,701 & 9\,023 & 25\,212 \;(40,4~\%) \\
0,90 & \phantom{0}34\,480 & 8\,440 & 22\,782 \;(36,5~\%) \\
0,95 & \phantom{0}30\,679 & 7\,840 & 20\,934 \;(33,5~\%) \\
\bottomrule
\end{tabular}
\end{center}
La sensibilidad es \emph{moderada}: $\pm0{,}05$ alrededor de $0{,}80$ desplaza el
número de enmiendas agrupadas en $\approx\pm10$ a $\pm15~\%$. Pero la distribución
de las similitudes de los pares candidatos (anexo~\ref{sec:annexes-es}) muestra que
$0{,}80$ cae en una \emph{zona de transición}: hay una gran masa de pares en
$[0{,}75;0{,}85)$ (enmiendas que comparten mucho texto estándar sin ser copias),
mientras que las copias casi literales forman un pico distinto por encima de
$0{,}95$. Un umbral más estricto ($0{,}90$) solo retendría las copias casi
verbatim; uno más laxo ($0{,}70$--$0{,}75$) absorbería también las reescrituras
pesadas. La elección de $0{,}80$ es un compromiso asumido, no un punto
\og{}natural\fg{} de la nube. Para $\theta<0{,}80$, la cobertura de los pares
candidatos por el LSH 16$\times$4 disminuye (véase abajo): los recuentos mostrados
ahí son cotas inferiores.

\emph{Tamaño de los shingles $k$} (muestra de $\approx 3\,500$ enmiendas, Jaccard
exacto sobre todos los pares, $\theta=0{,}80$): $k=3 \Rightarrow 12{,}4~\%$
agrupadas; $k=5 \Rightarrow 8{,}3~\%$; $k=7 \Rightarrow 6{,}5~\%$. Shingles más
cortos son más permisivos; $k=5$ es el ajuste mediano.

\emph{Configuración LSH} ($b\times r$): $P(s)=1-(1-s^r)^b$. Para el ajuste de
producción 16$\times$4: $P(0{,}80)=0{,}9998$, $P(0{,}70)=0{,}988$, $P(0{,}50)=0{,}64$.
El LSH solo afecta a la \emph{cobertura} de los pares candidatos; la
\emph{precisión} final está garantizada por la verificación Jaccard exacta que
sigue. En el umbral de producción, la cobertura es por tanto casi perfecta.

\emph{Truncado $L_{\max}$ y umbral mínimo $s_{\min}$}: solo $192$ enmiendas
($0{,}18~\%$) superan los $12\,000$ caracteres tras la normalización; elevar
$s_{\min}$ de $8$ a $50$ excluiría el $38~\%$ de las enmiendas actualmente retenidas
(el valor $8$ es deliberadamente permisivo, para conservar los textos dispositivos
cortos).

\paragraph{(d) Literatura.} La representación por \emph{shingles} y la estimación de
la similitud de Jaccard mediante min-wise hashing se deben a
Broder~\cite{broder1997,broder2000}; el bloqueo por bandas (LSH) sigue a
Gionis-Indyk-Motwani~\cite{gionis1999} y la exposición de
Leskovec-Rajaraman-Ullman~\cite{lru2014}. La aplicación a la detección de
reutilización de texto legislativo está explorada notablemente por el
\emph{Legislative Influence Detector}~\cite{burgess2016} y los trabajos sobre la
difusión de \og{}model bills\fg{} (Hertel-Fernandez~\cite{hf2014},
Jansa et al.~\cite{jansa2019}, Linder et al.~\cite{linder2020}). Nuestro aporte no
es metodológico: es la \emph{aplicación} de piezas conocidas al flujo de enmiendas
de la Asamblea, con una tipología descriptiva.

\paragraph{(e) Límites.} (i)~La normalización \emph{trunca} más allá de $12\,000$
caracteres: dos enmiendas muy largas solo se comparan en su prefijo. (ii)~El
shingling de palabras ignora el orden más allá de la ventana $k$: una reordenación
de párrafos puede hacer caer el Jaccard. (iii)~El umbral $0{,}80$ está en la zona
ambigua de la nube: la frontera exacta de los clústeres es convencional. (iv)~Un
\og{}clúster\fg{} no dice nada de una coordinación intencional: un mismo modelo
puede circular vía una federación de asociaciones, un sindicato, un sitio
militante, sin concertación entre diputados.

\subsection{Tipología de los clústeres}\label{sec:typo-es}

\paragraph{(a) Definición.} Para un clúster $C$, sean $n(C)=|C|$ su tamaño y $g(C)$
el número de grupos parlamentarios distintos entre sus enmiendas. Se clasifica:
\[
\tau(C)=\begin{cases}
\text{presentación masiva} & \text{si } n(C)\ge \param{T_{masa}}=15,\\[2pt]
\text{convergencia intergrupos} & \text{si no, si } g(C)\ge 2,\\[2pt]
\text{amplificación intragrupo} & \text{si no, si } g(C)\le 1 \text{ y } n(C)\ge \param{T_{ampl}}=3,\\[2pt]
\text{reutilización simple} & \text{en otro caso.}
\end{cases}
\]
El tamaño siempre prevalece; no se infiere ninguna intención (la etiqueta interna
\og{}obstrucción\fg{} se neutralizó a \og{}presentación masiva\fg{}).

\paragraph{(b) Algoritmo y parámetros.} Cálculo directo sobre la tabla de clústeres
($\param{T_{masa}}=15$, $\param{T_{ampl}}=3$). En la instantánea (clústeres de
$\theta=0{,}80$): $91$ clústeres \og{}presentación masiva\fg{} ($2\,578$ enmiendas),
$4\,314$ \og{}convergencia\fg{} ($13\,834$), $773$ \og{}amplificación\fg{}
($3\,146$), $4\,283$ \og{}reutilización\fg{} ($8\,566$).

\paragraph{(c) Sensibilidad.}
\begin{center}\small
\begin{tabular}{rr rr rr rr}
\toprule
$\param{T_{masa}}$ & $\param{T_{ampl}}$ & \multicolumn{2}{c}{pres. masiva} & \multicolumn{2}{c}{convergencia} & \multicolumn{2}{c}{amplificación} \\
 & & cl. & enm. & cl. & enm. & cl. & enm. \\
\midrule
10 & 3 & 203 & 3\,875 & 4\,226 & 12\,812 & 749 & 2\,871 \\
\textbf{15} & \textbf{3} & \textbf{91} & \textbf{2\,578} & \textbf{4\,314} & \textbf{13\,834} & \textbf{773} & \textbf{3\,146} \\
20 & 3 & 43 & 1\,776 & 4\,352 & 14\,472 & 783 & 3\,310 \\
30 & 3 & 19 & 1\,218 & 4\,370 & 14\,880 & 789 & 3\,460 \\
50 & 3 & \phantom{0}7 & \phantom{0,}784 & 4\,377 & 15\,121 & 794 & 3\,653 \\
\midrule
15 & 2 & 91 & 2\,578 & 4\,314 & 13\,834 & 5\,056 & 11\,712 \\
15 & 5 & 91 & 2\,578 & 4\,314 & 13\,834 & 160 & 1\,138 \\
\bottomrule
\end{tabular}
\end{center}
El \emph{número} de clústeres \og{}presentación masiva\fg{} es sensible a
$\param{T_{masa}}$ ($91$ a $7$ entre los umbrales $15$ y $50$) pero el
\emph{volumen} de enmiendas afectado lo es menos; el umbral $\param{T_{ampl}}$
sobre todo hace deslizar clústeres entre \og{}amplificación\fg{} y
\og{}reutilización\fg{}. La distribución por tamaño está en el
anexo~\ref{sec:annexes-es}.

\paragraph{(d) Literatura.} La tipología es ad~hoc y descriptiva; se inscribe en el
campo del estudio de la \emph{obstrucción parlamentaria} y de la presentación
masiva de enmiendas como táctica de debate (cf.~Koger~\cite{koger2010} para el caso
estadounidense; el propio Reglamento de la Asamblea Nacional regula la
\og{}discusión común\fg{} de enmiendas idénticas).

\paragraph{(e) Límites.} Cuatro casillas, umbrales redondos: un clúster de $14$
enmiendas idénticas y uno de $16$ caen a uno y otro lado de una frontera sin
fundamento teórico. La casilla \og{}convergencia\fg{} no distingue un consenso
técnico de una alianza puntual.

\subsection{Disciplina de voto (cohesión individual)}\label{sec:discipline-es}

\paragraph{(a) Definición.} Sea $G_d$ el grupo del diputado $d$. Sea $V_d$ el
conjunto de las votaciones donde $d$ ha emitido un voto $p_d(v)\in\{\text{a favor},\text{en contra},\text{abstención}\}$
\emph{y} donde $p^\star(v,G_d)$ está definida. Se pone
\[
\mathrm{emitidos}(d)=|V_d|,\quad
\mathrm{alineados}(d)=\big|\{v\in V_d : p_d(v)=p^\star(v,G_d)\}\big|,\quad
\mathrm{disciplina}(d)=\frac{\mathrm{alineados}(d)}{\mathrm{emitidos}(d)}.
\]
Se muestra en las clasificaciones solo si $\mathrm{emitidos}(d)\ge 50$ (o $\ge 100$
para la lista de \og{}disidentes\fg{}, ordenada por disciplina creciente); los
diputados no inscritos (NI) quedan excluidos.

\paragraph{(b) Algoritmo y parámetros.} Todo se apoya en el campo oficial
$p^\star$: \emph{no} calculamos nosotros la posición mayoritaria de un grupo. En la
instantánea, $\mathrm{disciplina}$ tiene una media de $\approx 91{,}9~\%$, una
mediana de $\approx 92{,}5~\%$, un decil inferior de $\approx 85~\%$.

\paragraph{(c) Sensibilidad.} Efecto del umbral mínimo de votos emitidos:
\begin{center}\small
\begin{tabular}{rrrrrr}
\toprule
umbral & elegibles & media & mediana & $p_{10}$ & $p_{90}$ \\
\midrule
0   & 639 & 91,92 & 92,56 & 84,92 & 98,30 \\
20  & 628 & 91,87 & 92,50 & 84,93 & 98,25 \\
\textbf{50}  & \textbf{623} & \textbf{91,85} & \textbf{92,50} & \textbf{84,99} & \textbf{98,21} \\
100 & 611 & 91,83 & 92,45 & 85,02 & 98,18 \\
200 & 600 & 91,87 & 92,46 & 85,06 & 98,18 \\
\bottomrule
\end{tabular}
\end{center}
El indicador es \emph{muy robusto} al umbral: la media y los cuantiles solo varían
unas centésimas de punto entre los umbrales $0$ y $200$. La elección $\ge 50$ no es,
pues, un parámetro sensible — solo sirve para evitar mostrar un \og{}100~\%\fg{}
basado en tres votos.

\paragraph{(d) El \og{}sesgo circular\fg{}.} La posición $p^\star(v,G_d)$ es la
pluralidad de los votos de los miembros de $G_d$, y $d$ forma parte de él: $d$ es,
pues, comparado, en parte, con una posición que él ha contribuido a definir. Es un
sesgo \emph{real} pero (i)~no es obra nuestra — es la convención oficial de la
Asamblea; (ii)~para un grupo de $50$ a $120$ miembros, el peso de un voto individual
en la pluralidad está acotado por $\le 2~\%$; (iii)~el caso de un empate exacto, que
haría ambigua la pluralidad, lo arbitra aguas arriba la Asamblea y nunca aparece en
los datos (el campo $p^\star$ está siempre cumplimentado y siempre en
$\{\text{a favor},\text{en contra},\text{abstención}\}$, cf.~§\ref{sec:notations-es}).

\paragraph{(e) Literatura.} La disciplina así definida es la versión
\emph{individual} (por diputado) de los índices de cohesión partidaria: índice de
Rice~\cite{rice1925}, \emph{Agreement Index} de Hix-Noury-Roland~\cite{hnr2005,hnr2007},
discusión de Carey~\cite{carey2007} sobre los \og{}principales en competencia\fg{}.
\emph{No} estima puntos ideales: para ello harían falta métodos tipo NOMINATE o
\emph{Optimal Classification} (Poole-Rosenthal~\cite{pr1997}, Poole~\cite{poole2005}),
fuera de alcance.

\paragraph{(e') Límites.} (i)~Perímetro restringido a las votaciones públicas
nominales: la \og{}disciplina\fg{} medida es la de los votos \emph{emitidos y
registrados}, no la unidad de grupo en general. (ii)~Una abstención conforme con la
posición del grupo cuenta como \og{}alineada\fg{}. (iii)~Un diputado a menudo
ausente tiene un $\mathrm{emitidos}$ bajo: su disciplina se mide sobre pocos casos
(de ahí el umbral).

\subsection{Matriz de cohesión entre grupos}\label{sec:matrice-es}

\paragraph{(a) Definición.} Para dos grupos $G_1,G_2$, sea $V_{12}$ el conjunto de
las votaciones donde ambos tienen una posición mayoritaria definida. La cohesión es
\[
\mathrm{coh}(G_1,G_2)=\frac{\big|\{v\in V_{12}:p^\star(v,G_1)=p^\star(v,G_2)\}\big|}{|V_{12}|},
\qquad \mathrm{coh}(G,G)=1.
\]

\paragraph{(b) Algoritmo y parámetros.} Ningún parámetro libre; una auto-unión
sobre la tabla de posiciones mayoritarias.

\paragraph{(c) Sensibilidad.} Una elección discutible: una \og{}abstención $=$
abstención\fg{} cuenta como acuerdo. Recalculando la matriz sin las votaciones donde
ambos grupos se abstienen, la diferencia absoluta media es de $0{,}21$~puntos y el
máximo de $0{,}83$~puntos. El efecto es, pues, \emph{despreciable}.

\paragraph{(d) Literatura.} Es una matriz de \emph{acuerdo} entre bloques en el
sentido de Hix-Noury-Roland~\cite{hnr2007}; no proporciona ningún orden
unidimensional (para ello, véanse los métodos espaciales citados en
§\ref{sec:discipline-es}).

\paragraph{(e) Límites.} Cohesión calculada sobre \emph{posiciones de grupo}, no
sobre votos individuales: un grupo dividido $51/49$ y un grupo unánime cuentan
igual. Dos \og{}en contra\fg{} por motivos opuestos son indistinguibles de un
acuerdo de fondo.

\subsection{Diagrama ternario de los bloques}\label{sec:ternaire-es}

\paragraph{(a) Definición.} Dados tres grupos \emph{ancla} $A_1,A_2,A_3$ situados en
los vértices de un triángulo equilátero en posiciones $\mathbf{x}_{A_1},\mathbf{x}_{A_2},\mathbf{x}_{A_3}$,
un grupo $G$ de cohesiones $c_i=\mathrm{coh}(G,A_i)$ se sitúa en
\[
(\alpha,\beta,\gamma)=\frac{(c_1,c_2,c_3)}{c_1+c_2+c_3},\qquad
\mathbf{x}_G=\alpha\,\mathbf{x}_{A_1}+\beta\,\mathbf{x}_{A_2}+\gamma\,\mathbf{x}_{A_3}.
\]
Las anclas por defecto son LFI, EPR, RN; el usuario puede cambiarlas.

\paragraph{(b)--(c) Algoritmo y sensibilidad.} El único \og{}parámetro\fg{} es la
elección de las tres anclas: el diagrama no es más que una \emph{re-proyección} de
la matriz de cohesión §\ref{sec:matrice-es} sobre tres ejes elegidos, y cambia con
ellos.

\paragraph{(d)--(e) Literatura y límites.} Herramienta de visualización, sin
pretensión de medida: la posición de un grupo \emph{depende del marco de
referencia}. No leer las distancias como distancias ideológicas.

\subsection{Coaliciones por votación (por tema)}\label{sec:coalitions-es}

\paragraph{(a) Definición.} Se seleccionan las votaciones \og{}muy concurridas y
divisivas\fg{}: (a)~$\text{emitidos}(v)\ge \param{T_v}=300$ ($\approx 52~\%$ de los
577 escaños); (b)~$|\mathrm{a\_favor}(v)-\mathrm{en\_contra}(v)|\le \param{T_e}=50$.
Se ordena por (b) creciente y luego fecha decreciente y se conservan las $N=10$
primeras; si pasan menos de $N/2$ votaciones, se relaja el criterio (b). Para cada
votación retenida, se muestra la posición mayoritaria de cada grupo.

\paragraph{(b)--(c) Parámetros y sensibilidad.} Número de votaciones elegibles según
$(\param{T_v},\param{T_e})$:
\begin{center}\small
\begin{tabular}{r|rrrrr}
\toprule
$\param{T_v}\backslash\param{T_e}$ & 10 & 25 & 50 & 75 & 100 \\
\midrule
200 & \phantom{0}95 & 252 & 565 & 913 & 1\,200 \\
250 & \phantom{0}25 & \phantom{0}81 & 193 & 328 & \phantom{0,}440 \\
\textbf{300} & \phantom{0}\textbf{11} & \phantom{0}\textbf{36} & \phantom{00}\textbf{79} & 118 & \phantom{0,}155 \\
350 & \phantom{00}3 & \phantom{0}14 & \phantom{00}24 & \phantom{0}47 & \phantom{0,0}55 \\
400 & \phantom{00}2 & \phantom{00}7 & \phantom{00}11 & \phantom{0}19 & \phantom{0,0}20 \\
\bottomrule
\end{tabular}
\end{center}
En el ajuste de producción $(300,50)$, $79$ votaciones pasan: el repliegue (b) no se
activa. La selección es sensible a ambos umbrales, como cabía esperar (es un filtro,
no un estimador).

\paragraph{(d)--(e) Literatura y límites.} Enfoque puramente descriptivo: se muestra
\emph{quién vota como quién} en las votaciones ajustadas, sin modelo. Límite
principal: una votación \og{}ajustada en votos\fg{} puede serlo por baja
concurrencia asimétrica tanto como por una verdadera fractura.

\subsection{Respuestas ministeriales \og{}tipo\fg{}}\label{sec:templates-es}

\paragraph{(a) Definición.} Para cada respuesta escrita $q$ cuyo texto en bruto
supera los $200$ caracteres, se normaliza (supresión HTML, espacios), se retira
eventualmente la fórmula de cortesía inicial, se pasa a minúsculas: $t_q$. La
respuesta se descarta si $|t_q|<100$. La \emph{clave de prefijo} es
$\pi_L(q)=t_q[1:L]$ con $L=250$. Una \og{}respuesta tipo\fg{} es un prefijo
compartido por al menos $m=3$ respuestas distintas. La \emph{tasa de respuestas
tipo} es
\[
\rho \;=\; \frac{1}{|Q_{\text{eleg}}|}\sum_{\pi \,:\, |\{q:\pi_L(q)=\pi\}|\ge m} \big|\{q:\pi_L(q)=\pi\}\big|.
\]
En la instantánea, $|Q_{\text{eleg}}|=10\,766$ y $\rho=27{,}9~\%$.

\paragraph{(b)--(c) Parámetros y sensibilidad.} Rejilla de la tasa $\rho$ (en \%)
según la longitud de prefijo $L$ y el umbral de ocurrencias $m$:
\begin{center}\small
\begin{tabular}{r|rrrr}
\toprule
$L\backslash m$ & 2 & \textbf{3} & 5 & 10 \\
\midrule
100  & 42,1 & 31,6 & 23,1 & 13,4 \\
\textbf{250}  & 37,0 & \textbf{27,9} & 20,3 & 11,6 \\
500  & 33,7 & 24,6 & 17,2 & \phantom{0}9,9 \\
1000 & 30,2 & 22,1 & 15,3 & \phantom{0}9,0 \\
\bottomrule
\end{tabular}
\end{center}
La tasa es \emph{fuertemente sensible al umbral $m$} (cociente $\approx 3$--$4$
entre $m=2$ y $m=10$) y \emph{moderadamente sensible a $L$} (de $\approx 32~\%$ a
$\approx 22~\%$ con $m$ fijo). En este corpus, retirar la fórmula de cortesía no
tiene \emph{ningún} efecto (las variantes \og{}con\fg{} y \og{}sin\fg{} retirada
coinciden exactamente, el texto almacenado no empieza por esas fórmulas). El valor
mostrado en el sitio debe, pues, leerse como \emph{un} punto de una familia de
medidas, no como una constante: la formulación honesta es \og{}$\approx 28~\%$ de
las respuestas comparten un prefijo de $250$~caracteres con al menos otras dos\fg{}.

\paragraph{(d) Literatura.} Las preguntas parlamentarias como instrumento de
control, y sus respuestas como objeto de estudio, son tratadas por
Wiberg~\cite{wiberg1995}, Russo-Wiberg~\cite{rw2010}, Martin~\cite{martin2011}. La
detección de respuestas estandarizadas por agrupamiento de prefijos es una
heurística \emph{casera}, próxima a la detección de \og{}boilerplate\fg{}; no es una
medida validada de \og{}lenguaje de madera\fg{}.

\paragraph{(e) Límites.} (i)~Un prefijo común es una condición \emph{necesaria} pero
no \emph{suficiente} de una respuesta \og{}vacía\fg{}: dos respuestas pueden
compartir un encabezado de encuadre y luego divergir por completo. (ii)~El umbral
$m=3$ y la longitud $L=250$ son arbitrarios (cf.~rejilla). (iii)~El corpus se limita
a las respuestas publicadas de más de $200$ caracteres; las preguntas sin respuesta
no entran en el denominador de $\rho$.

\subsection{Absentismo estratégico}\label{sec:absence-es}

\paragraph{(a) Definición.} Entre las votaciones con $\text{emitidos}(v)\ge \param{T_v}=200$,
se dice que $v$ es \emph{divisiva} si
$|\mathrm{a\_favor}(v)-\mathrm{en\_contra}(v)|/\text{emitidos}(v)<\param{c}=0{,}10$,
\emph{consensual} si $>\param{C}=0{,}50$, en otro caso sin clasificar. Para un
diputado $d$, sea $\mathrm{pres}_{div}(d)$ (resp.\ $\mathrm{pres}_{con}(d)$) la
\emph{proporción} de votaciones divisivas (resp.\ consensuales) en las que $d$ ha
emitido un voto. Se pone $\mathrm{gap}(d)=\mathrm{pres}_{con}(d)-\mathrm{pres}_{div}(d)$
(en puntos) y $\mathrm{ratio}(d)=\mathrm{pres}_{con}(d)/\mathrm{pres}_{div}(d)$. Se
muestra si $d$ está en ejercicio y $\mathrm{tot}_{div}(d)\ge \param{T_{div}}=30$.

\paragraph{(b)--(c) Parámetros y sensibilidad.} Cuatro umbrales
($\param{T_v},\param{c},\param{C},\param{T_{div}}$). Los valores exactos de las
$4^4$ combinaciones exploradas figuran en el anexo~\ref{sec:annexes-es}; se observa
que la mediana del \emph{gap} está cerca de $0$ — la mayoría de los diputados tienen
una presencia similar en votaciones divisivas y consensuales — y que solo los
\emph{outliers} tienen un \emph{gap} marcado, cuya amplitud varía de $\approx 13$ a
$\approx 31$ puntos según los umbrales.

\paragraph{(d)--(e) Literatura y límites.} La idea se relaciona con la literatura
sobre el \emph{shirking} legislativo y la participación selectiva en los votos.
Límites fuertes: (i)~no votar $\ne$ estar ausente — un diputado presente puede
abstenerse de participar; (ii)~\og{}divisiva\fg{} se define por el \emph{resultado}
del voto, no por su relevancia política \emph{a priori}; (iii)~un \emph{gap}
positivo tiene muchas causas no estratégicas (agenda, enfermedad, circunscripción
lejana). Una señal a explorar, no un veredicto.

\subsection{Enmiendas \og{}fantasma\fg{}}\label{sec:fantomes-es}

\paragraph{(a) Definición.} Para el primer firmante $d$, sobre el conjunto de sus
enmiendas:
\[
\begin{aligned}
\mathrm{defendidas}(d) &= \big|\{a : \mathrm{sort}(a)\in\{\text{Aprobada, Rechazada, Retirada, Discutida}\}\}\big|,\\[2pt]
\mathrm{fantasma}(d) &= \big|\{a : \mathrm{sort}(a)=\text{No defendida, o vacío}\}\big|,\\[2pt]
\mathrm{pct\_fantasma}(d) &= \mathrm{fantasma}(d)\,/\,\mathrm{total}(d).
\end{aligned}
\]
Se muestra si $\mathrm{total}(d)\ge \param{T_a}=50$; ordenado por
$\mathrm{pct\_fantasma}$ decreciente.

\paragraph{(b)--(c) Parámetros y sensibilidad.} Censo de los códigos \og{}sort\fg{}
(porcentaje del total de enmiendas): Rechazada 22,9~\%; Inadmisible 16,3~\%; En
tramitación 14,1~\%; Aprobada 13,0~\%; Retirada 10,2~\%; Discutida 9,0~\%; Decaída
8,7~\%; No defendida 5,9~\%; Borrada $\approx$0~\%. (Ninguna enmienda tiene un sort
vacío o nulo en la instantánea.) Por tanto $\mathrm{fantasma}$ se reduce en la
práctica a \og{}No defendida\fg{} ($5{,}9~\%$), $\mathrm{defendidas}$ cubre
$\approx 55~\%$, y el resto (Inadmisible, En tramitación, Decaída) no se cuenta
\emph{ni} como defendida \emph{ni} como fantasma. Efecto del umbral $\param{T_a}$:
\begin{center}\small
\begin{tabular}{rrrrr}
\toprule
$\param{T_a}$ & autores elegibles & de los cuales en ejercicio & $\overline{\mathrm{pct\_fant.}}$ & $\mathrm{pct\_fant.}^{\max}$ \\
\midrule
10  & 558 & 513 & 7,4 & 61,5 \\
20  & 530 & 491 & 7,2 & 52,4 \\
\textbf{50}  & \textbf{439} & \textbf{412} & \textbf{6,5} & \textbf{50,0} \\
100 & 313 & 300 & 6,6 & 50,0 \\
200 & 167 & 163 & 5,8 & 34,7 \\
\bottomrule
\end{tabular}
\end{center}
La media de la tasa es estable ($\approx 6$--$7{,}5~\%$); el \emph{máximo} decrece
con el umbral (los valores extremos los llevan autores de bajo volumen). Una
\emph{vista continua} (diagrama de dispersión $\mathrm{total}\times\mathrm{pct\_fantasma}$,
sin umbral) se proporciona en el depósito para no ocultar esta dependencia.

\paragraph{(d)--(e) Literatura y límites.} Sin literatura directa; cuantificamos una
asimetría entre \emph{presentación} y \emph{defensa} de enmiendas. Límites: (i)~la
categoría \og{}No defendida\fg{} también cubre restricciones de tiempo en el pleno,
independientes de la voluntad del autor; (ii)~la elección de incluir
\og{}Discutida\fg{} en \og{}defendida\fg{} y de excluir \og{}Decaída\fg{} es
discutible (la rejilla de sensibilidad sobre la definición de \og{}defendida\fg{}
está en el depósito); (iii)~\og{}primer firmante\fg{} $\ne$ \og{}autor real\fg{} en
el caso de enmiendas colectivas.

\subsection{Plazos de respuesta ministeriales y clasificaciones}\label{sec:delais-es}

\paragraph{(a) Definición.} Para una pregunta $q$, $\mathrm{plazo}(q)$ es el número
de días entre la presentación y la respuesta (indefinido si no se responde). Para un
ministerio (en el sentido del \og{}ministerio interpelado\fg{} corto) $\mu$:
\[
\overline{\mathrm{plazo}}(\mu)=\operatorname*{media}_{q:\,\mathrm{plazo}(q)\ \text{definido}}\mathrm{plazo}(q),\qquad
\mathrm{tasa\_resp}(\mu)=\frac{|\{q:\text{estado}=\text{con respuesta}\}|}{|\{q\ \text{dirigidas a }\mu\}|}.
\]
Se muestra si $|\{q\}|\ge \param{T_q}=30$ \emph{y} $|\{q:\text{estado}=\text{con respuesta}\}|\ge \param{T_r}=5$.
La tasa de respuesta se muestra junto al plazo: un plazo corto no significa lo mismo
si solo $1$ pregunta de cada $10$ ha obtenido respuesta.

\paragraph{(b)--(c) Parámetros y sensibilidad.} Número de ministerios retenidos
según $(\param{T_q},\param{T_r})$: $(10,1)\to 73$; $(20,*)\to 62$; $(30,*)\to 57$;
$(50,*)\to 52$; $(100,*)\to 40$. El umbral $\param{T_r}$ es casi inoperante (un
ministerio que recibe $\ge 30$ preguntas casi siempre ha respondido $\ge 5$).

\paragraph{(d)--(e) Literatura y límites.} Cf.~§\ref{sec:templates-es} para las
preguntas como objeto de estudio. Límites: la media es sensible a los valores
extremos (una respuesta a $700$~días arrastra la media); el \og{}ministerio
interpelado\fg{} puede diferir del \og{}ministerio asignado\fg{} que efectivamente
responde.

\subsection{Indicadores derivados simples}\label{sec:simples-es}

\paragraph{(a) Actividad acumulada.} $\mathrm{act}(d)=n_Q(d)+n_A(d)+n_V(d)$
(preguntas formuladas $+$ enmiendas presentadas $+$ votos emitidos).
\emph{Advertencia estructural}: $n_V$ es de un orden de magnitud superior a $n_Q$ y
$n_A$ (un diputado asiduo emite miles de votos pero presenta decenas de enmiendas y
formula decenas de preguntas); el índice acumulado está, pues, \emph{mecánicamente
dominado por la presencia en las votaciones}. El sitio lo acompaña sistemáticamente
de las subclasificaciones aisladas (solo preguntas, solo enmiendas, solo presencia).

\paragraph{(b) Tasa de aprobación de las enmiendas.} $\mathrm{aprobadas}(d)/\mathrm{total}(d)$
(por diputado o por grupo).

\paragraph{(c) Tasa de respuesta a las preguntas.} cf.~§\ref{sec:delais-es}.

\paragraph{(d) Edad, antigüedad.} Derivadas de la fecha de nacimiento y de la fecha
de inicio del mandato. Triviales.

\paragraph{(e) Constructor de rankings parametrizables.} El sitio ofrece una
treintena de \og{}métricas\fg{} combinables con filtros opcionales (por grupo,
período, tipo de texto, etc.). Todas son \emph{re-exposiciones} de las primitivas
anteriores (recuentos, ratios, medias); no hay ningún método adicional.

\paragraph{(f) Estadísticas por circunscripción.} $\mathrm{participación}=\mathrm{votantes}/\mathrm{inscritos}$,
población (INSEE), inscritos (Ministerio del Interior), unidos por
$(\text{código de departamento},\ \text{número de circunscripción})$. Cobertura en
la instantánea: $566$ circunscripciones con población INSEE (las $11$
circunscripciones de los franceses en el extranjero no la tienen), $577$ con
inscritos.

\section{Reproducibilidad}\label{sec:repro-es}
\begin{itemize}
  \item \textbf{Código de la aplicación}: público, con licencia MIT
        (\href{https://github.com/RogericBot/577deputes}{github.com/RogericBot/577deputes}).
        La base de datos y las fotografías no se redistribuyen; se reconstruyen con
        el comando \code{anqp bootstrap} (descarga e ingestión de las fuentes).
  \item \textbf{Cachés derivados}: las tablas \code{deputy\_discipline\_cache},
        \code{groupe\_discipline\_cache}, \code{deputy\_amd\_cache},
        \code{groupe\_amd\_cache}, \code{scrutin\_groupes}, \code{amendement\_clusters},
        \code{circo\_stats} se recalculan en la ingestión; no contienen nada que no
        derive de las fuentes.
  \item \textbf{Análisis de sensibilidad}: el script \code{scripts/methodo\_sensitivity.py}
        (solo biblioteca estándar, solo lectura) reproduce todas las tablas de este
        documento; está adjunto al depósito Zenodo con sus salidas (CSV y Markdown).
        Todos los sorteos aleatorios usan una semilla fija.
  \item \textbf{Instantánea}: las cifras de este documento corresponden a una copia
        de la base tomada para su redacción (volumen en §1.3); pueden diferir
        ligeramente del estado actual del sitio, que se actualiza a diario.
\end{itemize}

\section{Límites generales y sesgos conocidos}\label{sec:limites-es}
\begin{enumerate}[label=L\arabic*.]
  \item \textbf{Perímetro de las votaciones.} Solo existen en los datos las
        votaciones públicas nominales celebradas en el pleno. Las votaciones a mano
        alzada y las votaciones en comisión no figuran. Todo indicador de voto
        (disciplina, cohesión, absentismo, coaliciones) hereda este límite. Además,
        la Asamblea vota muy a menudo con solo una fracción del hemiciclo presente
        (mediana $\approx 150$ votantes de $577$): un recuento bajo refleja la
        concurrencia de ese día, no un voto truncado, y no se exige quórum para la
        mayoría de los votos.
  \item \textbf{\og{}Autor\fg{} de una enmienda.} El campo usado es el \emph{primer
        firmante}; para una enmienda colectiva no refleja necesariamente la
        redacción.
  \item \textbf{Legislaturas a caballo.} Algunos expedientes abiertos bajo la
        XVI\textsuperscript{a} legislatura y continuados bajo la XVII\textsuperscript{a}
        están asignados a su legislatura de origen; de ahí pequeñas discrepancias
        entre contadores según el perímetro elegido.
  \item \textbf{Carácter heurístico de las tipologías.} \og{}Presentación masiva\fg{},
        \og{}convergencia\fg{}, \og{}amplificación\fg{}, \og{}reutilización\fg{},
        \og{}respuesta tipo\fg{} son etiquetas definidas por umbrales redondos, no
        categorías teóricas.
  \item \textbf{Ninguna inferencia de intención.} Los indicadores describen
        regularidades medibles. No dicen \emph{por qué}: una presentación masiva
        puede enriquecer el debate o ralentizarlo; un \emph{gap} de presencia tiene
        a menudo causas prosaicas; una respuesta con prefijo común puede ser
        sustancial tras ese prefijo.
  \item \textbf{Sin estimación de puntos ideales.} Ningún indicador produce una
        escala izquierda-derecha ni un mapa espacial; la matriz de cohesión y el
        diagrama ternario son solo acuerdos por pares re-proyectados.
  \item \textbf{Dependencia de los productores de datos.} Una laguna, un retraso o
        un cambio de formato en la Asamblea, el INSEE o el Ministerio del Interior se
        repercute directamente; el sitio señala explícitamente cuándo una fuente no
        está disponible en lugar de inventar.
  \item \textbf{Instantánea.} El sitio se actualiza a diario; toda cifra citada está
        fechada en el día de su cálculo.
\end{enumerate}

\section{Anexos}\label{sec:annexes-es}
Las tablas completas siguientes se proporcionan también en CSV/Markdown en el
depósito (un archivo por tabla, más \code{SUMMARY.md}).
\begin{itemize}\small
  \item \textbf{A1} umbral Jaccard sobre el conjunto completo; \textbf{A1bis}
        histograma de las similitudes candidatas; \textbf{A2} $k$ y $\theta$ sobre
        muestra (Jaccard exacto); \textbf{A3} curva LSH $P(s)=1-(1-s^r)^b$;
        \textbf{A4} exposición a $L_{\max}$; \textbf{A4bis} umbral $s_{\min}$.
  \item \textbf{B} umbrales de tipología; \textbf{B2} distribución de los tamaños de
        clúster.
  \item \textbf{C0} valores de \code{positionMajoritaire}; \textbf{C1} umbral de
        disciplina.
  \item \textbf{D} matriz de cohesión, efecto \og{}abstención compartida\fg{}.
  \item \textbf{F} umbrales \og{}coaliciones por votación\fg{}.
  \item \textbf{G} rejilla \og{}respuestas tipo\fg{} ($L\times m\times$ retirada de
        las fórmulas).
  \item \textbf{H} rejilla \og{}absentismo estratégico\fg{} ($\param{T_v}\times\param{c}\times\param{C}\times\param{T_{div}}$).
  \item \textbf{I0} censo de los sorts; \textbf{I1} umbral fantasma; \textbf{I2}
        vista continua (diagrama de dispersión).
  \item \textbf{J} umbrales del filtro de ministerios; \textbf{K} cobertura de las
        estadísticas de circunscripción.
\end{itemize}
).
%==============================================================================
\documentclass[11pt,a4paper]{article}

\usepackage[utf8]{inputenc}
\usepackage[T1]{fontenc}
\usepackage{lmodern}
% Trilingue : français (principal), anglais, espagnol. On peut basculer avec
% \selectlanguage{...} dans en.tex / es.tex.
\usepackage[english,spanish,french]{babel}
\usepackage[a4paper,margin=2.4cm]{geometry}
\usepackage{amsmath,amssymb}
\usepackage{booktabs}
\usepackage{array}
\usepackage{longtable}
\usepackage{enumitem}
\usepackage{xcolor}
\usepackage{microtype}
\usepackage{csquotes}
\usepackage[hidelinks]{hyperref}
\usepackage{titlesec}

\setlist{itemsep=2pt,topsep=3pt}
\renewcommand{\arraystretch}{1.15}
\titleformat{\part}[display]{\normalfont\huge\bfseries}{\partname~\thepart}{12pt}{\Huge}
\newcommand{\code}[1]{\texttt{\small #1}}
\newcommand{\param}[1]{\ensuremath{\mathsf{#1}}}
\hypersetup{
  pdftitle={577deputes.fr — Note methodologique},
  pdfauthor={Eric Martinez},
}

\title{\bfseries 577deputes.fr — Note méthodologique\\[4pt]
\large Formulation des indicateurs, analyse de sensibilité aux paramètres,\\
rapprochement à la littérature, limites}
\author{Eric Martinez\\\small \href{https://577deputes.fr}{577deputes.fr} ·
code source : \href{https://github.com/RogericBot/577deputes}{github.com/RogericBot/577deputes}}
\date{\today\\[2pt]\small Version 1.0 \quad·\quad DOI :
\href{https://doi.org/10.5281/zenodo.20127509}{\texttt{10.5281/zenodo.20127509}}}

\begin{document}
\maketitle
\thispagestyle{empty}

%------------------------------------------------------------------------------
% Résumés trilingues
%------------------------------------------------------------------------------
\begin{abstract}
\noindent\textbf{Résumé (français).} \emph{577deputes.fr} est un site web qui
ingère les jeux de données ouverts de l'Assemblée nationale française (XVII\textsuperscript{e}
législature, depuis juillet~2024), de l'INSEE et du ministère de l'Intérieur, et
en dérive une douzaine d'indicateurs : détection d'amendements quasi-identiques
(MinHash~+~LSH) avec typologie, cohésion et discipline de vote, matrice de
cohésion inter-groupes, diagramme ternaire des blocs, détection de réponses
ministérielles dupliquées, absentéisme stratégique, amendements jamais défendus,
délais de réponse, statistiques par circonscription. Ce document donne, pour
chaque indicateur~: (i)~sa définition formelle ; (ii)~l'algorithme et les valeurs
de paramètres retenues ; (iii)~une analyse de sensibilité à ces paramètres ;
(iv)~un rapprochement à la littérature existante ; (v)~ses limites. Il s'accompagne
d'un script de calcul reproductible. Le code de l'application est sous licence MIT ;
les données restent soumises à la Licence Ouverte~2.0 (Etalab) de leurs producteurs.

\medskip\noindent\textbf{Abstract (English).} \emph{577deputes.fr} is a website
that ingests the open datasets of the French National Assembly (17\textsuperscript{th}
legislature, since July~2024), INSEE and the Ministry of the Interior, and derives
about a dozen indicators: near-duplicate amendment detection (MinHash~+~LSH) with a
typology, vote cohesion and discipline, an inter-group cohesion matrix, a ternary
diagram of blocs, detection of duplicated ministerial answers, strategic absenteeism,
amendments never defended, response delays, constituency statistics. For each
indicator this document gives: (i)~its formal definition; (ii)~the algorithm and
the parameter values used; (iii)~a sensitivity analysis to those parameters;
(iv)~a comparison with the existing literature; (v)~its limitations. A reproducible
computation script is provided. The application code is MIT-licensed; the data
remain subject to their producers' \emph{Licence Ouverte~2.0} (Etalab).

\medskip\noindent\textbf{Resumen (español).} \emph{577deputes.fr} es un sitio web
que ingiere los conjuntos de datos abiertos de la Asamblea Nacional francesa
(XVII\textsuperscript{a} legislatura, desde julio de~2024), del INSEE y del
Ministerio del Interior, y deriva una docena de indicadores: detección de enmiendas
casi idénticas (MinHash~+~LSH) con tipología, cohesión y disciplina de voto, matriz
de cohesión entre grupos, diagrama ternario de bloques, detección de respuestas
ministeriales duplicadas, absentismo estratégico, enmiendas nunca defendidas,
plazos de respuesta, estadísticas por circunscripción. Para cada indicador este
documento ofrece: (i)~su definición formal; (ii)~el algoritmo y los valores de los
parámetros usados; (iii)~un análisis de sensibilidad a esos parámetros; (iv)~una
comparación con la literatura existente; (v)~sus límites. Se acompaña de un script
de cálculo reproducible. El código de la aplicación tiene licencia MIT; los datos
siguen sujetos a la \emph{Licence Ouverte~2.0} (Etalab) de sus productores.

\medskip\noindent\textbf{Mots-clés / Keywords / Palabras clave~:}
parlement, données ouvertes, votes nominaux, cohésion partisane, détection de
réutilisation de texte, MinHash, LSH, questions parlementaires, analyse de
sensibilité, reproductibilité.
\end{abstract}

%------------------------------------------------------------------------------
% Note sur le multilinguisme
%------------------------------------------------------------------------------
\section*{Note sur le choix trilingue \textnormal{\small/ On the trilingual choice / Sobre la elección trilingüe}}
\small
\textbf{FR.} Ce document est publié en français, en anglais et en espagnol. Ce
n'est pas l'usage pour une note méthodologique sur des données françaises ; nous
faisons ce choix délibérément. L'objet (l'activité parlementaire d'une démocratie)
et les sources (des données publiques) gagnent à être discutables par le plus grand
nombre — étudiants, journalistes, chercheurs, citoyens — sans barrière de langue.
Les trois versions visent à être autonomes et équivalentes ; en cas de doute, la
version française fait foi.

\textbf{EN.} This document is published in French, English and Spanish. This is
not the usual practice for a methodological note on French data; the choice is
deliberate. The subject (the parliamentary activity of a democracy) and the
sources (public data) deserve to be open to scrutiny by as many people as
possible — students, journalists, researchers, citizens — without a language
barrier. The three versions are meant to be self-contained and equivalent; in
case of doubt, the French version prevails.

\textbf{ES.} Este documento se publica en francés, inglés y español. No es la
práctica habitual para una nota metodológica sobre datos franceses; la elección
es deliberada. El tema (la actividad parlamentaria de una democracia) y las
fuentes (datos públicos) merecen poder ser examinados por el mayor número posible
de personas — estudiantes, periodistas, investigadores, ciudadanía — sin barrera
idiomática. Las tres versiones pretenden ser autónomas y equivalentes; en caso de
duda, prevalece la versión francesa.
\normalsize

\vfill
\noindent\rule{\linewidth}{0.4pt}\\
\small Ce document est déposé sur Zenodo —
\href{https://doi.org/10.5281/zenodo.20127509}{\texttt{doi:10.5281/zenodo.20127509}}.
Le script de calcul des analyses de sensibilité (\code{scripts/methodo\_sensitivity.py})
et les tableaux de sortie sont joints au dépôt ; la source LaTeX de ce document est
publique (\href{https://github.com/RogericBot/577deputes}{github.com/RogericBot/577deputes},
dossier \code{methodo/}). Toute erreur ou imprécision peut être signalée via le canal
de contact du site (\code{/.well-known/security.txt}).

\par\medskip\noindent\textbf{Pour citer ce document / How to cite / Para citar:}
Martinez,~E.\ (\the\year). \emph{577deputes.fr — Note méthodologique : formulation des
indicateurs, analyse de sensibilité aux paramètres, limites.} Zenodo.
\href{https://doi.org/10.5281/zenodo.20127509}{https://doi.org/10.5281/zenodo.20127509}
\normalsize

\clearpage
\tableofcontents
\clearpage

%==============================================================================
\part{Version française}
%==============================================================================

\section{Introduction et périmètre}

\subsection{Objet}
\emph{577deputes.fr} est une application web \emph{local-first} : elle télécharge
les jeux de données ouverts officiels, les ingère dans une base SQLite, et expose
des pages de consultation et des indicateurs dérivés pour la XVII\textsuperscript{e}
législature de l'Assemblée nationale française (en cours depuis juillet~2024).
L'application ne produit aucune donnée primaire ; tout ce qu'elle affiche est soit
une donnée brute reprise telle quelle d'une source publique, soit le résultat d'un
calcul documenté ici.

\subsection{Sources}
\begin{itemize}
  \item \textbf{Assemblée nationale} (\href{https://data.assemblee-nationale.fr}{data.assemblee-nationale.fr},
        Licence Ouverte~2.0) : acteurs, mandats et organes ; dossiers législatifs,
        documents et actes ; amendements ; scrutins publics nominatifs et votes
        individuels ; questions écrites, orales et au Gouvernement ; agenda et
        comptes rendus de séance.
  \item \textbf{INSEE} : population légale par circonscription législative.
  \item \textbf{Ministère de l'Intérieur} (data.gouv.fr) : inscrits, votants et
        abstentions par circonscription (élections législatives de 2024).
\end{itemize}
Les jeux de données de l'Assemblée sont rafraîchis quotidiennement par requête
conditionnelle (ETag / Last-Modified) ; un échec de téléchargement n'altère pas
les autres sources.

\subsection{Volumétrie (instantané ayant servi à ce document)}
\begin{center}\small
\begin{tabular}{@{}lr@{}}
\toprule
Députés indexés (dont 577 en exercice) & $\approx 630$ \\
Questions parlementaires & $\approx 17\,000$ \\
\quad dont avec réponse publiée $>200$ caractères & $10\,766$ \\
Amendements (XVII\textsuperscript{e} législature) & $108\,115$ \\
Scrutins publics nominatifs & $6\,490$ \\
Votes individuels & $1\,043\,305$ \\
Couples (groupe, scrutin) avec position majoritaire & $77\,748$ \\
\bottomrule
\end{tabular}
\end{center}

\subsection{Ce qui est exclu du périmètre}
\begin{itemize}
  \item Les \textbf{votes à main levée} et les votes en commission : le jeu de
        données \og{}Scrutins\fg{} ne contient que les scrutins publics nominatifs
        de séance publique. Tous les chiffres de discipline, de cohésion et
        d'absentéisme reposent sur ce périmètre-là (cf.~§\ref{sec:limites}).
  \item Les dossiers ouverts sous la XVI\textsuperscript{e} législature et
        poursuivis sous la XVII\textsuperscript{e} : ils sont indexés mais rattachés
        à leur législature d'origine, d'où de légers écarts entre compteurs.
  \item L'inférence d'intention : aucun indicateur de ce document ne prétend
        établir une motivation. \og{}Dépôt en masse\fg{}, \og{}amendement fantôme\fg{},
        \og{}absentéisme stratégique\fg{} décrivent des \emph{régularités mesurables},
        pas des desseins.
\end{itemize}

\section{Notations et conventions}\label{sec:notations}
On note $J(A,B) = |A\cap B| / |A\cup B|$ le coefficient de Jaccard de deux
ensembles. Pour un scrutin $v$, $\mathrm{pour}(v)$, $\mathrm{contre}(v)$,
$\mathrm{abst}(v)$ désignent les comptes officiels et
$\text{exprimés}(v)=\mathrm{pour}(v)+\mathrm{contre}(v)+\mathrm{abst}(v)$.
Pour un groupe parlementaire $G$ et un scrutin $v$, $p^\star(v,G)$ est la
\emph{position majoritaire} du groupe telle que publiée par l'Assemblée (champ
\code{positionMajoritaire}), à valeurs dans $\{\text{pour},\text{contre},\text{abstention}\}$.
Sur l'instantané utilisé, ce champ est renseigné pour la totalité des
$77\,748$ couples $(G,v)$ et ne prend jamais d'autre valeur (47,9~\% \og{}pour\fg{},
46,8~\% \og{}contre\fg{}, 5,3~\% \og{}abstention\fg{}) : il n'existe donc, à notre
niveau, aucun cas d'égalité parfaite à arbitrer.

Chaque sous-section d'indicateur suit le même plan : \textbf{(a)~définition},
\textbf{(b)~algorithme \& paramètres}, \textbf{(c)~sensibilité}, \textbf{(d)~littérature},
\textbf{(e)~limites}. Les chiffres de sensibilité proviennent du script
\code{methodo\_sensitivity.py} exécuté sur l'instantané ci-dessus ; les tableaux
complets sont en annexe~\ref{sec:annexes} et dans le dépôt.

\section{Les indicateurs}

%------------------------------------------------------------------------------
\subsection{Détection d'amendements quasi-identiques}\label{sec:clusters}

\paragraph{(a) Définition.} On cherche à regrouper les amendements dont le
\emph{dispositif} est \og{}essentiellement le même texte\fg{}, qu'ils aient ou non le
même auteur. Pour un amendement $a$, on construit $S_a$, l'ensemble des
$k$-\emph{shingles} de mots (séquences de $k$ mots consécutifs) du texte normalisé,
chacun haché sur 64~bits. Deux amendements $a,a'$ sont jugés quasi-identiques si
$J(S_a,S_{a'})\ge\theta$. Un \emph{cluster} est une composante connexe (au sens de
l'union de ces paires) de taille $\ge 2$.

\paragraph{(b) Algorithme \& paramètres.}
\begin{enumerate}[label=\arabic*.]
  \item \emph{Normalisation} : suppression des balises HTML ; décomposition
        Unicode (NFKD) et suppression des diacritiques ; passage en minuscules ;
        suppression de tout caractère hors \texttt{[a-z0-9 ]} ; normalisation des
        espaces ; troncature à $L_{\max}=12\,000$ caractères.
  \item \emph{Shinglage} : $k = 5$ ; $S_a=\{h(w_i\cdots w_{i+k-1})\}$ avec $h$ un
        hachage BLAKE2b 64~bits. Amendement écarté si $|S_a| < s_{\min}=8$.
  \item \emph{Signature MinHash} : $K=64$ fonctions de hachage
        $g_j(x)=x\oplus\sigma_j$ ($\sigma_j$ graines déterministes), signature
        $\mathrm{sig}_j(a)=\min_{x\in S_a} g_j(x)$.
  \item \emph{Blocage LSH} : la signature est découpée en $b=16$ bandes de $r=4$
        composantes ; $a$ et $a'$ sont \emph{candidats} s'ils coïncident sur au
        moins une bande. La probabilité qu'une paire de similarité réelle $s$
        devienne candidate est $P(s)=1-(1-s^{\,r})^{b}$ (courbe en~S, seuil
        $\approx (1/b)^{1/r}\approx 0{,}55$).
  \item \emph{Vérification} : pour chaque paire candidate, calcul du Jaccard exact
        sur $S_a$, $S_{a'}$ ; paire retenue si $J\ge\theta=0{,}80$.
  \item \emph{Clusterisation} : union-find sur les paires retenues ; clusters de
        taille $\ge 2$.
\end{enumerate}
Sur l'instantané : $62\,476$ amendements passent le seuil $s_{\min}$ ; on obtient
$\approx 9\,460$ clusters regroupant $28\,124$ amendements (la base de production en
recense $28\,055$ — l'écart de $0{,}25~\%$ vient d'un garde-fou de taille de
\emph{bucket} introduit dans le script de vérification, et constitue notre contrôle
de reproductibilité).

\paragraph{(c) Sensibilité.}
\emph{Seuil $\theta$} (jeu complet, LSH de production)~:
\begin{center}\small
\begin{tabular}{rrrr}
\toprule
$\theta$ & paires retenues & clusters & amendements clusterisés (\%) \\
\midrule
0,70 & 180\,407 & 9\,638 & 32\,953 \;(52,8~\%) \\
0,75 & 147\,438 & 9\,334 & 30\,794 \;(49,3~\%) \\
0,78 & \phantom{0}88\,007 & 9\,649 & 29\,159 \;(46,7~\%) \\
\textbf{0,80} & \phantom{0}\textbf{60\,381} & \textbf{9\,461} & \textbf{28\,124 \;(45,0~\%)} \\
0,82 & \phantom{0}54\,380 & 9\,253 & 26\,834 \;(43,0~\%) \\
0,85 & \phantom{0}43\,701 & 9\,023 & 25\,212 \;(40,4~\%) \\
0,90 & \phantom{0}34\,480 & 8\,440 & 22\,782 \;(36,5~\%) \\
0,95 & \phantom{0}30\,679 & 7\,840 & 20\,934 \;(33,5~\%) \\
\bottomrule
\end{tabular}
\end{center}
La sensibilité est \emph{modérée} : $\pm0{,}05$ autour de $0{,}80$ déplace le nombre
d'amendements clusterisés de $\approx \pm10$ à $\pm15~\%$. Mais la distribution des
similarités des paires candidates (annexe~\ref{sec:annexes}) montre que $0{,}80$
tombe dans une \emph{zone de transition} : il y a un fort amas de paires dans
$[0{,}75;0{,}85)$ (amendements qui partagent beaucoup de texte standard sans être
des copies), alors que les copies quasi-littérales forment un pic distinct au-delà
de $0{,}95$. Un seuil plus strict ($0{,}90$) ne retiendrait que les copies presque
verbatim ; un seuil plus lâche ($0{,}70$--$0{,}75$) absorberait aussi les
reformulations lourdes. Le choix de $0{,}80$ est un compromis assumé, non un point
\og{}naturel\fg{} du nuage. Pour $\theta<0{,}80$, le rappel des paires candidates
par le LSH 16$\times$4 décroît (voir ci-dessous) : les comptes affichés y sont des
bornes basses.

\emph{Taille des shingles $k$} (échantillon de $\approx 3\,500$ amendements,
Jaccard exact sur toutes les paires, $\theta=0{,}80$) : $k=3 \Rightarrow 12{,}4~\%$
d'amendements clusterisés ; $k=5 \Rightarrow 8{,}3~\%$ ; $k=7 \Rightarrow 6{,}5~\%$.
Des shingles plus courts sont plus permissifs ; $k=5$ est le réglage médian.

\emph{Configuration LSH} ($b\times r$) : $P(s)=1-(1-s^r)^b$. Pour le réglage de
production 16$\times$4 : $P(0{,}80)=0{,}9998$, $P(0{,}70)=0{,}988$, $P(0{,}50)=0{,}64$.
Le LSH n'affecte que le \emph{rappel} des paires candidates ; la \emph{précision}
finale est garantie par la vérification Jaccard exacte qui suit. Au seuil de
production, le rappel est donc quasi parfait.

\emph{Troncature $L_{\max}$ et plancher $s_{\min}$} : seuls $192$ amendements
($0{,}18~\%$) dépassent $12\,000$ caractères après normalisation ; relever
$s_{\min}$ de $8$ à $50$ exclurait $38~\%$ des amendements actuellement retenus
(la valeur $8$ est volontairement permissive, pour conserver les dispositifs
courts).

\paragraph{(d) Littérature.} La représentation par \emph{shingles} et l'estimation
de la similarité de Jaccard par \emph{min-wise hashing} sont dues à
Broder~\cite{broder1997,broder2000} ; le blocage par bandes (LSH) suit
Gionis-Indyk-Motwani~\cite{gionis1999} et l'exposé de
Leskovec-Rajaraman-Ullman~\cite{lru2014}. L'application à la détection de
réutilisation de texte législatif est explorée notamment par le \emph{Legislative
Influence Detector}~\cite{burgess2016} et les travaux sur la diffusion de \og{}model
bills\fg{} (Hertel-Fernandez~\cite{hf2014}, Jansa et al.~\cite{jansa2019},
Linder et al.~\cite{linder2020}). Notre apport n'est pas méthodologique : c'est
l'\emph{application} de briques connues au flux d'amendements de l'Assemblée, avec
une typologie descriptive.

\paragraph{(e) Limites.} (i)~La normalisation \emph{tronque} au-delà de $12\,000$
caractères : deux très longs amendements ne sont comparés que sur leur préfixe.
(ii)~Le shinglage de mots ignore l'ordre au-delà de la fenêtre $k$ : un
réagencement de paragraphes peut faire chuter le Jaccard. (iii)~Le seuil $0{,}80$
est dans la zone ambiguë du nuage : la frontière exacte des clusters est
conventionnelle. (iv)~Un \og{}cluster\fg{} ne dit rien d'une coordination
intentionnelle : un même modèle peut circuler via une fédération d'associations,
un syndicat, un site militant, sans concertation entre députés.

%------------------------------------------------------------------------------
\subsection{Typologie des clusters}\label{sec:typo}

\paragraph{(a) Définition.} Pour un cluster $C$, soit $n(C)=|C|$ sa taille et
$g(C)$ le nombre de groupes parlementaires distincts parmi ses amendements. On
classe~:
\[
\tau(C)=\begin{cases}
\text{dépôt en masse} & \text{si } n(C)\ge \param{T_{masse}}=15,\\[2pt]
\text{convergence inter-groupes} & \text{sinon si } g(C)\ge 2,\\[2pt]
\text{amplification intra-groupe} & \text{sinon si } g(C)\le 1 \text{ et } n(C)\ge \param{T_{ampl}}=3,\\[2pt]
\text{réutilisation simple} & \text{sinon.}
\end{cases}
\]
La taille prime toujours ; aucune intention n'est inférée (le libellé interne
\og{}obstruction\fg{} a été neutralisé en \og{}dépôt en masse\fg{}).

\paragraph{(b) Algorithme \& paramètres.} Calcul direct sur la table des clusters
($\param{T_{masse}}=15$, $\param{T_{ampl}}=3$). Sur l'instantané (clusters de
$\theta=0{,}80$) : $91$ clusters \og{}dépôt en masse\fg{} ($2\,578$ amendements),
$4\,314$ \og{}convergence\fg{} ($13\,834$), $773$ \og{}amplification\fg{}
($3\,146$), $4\,283$ \og{}réutilisation\fg{} ($8\,566$).

\paragraph{(c) Sensibilité.}
\begin{center}\small
\begin{tabular}{rr rr rr rr}
\toprule
$\param{T_{masse}}$ & $\param{T_{ampl}}$ & \multicolumn{2}{c}{dépôt masse} & \multicolumn{2}{c}{convergence} & \multicolumn{2}{c}{amplification} \\
 & & cl. & amdt & cl. & amdt & cl. & amdt \\
\midrule
10 & 3 & 203 & 3\,875 & 4\,226 & 12\,812 & 749 & 2\,871 \\
\textbf{15} & \textbf{3} & \textbf{91} & \textbf{2\,578} & \textbf{4\,314} & \textbf{13\,834} & \textbf{773} & \textbf{3\,146} \\
20 & 3 & 43 & 1\,776 & 4\,352 & 14\,472 & 783 & 3\,310 \\
30 & 3 & 19 & 1\,218 & 4\,370 & 14\,880 & 789 & 3\,460 \\
50 & 3 & \phantom{0}7 & \phantom{0,}784 & 4\,377 & 15\,121 & 794 & 3\,653 \\
\midrule
15 & 2 & 91 & 2\,578 & 4\,314 & 13\,834 & 5\,056 & 11\,712 \\
15 & 5 & 91 & 2\,578 & 4\,314 & 13\,834 & 160 & 1\,138 \\
\bottomrule
\end{tabular}
\end{center}
Le \emph{nombre} de clusters \og{}dépôt en masse\fg{} est sensible à
$\param{T_{masse}}$ (de $91$ à $7$ entre les seuils $15$ et $50$) mais le
\emph{volume} d'amendements concernés l'est moins ; le seuil $\param{T_{ampl}}$
fait surtout glisser des clusters entre \og{}amplification\fg{} et
\og{}réutilisation\fg{}. La répartition par taille est en annexe~\ref{sec:annexes}.

\paragraph{(d) Littérature.} La typologie est ad~hoc et descriptive ; elle
s'inscrit dans le champ de l'étude de l'\emph{obstruction parlementaire} et du
dépôt massif d'amendements comme tactique de débat (cf.~Koger~\cite{koger2010}
pour le cas américain ; le Règlement de l'Assemblée nationale encadre lui-même la
\og{}discussion commune\fg{} d'amendements identiques).

\paragraph{(e) Limites.} Quatre cases, des seuils ronds : un cluster de $14$
amendements identiques et un cluster de $16$ tombent de part et d'autre d'une
frontière qui n'a pas de fondement théorique. La case \og{}convergence\fg{} ne
distingue pas un consensus technique d'une alliance ponctuelle.

%------------------------------------------------------------------------------
\subsection{Discipline de vote (cohésion individuelle)}\label{sec:discipline}

\paragraph{(a) Définition.} Soit $G_d$ le groupe du député $d$. Soit $V_d$
l'ensemble des scrutins où $d$ a exprimé un vote $p_d(v)\in\{\text{pour},\text{contre},\text{abstention}\}$
\emph{et} où $p^\star(v,G_d)$ est défini. On pose
\[
\text{exprimés}(d)=|V_d|,\quad
\text{alignés}(d)=\big|\{v\in V_d : p_d(v)=p^\star(v,G_d)\}\big|,\quad
\mathrm{discipline}(d)=\frac{\text{alignés}(d)}{\text{exprimés}(d)}.
\]
Affichage dans les classements seulement si $\text{exprimés}(d)\ge 50$
(ou $\ge 100$ pour la liste des \og{}dissidents\fg{}, triée par discipline
croissante) ; les députés non inscrits (NI) en sont exclus.

\paragraph{(b) Algorithme \& paramètres.} Tout repose sur le champ officiel
$p^\star$ : nous ne calculons \emph{pas} la position majoritaire d'un groupe.
Sur l'instantané, $\mathrm{discipline}$ a une moyenne de $\approx 91{,}9~\%$, une
médiane de $\approx 92{,}5~\%$, un décile inférieur de $\approx 85~\%$.

\paragraph{(c) Sensibilité.} Effet du plancher d'exprimés~:
\begin{center}\small
\begin{tabular}{rrrrrr}
\toprule
plancher & éligibles & moyenne & médiane & $p_{10}$ & $p_{90}$ \\
\midrule
0   & 639 & 91,92 & 92,56 & 84,92 & 98,30 \\
20  & 628 & 91,87 & 92,50 & 84,93 & 98,25 \\
\textbf{50}  & \textbf{623} & \textbf{91,85} & \textbf{92,50} & \textbf{84,99} & \textbf{98,21} \\
100 & 611 & 91,83 & 92,45 & 85,02 & 98,18 \\
200 & 600 & 91,87 & 92,46 & 85,06 & 98,18 \\
\bottomrule
\end{tabular}
\end{center}
L'indicateur est \emph{très robuste} au plancher : la moyenne et les quantiles ne
bougent que de quelques centièmes de point entre les planchers $0$ et $200$. Le
choix $\ge 50$ n'est donc pas un paramètre sensible — il sert seulement à éviter
d'afficher un \og{}100~\%\fg{} reposant sur trois votes.

\paragraph{(d) Le \og{}biais circulaire\fg{}.} La position $p^\star(v,G_d)$ est la
pluralité des votes des membres de $G_d$, et $d$ en fait partie : $d$ est donc
comparé, en partie, à une position qu'il a contribué à définir. C'est un biais
\emph{réel} mais (i)~il n'est pas de notre fait — c'est la convention officielle de
l'Assemblée ; (ii)~pour un groupe de $50$ à $120$ membres, le poids d'un vote
individuel dans la pluralité est borné par $\le 2~\%$ ; (iii)~le cas d'une égalité
parfaite, qui rendrait la pluralité ambiguë, est arbitré en amont par l'Assemblée
et n'apparaît jamais dans les données (le champ $p^\star$ est toujours renseigné et
toujours dans $\{\text{pour},\text{contre},\text{abstention}\}$, cf.~§\ref{sec:notations}).

\paragraph{(e) Littérature.} La discipline ainsi définie est la version
\emph{individuelle} (par député) des indices de cohésion partisane : indice de
Rice~\cite{rice1925}, \emph{Agreement Index} de Hix-Noury-Roland~\cite{hnr2005,hnr2007},
discussions de Carey~\cite{carey2007} sur les \og{}principaux concurrents\fg{}.
Elle n'estime \emph{pas} de points idéaux : pour cela il faudrait des méthodes type
NOMINATE ou \emph{Optimal Classification} (Poole-Rosenthal~\cite{pr1997}, Poole~\cite{poole2005}),
hors périmètre.

\paragraph{(e') Limites.} (i)~Périmètre des seuls scrutins publics nominatifs :
la \og{}discipline\fg{} mesurée est celle des votes \emph{exprimés et nominatifs},
pas l'unité de groupe en général. (ii)~Une abstention conforme à la position du
groupe compte comme \og{}aligné\fg{}. (iii)~Un député souvent absent a un
$\text{exprimés}$ faible : sa discipline est mesurée sur peu de cas (d'où
le plancher).

%------------------------------------------------------------------------------
\subsection{Matrice de cohésion inter-groupes}\label{sec:matrice}

\paragraph{(a) Définition.} Pour deux groupes $G_1,G_2$, soit $V_{12}$ l'ensemble
des scrutins où les deux ont une position majoritaire définie. La cohésion est
\[
\mathrm{coh}(G_1,G_2)=\frac{\big|\{v\in V_{12}:p^\star(v,G_1)=p^\star(v,G_2)\}\big|}{|V_{12}|},
\qquad \mathrm{coh}(G,G)=1.
\]

\paragraph{(b) Algorithme \& paramètres.} Aucun paramètre libre ; auto-jointure
sur la table des positions majoritaires.

\paragraph{(c) Sensibilité.} Choix discutable : une \og{}abstention $=$ abstention\fg{}
est comptée comme accord. En recalculant la matrice sans les scrutins où les deux
groupes s'abstiennent, l'écart absolu moyen est de $0{,}21$~point et le maximum de
$0{,}83$~point. L'effet est donc \emph{négligeable}.

\paragraph{(d) Littérature.} C'est une matrice d'\emph{accord} entre blocs au sens
de Hix-Noury-Roland~\cite{hnr2007} ; elle ne fournit pas d'ordre unidimensionnel
(pour cela, voir les méthodes spatiales citées en §\ref{sec:discipline}).

\paragraph{(e) Limites.} Cohésion calculée sur \emph{positions de groupe}, pas sur
votes individuels : un groupe à $51/49$ et un groupe unanime comptent pareil. Deux
\og{}contre\fg{} pour des motifs opposés sont indistinguables d'un accord de fond.

%------------------------------------------------------------------------------
\subsection{Diagramme ternaire des blocs}\label{sec:ternaire}

\paragraph{(a) Définition.} Étant donné trois groupes \emph{ancres} $A_1,A_2,A_3$
placés aux sommets d'un triangle équilatéral en positions $\mathbf{x}_{A_1},\mathbf{x}_{A_2},\mathbf{x}_{A_3}$,
un groupe $G$ de cohésions $c_i=\mathrm{coh}(G,A_i)$ est placé en
\[
(\alpha,\beta,\gamma)=\frac{(c_1,c_2,c_3)}{c_1+c_2+c_3},\qquad
\mathbf{x}_G=\alpha\,\mathbf{x}_{A_1}+\beta\,\mathbf{x}_{A_2}+\gamma\,\mathbf{x}_{A_3}.
\]
Les ancres par défaut sont LFI, EPR, RN ; l'utilisateur peut en changer.

\paragraph{(b)--(c) Algorithme \& sensibilité.} Le seul \og{}paramètre\fg{} est le
choix des trois ancres : le diagramme n'est qu'une \emph{re-projection} de la
matrice de cohésion §\ref{sec:matrice} sur trois axes choisis, et change avec eux.

\paragraph{(d)--(e) Littérature \& limites.} Outil de visualisation, sans
prétention de mesure : la position d'un groupe \emph{dépend du repère}. Ne pas lire
les distances comme des distances idéologiques.

%------------------------------------------------------------------------------
\subsection{Coalitions par scrutin (par sujet)}\label{sec:coalitions}

\paragraph{(a) Définition.} On sélectionne les scrutins \og{}très suivis et
clivants\fg{} : (a)~$\text{exprimés}(v)\ge \param{T_v}=300$
($\approx 52~\%$ des 577 sièges) ; (b)~$|\mathrm{pour}(v)-\mathrm{contre}(v)|\le \param{T_e}=50$.
On trie par (b) croissant puis date décroissante et on garde les $N=10$ premiers ;
si moins de $N/2$ scrutins passent, on relâche le critère (b). Pour chaque scrutin
retenu, on affiche la position majoritaire de chaque groupe.

\paragraph{(b)--(c) Paramètres \& sensibilité.} Nombre de scrutins éligibles selon
$(\param{T_v},\param{T_e})$~:
\begin{center}\small
\begin{tabular}{r|rrrrr}
\toprule
$\param{T_v}\backslash\param{T_e}$ & 10 & 25 & 50 & 75 & 100 \\
\midrule
200 & \phantom{0}95 & 252 & 565 & 913 & 1\,200 \\
250 & \phantom{0}25 & \phantom{0}81 & 193 & 328 & \phantom{0,}440 \\
\textbf{300} & \phantom{0}\textbf{11} & \phantom{0}\textbf{36} & \phantom{00}\textbf{79} & 118 & \phantom{0,}155 \\
350 & \phantom{00}3 & \phantom{0}14 & \phantom{00}24 & \phantom{0}47 & \phantom{0,0}55 \\
400 & \phantom{00}2 & \phantom{00}7 & \phantom{00}11 & \phantom{0}19 & \phantom{0,0}20 \\
\bottomrule
\end{tabular}
\end{center}
Au réglage de production $(300,50)$, $79$ scrutins passent : le repli (b) n'est pas
activé. La sélection est sensible aux deux seuils, ce qui est attendu (c'est un
filtre, pas un estimateur).

\paragraph{(d)--(e) Littérature \& limites.} Approche purement descriptive : on
montre \emph{qui vote comme qui} sur les scrutins serrés, sans modèle. Limite
majeure : un scrutin \og{}serré en voix\fg{} peut l'être par faible affluence
asymétrique autant que par vrai clivage.

%------------------------------------------------------------------------------
\subsection{Réponses ministérielles \og{}types\fg{}}\label{sec:templates}

\paragraph{(a) Définition.} Pour chaque réponse écrite $q$ dont le texte brut
dépasse $200$ caractères, on normalise (suppression HTML, espaces), on retire
éventuellement la formule d'appel initiale, on passe en minuscules : $t_q$. La
réponse est écartée si $|t_q|<100$. La \emph{clé de préfixe} est
$\pi_L(q)=t_q[1:L]$ avec $L=250$. Une \og{}réponse type\fg{} est un préfixe partagé
par au moins $m=3$ réponses distinctes. Le \emph{taux de réponses types} est
\[
\rho \;=\; \frac{1}{|Q_{\text{élig}}|}\sum_{\pi \,:\, |\{q:\pi_L(q)=\pi\}|\ge m} \big|\{q:\pi_L(q)=\pi\}\big|.
\]
Sur l'instantané, $|Q_{\text{élig}}|=10\,766$ et $\rho=27{,}9~\%$.

\paragraph{(b)--(c) Paramètres \& sensibilité.} Grille du taux $\rho$ (en \%) selon
la longueur de préfixe $L$ et le seuil d'occurrences $m$ :
\begin{center}\small
\begin{tabular}{r|rrrr}
\toprule
$L\backslash m$ & 2 & \textbf{3} & 5 & 10 \\
\midrule
100  & 42,1 & 31,6 & 23,1 & 13,4 \\
\textbf{250}  & 37,0 & \textbf{27,9} & 20,3 & 11,6 \\
500  & 33,7 & 24,6 & 17,2 & \phantom{0}9,9 \\
1000 & 30,2 & 22,1 & 15,3 & \phantom{0}9,0 \\
\bottomrule
\end{tabular}
\end{center}
Le taux est \emph{fortement sensible au seuil $m$} (rapport $\approx 3$--$4$ entre
$m=2$ et $m=10$) et \emph{modérément sensible à $L$} (de $\approx 32~\%$ à
$\approx 22~\%$ à $m$ fixé). Sur ce corpus, le retrait de la formule d'appel n'a
\emph{aucun} effet (les variantes \og{}avec\fg{} et \og{}sans\fg{} retrait
coïncident exactement, le texte stocké ne commençant pas par ces formules). La
valeur affichée sur le site est donc à lire comme \emph{un} point d'une famille de
mesures, pas comme une constante : la formulation honnête est \og{}$\approx 28~\%$
des réponses partagent un préfixe de $250$~caractères avec au moins deux autres\fg{}.

\paragraph{(d) Littérature.} Les questions parlementaires comme instrument de
contrôle, et leurs réponses comme objet d'étude, sont traitées par
Wiberg~\cite{wiberg1995}, Russo-Wiberg~\cite{rw2010}, Martin~\cite{martin2011}.
Le repérage de réponses standardisées par regroupement de préfixes est une
heuristique \emph{maison}, voisine de la détection de \og{}boilerplate\fg{} ;
ce n'est pas une mesure validée de \og{}langue de bois\fg{}.

\paragraph{(e) Limites.} (i)~Un préfixe commun est une condition \emph{nécessaire}
mais non \emph{suffisante} d'une réponse \og{}vide\fg{} : deux réponses peuvent
partager un chapeau de cadrage puis diverger entièrement. (ii)~Le seuil $m=3$ et la
longueur $L=250$ sont arbitraires (cf.~grille). (iii)~Le corpus se limite aux
réponses publiées de plus de $200$ caractères ; les questions sans réponse n'entrent
pas dans le dénominateur de $\rho$.

%------------------------------------------------------------------------------
\subsection{Absentéisme stratégique}\label{sec:absence}

\paragraph{(a) Définition.} Parmi les scrutins à $\text{exprimés}(v)\ge \param{T_v}=200$,
on dit que $v$ est \emph{clivant} si
$|\mathrm{pour}(v)-\mathrm{contre}(v)|/\text{exprimés}(v)<\param{c}=0{,}10$,
\emph{consensuel} si $>\param{C}=0{,}50$, sinon non classé. Pour un député $d$, soit
$\mathrm{pres}_{cli}(d)$ (resp.\ $\mathrm{pres}_{con}(d)$) la \emph{part} de
scrutins clivants (resp.\ consensuels) où $d$ a exprimé un vote. On pose
$\mathrm{gap}(d)=\mathrm{pres}_{con}(d)-\mathrm{pres}_{cli}(d)$ (en points) et
$\mathrm{ratio}(d)=\mathrm{pres}_{con}(d)/\mathrm{pres}_{cli}(d)$. Affiché si $d$
est en exercice et $\mathrm{tot}_{cli}(d)\ge \param{T_{cli}}=30$.

\paragraph{(b)--(c) Paramètres \& sensibilité.} Quatre seuils ($\param{T_v},\param{c},\param{C},\param{T_{cli}}$).
Extrait de la grille (table complète en annexe~\ref{sec:annexes})~:
\begin{center}\small
\begin{tabular}{rrrr rrr}
\toprule
$\param{T_v}$ & $\param{c}$ & $\param{C}$ & $\param{T_{cli}}$ & $n_{cli}$ & $n_{con}$ & $\mathrm{gap}_{\max}$ (pts) \\
\midrule
200 & 0,10 & 0,50 & 30 & $\approx$ & $\approx$ & $\approx 20$--$30$ \\
\bottomrule
\end{tabular}
\end{center}
\emph{(Les valeurs exactes pour les $4^4$ combinaisons explorées figurent en
annexe ; on observe que la médiane du \emph{gap} est proche de $0$ — la majorité
des députés ont une présence semblable sur clivants et consensuels — et que seuls
les outliers ont un \emph{gap} marqué, dont l'amplitude varie de $\approx 13$ à
$\approx 31$ points selon les seuils.)}

\paragraph{(d)--(e) Littérature \& limites.} L'idée se rattache à la littérature
sur le \emph{shirking} législatif et la participation sélective aux votes. Limites
fortes : (i)~ne pas voter $\ne$ être absent — un député présent peut s'abstenir de
prendre part ; (ii)~\og{}clivant\fg{} est défini par l'\emph{issue} du vote, pas par
sa saillance politique \emph{a priori} ; (iii)~un \emph{gap} positif a de
nombreuses causes non stratégiques (agenda, maladie, circonscription lointaine).
C'est un signal à explorer, pas un verdict.

%------------------------------------------------------------------------------
\subsection{Amendements \og{}fantômes\fg{}}\label{sec:fantomes}

\paragraph{(a) Définition.} Pour le premier signataire $d$, sur l'ensemble de ses
amendements~:
\[
\begin{aligned}
\text{défendus}(d) &= \big|\{a : \mathrm{sort}(a)\in\{\text{Adopté, Rejeté, Retiré, Discuté}\}\}\big|,\\[2pt]
\text{fantômes}(d) &= \big|\{a : \mathrm{sort}(a)=\text{Non soutenu, ou sort vide}\}\big|,\\[2pt]
\text{pct\_fantômes}(d) &= \text{fantômes}(d)\,/\,\text{total}(d).
\end{aligned}
\]
Affiché si $\text{total}(d)\ge \param{T_a}=50$ ; trié par $\text{pct\_fantômes}$
décroissant.

\paragraph{(b)--(c) Paramètres \& sensibilité.} Recensement des codes \og{}sort\fg{}
(part du total des amendements) : Rejeté 22,9~\% ; Irrecevable 16,3~\% ; En
traitement 14,1~\% ; Adopté 13,0~\% ; Retiré 10,2~\% ; Discuté 9,0~\% ; Tombé
8,7~\% ; Non soutenu 5,9~\% ; Effacé $\approx$0~\%. (Aucun amendement n'a de sort
vide ou nul sur l'instantané.) Donc $\text{fantômes}$ se réduit en pratique à
\og{}Non soutenu\fg{} ($5{,}9~\%$), $\text{défendus}$ couvre $\approx 55~\%$,
et le reste (Irrecevable, En traitement, Tombé) n'est compté \emph{ni} comme
défendu \emph{ni} comme fantôme. Effet du plancher $\param{T_a}$~:
\begin{center}\small
\begin{tabular}{rrrrr}
\toprule
$\param{T_a}$ & auteurs éligibles & dont en exercice & $\overline{\mathrm{pct\_fant.}}$ & $\mathrm{pct\_fant.}^{\max}$ \\
\midrule
10  & 558 & 513 & 7,4 & 61,5 \\
20  & 530 & 491 & 7,2 & 52,4 \\
\textbf{50}  & \textbf{439} & \textbf{412} & \textbf{6,5} & \textbf{50,0} \\
100 & 313 & 300 & 6,6 & 50,0 \\
200 & 167 & 163 & 5,8 & 34,7 \\
\bottomrule
\end{tabular}
\end{center}
La moyenne du taux est stable ($\approx 6$--$7{,}5~\%$) ; le \emph{maximum} décroît
avec le plancher (les valeurs extrêmes sont portées par des auteurs à faible
volume). Une \emph{vue continue} (nuage de points $\mathrm{total}\times\text{pct\_fantômes}$,
sans plancher) est fournie dans le dépôt pour ne pas masquer cette dépendance.

\paragraph{(d)--(e) Littérature \& limites.} Pas de littérature directe ; on
quantifie une asymétrie entre \emph{dépôt} et \emph{défense} d'amendements.
Limites : (i)~la catégorie \og{}Non soutenu\fg{} recouvre aussi des contraintes de
temps en séance, indépendantes de la volonté de l'auteur ; (ii)~le choix d'inclure
\og{}Discuté\fg{} dans \og{}défendu\fg{} et d'exclure \og{}Tombé\fg{} est
discutable (la grille de sensibilité sur la définition de \og{}défendu\fg{} est dans
le dépôt) ; (iii)~\og{}premier signataire\fg{} $\ne$ \og{}auteur réel\fg{} dans le
cas d'amendements collectifs.

%------------------------------------------------------------------------------
\subsection{Délais de réponse ministériels et classements}\label{sec:delais}

\paragraph{(a) Définition.} Pour une question $q$, $\text{délai}(q)$ est le
nombre de jours entre dépôt et réponse (indéfini si non répondue). Pour un ministère
(au sens du \og{}ministère interrogé\fg{} court) $\mu$~:
\[
\overline{\text{délai}}(\mu)=\operatorname*{moyenne}_{q:\,\text{délai}(q)\ \text{défini}}\text{délai}(q),\qquad
\text{taux\_rép}(\mu)=\frac{|\{q:\text{statut}=\text{avec\_réponse}\}|}{|\{q\ \text{adressées à }\mu\}|}.
\]
Affiché si $|\{q\}|\ge \param{T_q}=30$ \emph{et} $|\{q:\text{statut}=\text{avec\_réponse}\}|\ge \param{T_r}=5$.
Le taux de réponse est affiché à côté du délai : un délai court n'a pas le même sens
si $1$~question sur $10$ seulement a obtenu une réponse.

\paragraph{(b)--(c) Paramètres \& sensibilité.} Nombre de ministères retenus selon
$(\param{T_q},\param{T_r})$ : $(10,1)\to 73$ ; $(20,*)\to 62$ ; $(30,*)\to 57$ ;
$(50,*)\to 52$ ; $(100,*)\to 40$. Le seuil $\param{T_r}$ est quasi inopérant (un
ministère recevant $\ge 30$ questions en a presque toujours répondu à $\ge 5$).

\paragraph{(d)--(e) Littérature \& limites.} Cf.~§\ref{sec:templates} pour les
questions comme objet d'étude. Limites : la moyenne est sensible aux valeurs
extrêmes (une réponse à $700$~jours tire la moyenne) ; le \og{}ministère interrogé\fg{}
peut différer du \og{}ministère attributaire\fg{} qui répond effectivement.

%------------------------------------------------------------------------------
\subsection{Indicateurs dérivés simples}\label{sec:simples}

\paragraph{(a) Activité cumulée.} $\mathrm{act}(d)=n_Q(d)+n_A(d)+n_V(d)$ (questions
posées $+$ amendements déposés $+$ votes exprimés). \emph{Avertissement structurel} :
$n_V$ est d'un ordre de grandeur supérieur à $n_Q$ et $n_A$ (un député assidu
exprime des milliers de votes mais dépose des dizaines d'amendements et pose des
dizaines de questions) ; l'indice cumulé est donc \emph{mécaniquement dominé par la
présence en scrutin}. Le site l'accompagne systématiquement des sous-classements
isolés (questions seules, amendements seuls, présence seule).

\paragraph{(b) Taux d'adoption des amendements.} $\text{adoptés}(d)/\mathrm{total}(d)$
(par député ou par groupe).

\paragraph{(c) Taux de réponse aux questions.} cf.~§\ref{sec:delais}.

\paragraph{(d) Âge, ancienneté.} Dérivés de la date de naissance et de la date de
début de mandat. Triviaux.

\paragraph{(e) Constructeur de tops paramétrables.} Le site propose une trentaine
de \og{}métriques\fg{} combinables à des filtres optionnels (par groupe, période,
type de texte, etc.). Toutes sont des \emph{ré-expositions} des primitives
ci-dessus (comptes, ratios, moyennes) ; il n'y a pas de méthode supplémentaire.

\paragraph{(f) Statistiques par circonscription.} $\mathrm{participation}=\mathrm{votants}/\mathrm{inscrits}$,
population (INSEE), inscrits (ministère de l'Intérieur), joints par
$(\text{code département},\ \text{numéro de circonscription})$. Couverture sur
l'instantané : $566$ circonscriptions avec population INSEE (les $11$
circonscriptions des Français de l'étranger n'en ont pas), $577$ avec inscrits.

\section{Reproductibilité}\label{sec:repro}
\begin{itemize}
  \item \textbf{Code de l'application} : public, sous licence MIT
        (\href{https://github.com/RogericBot/577deputes}{github.com/RogericBot/577deputes}).
        La base de données et les photographies ne sont pas redistribuées ; elles se
        reconstruisent avec la commande \code{anqp bootstrap} (téléchargement et
        ingestion des sources).
  \item \textbf{Caches dérivés} : les tables \code{deputy\_discipline\_cache},
        \code{groupe\_discipline\_cache}, \code{deputy\_amd\_cache},
        \code{groupe\_amd\_cache}, \code{scrutin\_groupes}, \code{amendement\_clusters},
        \code{circo\_stats} sont recalculées à l'ingestion ; elles ne contiennent
        rien qui ne dérive des sources.
  \item \textbf{Analyses de sensibilité} : le script \code{scripts/methodo\_sensitivity.py}
        (bibliothèque standard uniquement, lecture seule) reproduit tous les
        tableaux de ce document ; il est joint au dépôt Zenodo avec ses sorties
        (CSV et Markdown). Tous les tirages aléatoires utilisent une graine fixe.
  \item \textbf{Instantané} : les chiffres de ce document correspondent à une copie
        de la base prise pour sa rédaction (volumétrie en §1.3) ; ils peuvent
        différer légèrement de l'état courant du site, qui se met à jour
        quotidiennement.
\end{itemize}

\section{Limites générales et biais connus}\label{sec:limites}
\begin{enumerate}[label=L\arabic*.]
  \item \textbf{Périmètre des scrutins.} Seuls les scrutins publics nominatifs de
        séance publique existent dans les données. Les votes à main levée et les
        votes en commission n'y figurent pas. Tout indicateur de vote (discipline,
        cohésion, absentéisme, coalitions) hérite de cette limite. De plus,
        l'Assemblée vote très souvent avec une fraction seulement de l'hémicycle
        (médiane $\approx 150$ votants sur $577$) : un faible décompte reflète
        l'affluence du jour, pas un vote tronqué, et il n'y a pas de quorum requis
        pour la plupart des votes.
  \item \textbf{\og{}Auteur\fg{} d'un amendement.} Le champ utilisé est le
        \emph{premier signataire} ; pour un amendement collectif, il ne reflète pas
        nécessairement la rédaction.
  \item \textbf{Législatures à cheval.} Certains dossiers ouverts sous la
        XVI\textsuperscript{e} législature et poursuivis sous la XVII\textsuperscript{e}
        sont rattachés à leur législature d'origine ; d'où de légers écarts entre
        compteurs selon le périmètre retenu.
  \item \textbf{Caractère heuristique des typologies.} \og{}Dépôt en masse\fg{},
        \og{}convergence\fg{}, \og{}amplification\fg{}, \og{}réutilisation\fg{},
        \og{}réponse type\fg{} sont des étiquettes définies par des seuils ronds,
        pas des catégories théoriques.
  \item \textbf{Aucune inférence d'intention.} Les indicateurs décrivent des
        régularités mesurables. Ils ne disent pas \emph{pourquoi} : un dépôt massif
        peut enrichir le débat comme le ralentir ; un \emph{gap} de présence a
        souvent des causes prosaïques ; une réponse à préfixe commun peut être
        substantielle après ce préfixe.
  \item \textbf{Pas d'estimation de points idéaux.} Aucun indicateur ne produit
        d'échelle gauche-droite ni de carte spatiale ; la matrice de cohésion et le
        diagramme ternaire ne sont que des accords deux-à-deux re-projetés.
  \item \textbf{Dépendance aux producteurs de données.} Une lacune, un retard ou un
        changement de format chez l'Assemblée, l'INSEE ou le ministère de
        l'Intérieur se répercute directement ; le site signale explicitement quand
        une source est indisponible plutôt que d'inventer.
  \item \textbf{Instantané.} Le site se met à jour quotidiennement ; tout chiffre
        cité est daté du jour de son calcul.
\end{enumerate}

\section{Annexes}\label{sec:annexes}
Les tableaux complets ci-dessous sont également fournis en CSV/Markdown dans le
dépôt (un fichier par tableau, plus \code{SUMMARY.md}).
\begin{itemize}\small
  \item \textbf{A1} seuil Jaccard sur le jeu complet ; \textbf{A1bis} histogramme
        des similarités candidates ; \textbf{A2} $k$ et $\theta$ sur échantillon
        (Jaccard exact) ; \textbf{A3} courbe LSH $P(s)=1-(1-s^r)^b$ ; \textbf{A4}
        exposition à $L_{\max}$ ; \textbf{A4bis} plancher $s_{\min}$.
  \item \textbf{B} seuils de typologie ; \textbf{B2} distribution des tailles de
        cluster.
  \item \textbf{C0} valeurs de \code{positionMajoritaire} ; \textbf{C1} plancher de
        discipline.
  \item \textbf{D} matrice de cohésion, effet \og{}abstention partagée\fg{}.
  \item \textbf{F} seuils \og{}coalitions par scrutin\fg{}.
  \item \textbf{G} grille \og{}réponses types\fg{} ($L\times m\times$ retrait des
        formules).
  \item \textbf{H} grille \og{}absentéisme stratégique\fg{} ($\param{T_v}\times\param{c}\times\param{C}\times\param{T_{cli}}$).
  \item \textbf{I0} recensement des sorts ; \textbf{I1} plancher fantômes ;
        \textbf{I2} vue continue (nuage de points).
  \item \textbf{J} seuils du filtre ministères ; \textbf{K} couverture des
        statistiques de circonscription.
\end{itemize}

%==============================================================================
%==============================================================================
%  Part II — English version (full translation, mirrors Part I).
%==============================================================================
\clearpage
\selectlanguage{english}
\part{English version}
\setcounter{section}{0}   % each language part is numbered 1, 2, 3, … on its own

\section{Introduction and scope}

\subsection{Object}
\emph{577deputes.fr} is a \emph{local-first} web application: it downloads the
official open datasets, ingests them into a SQLite database, and exposes browsing
pages and derived indicators for the 17\textsuperscript{th} legislature of the
French National Assembly (ongoing since July~2024). The application produces no
primary data; everything it displays is either raw data taken as-is from a public
source, or the result of a computation documented here.

\subsection{Sources}
\begin{itemize}
  \item \textbf{National Assembly} (\href{https://data.assemblee-nationale.fr}{data.assemblee-nationale.fr},
        \emph{Licence Ouverte~2.0}): actors, mandates and bodies; legislative files,
        documents and procedural acts; amendments; recorded (nominal) public
        ballots and individual votes; written, oral and questions to the
        Government; sitting agenda and verbatim records.
  \item \textbf{INSEE}: legal population by legislative constituency.
  \item \textbf{Ministry of the Interior} (data.gouv.fr): registered voters,
        turnout and abstention by constituency (2024 legislative elections).
\end{itemize}
The Assembly datasets are refreshed daily via conditional requests
(ETag / Last-Modified); a download failure does not affect the other sources.

\subsection{Volume (snapshot used for this document)}
\begin{center}\small
\begin{tabular}{@{}lr@{}}
\toprule
Deputies indexed (of which 577 currently sitting) & $\approx 630$ \\
Parliamentary questions & $\approx 17\,000$ \\
\quad of which with a published answer $>200$ characters & $10\,766$ \\
Amendments (17\textsuperscript{th} legislature) & $108\,115$ \\
Recorded public ballots & $6\,490$ \\
Individual votes & $1\,043\,305$ \\
(group, ballot) pairs with a majority position & $77\,748$ \\
\bottomrule
\end{tabular}
\end{center}

\subsection{What is out of scope}
\begin{itemize}
  \item \textbf{Show-of-hands votes} and \textbf{committee votes}: the
        \og{}Scrutins\fg{} dataset only contains recorded nominal ballots held in
        the plenary chamber. All cohesion, discipline and absenteeism figures rest
        on that perimeter (cf.~§\ref{sec:limites-en}).
  \item \textbf{Files opened under the 16\textsuperscript{th} legislature} and
        continued under the 17\textsuperscript{th}: they are indexed but attached
        to their legislature of origin, hence small discrepancies between counters.
  \item \textbf{Intent inference}: no indicator in this document claims to
        establish a motive. \og{}Mass filing\fg{}, \og{}phantom amendment\fg{},
        \og{}strategic absenteeism\fg{} describe \emph{measurable regularities}, not
        designs.
\end{itemize}

\section{Notation and conventions}\label{sec:notations-en}
We write $J(A,B) = |A\cap B| / |A\cup B|$ for the Jaccard coefficient of two sets.
For a ballot $v$, $\mathrm{for}(v)$, $\mathrm{against}(v)$, $\mathrm{abst}(v)$ are
the official counts and $\text{cast}(v)=\mathrm{for}(v)+\mathrm{against}(v)+\mathrm{abst}(v)$.
For a parliamentary group $G$ and a ballot $v$, $p^\star(v,G)$ is the
\emph{majority position} of the group as published by the Assembly (field
\code{positionMajoritaire}), taking values in $\{\text{for},\text{against},\text{abstention}\}$.
On the snapshot used, this field is populated for all $77\,748$ pairs $(G,v)$ and
never takes any other value (47.9~\% \og{}for\fg{}, 46.8~\% \og{}against\fg{},
5.3~\% \og{}abstention\fg{}): there is therefore, at our level, no exact tie to
arbitrate.

Each indicator subsection follows the same plan: \textbf{(a)~definition},
\textbf{(b)~algorithm \& parameters}, \textbf{(c)~sensitivity}, \textbf{(d)~literature},
\textbf{(e)~limitations}. The sensitivity figures come from the script
\code{methodo\_sensitivity.py} run on the snapshot above; the full tables are in
Appendix~\ref{sec:annexes-en} and in the deposit.

\section{The indicators}

\subsection{Near-duplicate amendment detection}\label{sec:clusters-en}

\paragraph{(a) Definition.} We seek to group amendments whose operative text is
\og{}essentially the same\fg{}, whether or not they share an author. For an
amendment $a$ we build $S_a$, the set of word $k$-\emph{shingles} (sequences of
$k$ consecutive words) of the normalised text, each hashed to 64~bits. Two
amendments $a,a'$ are deemed near-identical if $J(S_a,S_{a'})\ge\theta$. A
\emph{cluster} is a connected component (under the union of these pairs) of size
$\ge 2$.

\paragraph{(b) Algorithm \& parameters.}
\begin{enumerate}[label=\arabic*.]
  \item \emph{Normalisation}: strip HTML tags; Unicode decomposition (NFKD) and
        removal of diacritics; lowercasing; removal of any character outside
        \texttt{[a-z0-9 ]}; whitespace collapse; truncation to $L_{\max}=12\,000$
        characters.
  \item \emph{Shingling}: $k = 5$; $S_a=\{h(w_i\cdots w_{i+k-1})\}$ with $h$ a
        64-bit BLAKE2b hash. Amendment discarded if $|S_a| < s_{\min}=8$.
  \item \emph{MinHash signature}: $K=64$ hash functions $g_j(x)=x\oplus\sigma_j$
        ($\sigma_j$ deterministic seeds), signature $\mathrm{sig}_j(a)=\min_{x\in S_a} g_j(x)$.
  \item \emph{LSH banding}: the signature is split into $b=16$ bands of $r=4$
        components; $a$ and $a'$ are \emph{candidates} if they coincide on at
        least one band. The probability that a pair of true similarity $s$ becomes
        a candidate is $P(s)=1-(1-s^{\,r})^{b}$ (S-curve, threshold
        $\approx (1/b)^{1/r}\approx 0.55$).
  \item \emph{Verification}: for each candidate pair, compute the exact Jaccard on
        $S_a$, $S_{a'}$; pair kept if $J\ge\theta=0.80$.
  \item \emph{Clustering}: union-find on the kept pairs; clusters of size $\ge 2$.
\end{enumerate}
On the snapshot: $62\,476$ amendments pass the $s_{\min}$ threshold; we obtain
$\approx 9\,460$ clusters grouping $28\,124$ amendments (the production database
records $28\,055$ — the $0.25~\%$ gap comes from a bucket-size guard introduced in
the verification script, and constitutes our reproducibility check).

\paragraph{(c) Sensitivity.}
\emph{Threshold $\theta$} (full set, production LSH):
\begin{center}\small
\begin{tabular}{rrrr}
\toprule
$\theta$ & pairs kept & clusters & clustered amendments (\%) \\
\midrule
0.70 & 180\,407 & 9\,638 & 32\,953 \;(52.8~\%) \\
0.75 & 147\,438 & 9\,334 & 30\,794 \;(49.3~\%) \\
0.78 & \phantom{0}88\,007 & 9\,649 & 29\,159 \;(46.7~\%) \\
\textbf{0.80} & \phantom{0}\textbf{60\,381} & \textbf{9\,461} & \textbf{28\,124 \;(45.0~\%)} \\
0.82 & \phantom{0}54\,380 & 9\,253 & 26\,834 \;(43.0~\%) \\
0.85 & \phantom{0}43\,701 & 9\,023 & 25\,212 \;(40.4~\%) \\
0.90 & \phantom{0}34\,480 & 8\,440 & 22\,782 \;(36.5~\%) \\
0.95 & \phantom{0}30\,679 & 7\,840 & 20\,934 \;(33.5~\%) \\
\bottomrule
\end{tabular}
\end{center}
Sensitivity is \emph{moderate}: $\pm0.05$ around $0.80$ shifts the number of
clustered amendments by $\approx\pm10$ to $\pm15~\%$. But the distribution of
candidate-pair similarities (Appendix~\ref{sec:annexes-en}) shows that $0.80$
falls in a \emph{transition region}: there is a large mass of pairs in
$[0.75;0.85)$ (amendments sharing much boilerplate without being copies), whereas
near-verbatim copies form a distinct peak above $0.95$. A stricter threshold
($0.90$) would keep only near-verbatim copies; a looser one ($0.70$--$0.75$) would
also absorb heavy rewrites. The choice of $0.80$ is a deliberate compromise, not a
\og{}natural\fg{} point of the cloud. For $\theta<0.80$, the recall of candidate
pairs by the 16$\times$4 LSH declines (see below): the displayed counts are then
lower bounds.

\emph{Shingle size $k$} (sample of $\approx 3\,500$ amendments, exact Jaccard on
all pairs, $\theta=0.80$): $k=3 \Rightarrow 12.4~\%$ clustered; $k=5 \Rightarrow
8.3~\%$; $k=7 \Rightarrow 6.5~\%$. Shorter shingles are more permissive; $k=5$ is
the median setting.

\emph{LSH configuration} ($b\times r$): $P(s)=1-(1-s^r)^b$. For the production
setting 16$\times$4: $P(0.80)=0.9998$, $P(0.70)=0.988$, $P(0.50)=0.64$. The LSH
only affects the \emph{recall} of candidate pairs; the final \emph{precision} is
guaranteed by the exact Jaccard verification that follows. At the production
threshold, recall is therefore near-perfect.

\emph{Truncation $L_{\max}$ and floor $s_{\min}$}: only $192$ amendments
($0.18~\%$) exceed $12\,000$ characters after normalisation; raising $s_{\min}$
from $8$ to $50$ would exclude $38~\%$ of the currently kept amendments (the value
$8$ is deliberately permissive, to keep short operative texts).

\paragraph{(d) Literature.} The shingle representation and the estimation of
Jaccard similarity by min-wise hashing are due to Broder~\cite{broder1997,broder2000};
band-based blocking (LSH) follows Gionis-Indyk-Motwani~\cite{gionis1999} and the
exposition of Leskovec-Rajaraman-Ullman~\cite{lru2014}. Application to detecting
reuse of legislative text is explored notably by the \emph{Legislative Influence
Detector}~\cite{burgess2016} and the work on the diffusion of \og{}model
bills\fg{} (Hertel-Fernandez~\cite{hf2014}, Jansa et al.~\cite{jansa2019},
Linder et al.~\cite{linder2020}). Our contribution is not methodological: it is
the \emph{application} of known building blocks to the Assembly's amendment stream,
with a descriptive typology.

\paragraph{(e) Limitations.} (i)~Normalisation \emph{truncates} beyond $12\,000$
characters: two very long amendments are compared only on their prefix.
(ii)~Word shingling ignores order beyond the window $k$: a reordering of
paragraphs can collapse the Jaccard. (iii)~The threshold $0.80$ is in the
ambiguous region of the cloud: the exact cluster boundary is conventional.
(iv)~A \og{}cluster\fg{} says nothing about intentional coordination: the same
template may circulate via an association federation, a union, an activist site,
without any concertation between deputies.

\subsection{Cluster typology}\label{sec:typo-en}

\paragraph{(a) Definition.} For a cluster $C$, let $n(C)=|C|$ be its size and
$g(C)$ the number of distinct parliamentary groups among its amendments. We
classify:
\[
\tau(C)=\begin{cases}
\text{mass filing} & \text{if } n(C)\ge \param{T_{masse}}=15,\\[2pt]
\text{inter-group convergence} & \text{otherwise if } g(C)\ge 2,\\[2pt]
\text{intra-group amplification} & \text{otherwise if } g(C)\le 1 \text{ and } n(C)\ge \param{T_{ampl}}=3,\\[2pt]
\text{simple reuse} & \text{otherwise.}
\end{cases}
\]
Size always takes precedence; no intent is inferred (the internal label
\og{}obstruction\fg{} was neutralised to \og{}mass filing\fg{}).

\paragraph{(b) Algorithm \& parameters.} Direct computation on the cluster table
($\param{T_{masse}}=15$, $\param{T_{ampl}}=3$). On the snapshot (clusters from
$\theta=0.80$): $91$ \og{}mass filing\fg{} clusters ($2\,578$ amendments),
$4\,314$ \og{}convergence\fg{} ($13\,834$), $773$ \og{}amplification\fg{}
($3\,146$), $4\,283$ \og{}reuse\fg{} ($8\,566$).

\paragraph{(c) Sensitivity.}
\begin{center}\small
\begin{tabular}{rr rr rr rr}
\toprule
$\param{T_{masse}}$ & $\param{T_{ampl}}$ & \multicolumn{2}{c}{mass filing} & \multicolumn{2}{c}{convergence} & \multicolumn{2}{c}{amplification} \\
 & & cl. & amdt & cl. & amdt & cl. & amdt \\
\midrule
10 & 3 & 203 & 3\,875 & 4\,226 & 12\,812 & 749 & 2\,871 \\
\textbf{15} & \textbf{3} & \textbf{91} & \textbf{2\,578} & \textbf{4\,314} & \textbf{13\,834} & \textbf{773} & \textbf{3\,146} \\
20 & 3 & 43 & 1\,776 & 4\,352 & 14\,472 & 783 & 3\,310 \\
30 & 3 & 19 & 1\,218 & 4\,370 & 14\,880 & 789 & 3\,460 \\
50 & 3 & \phantom{0}7 & \phantom{0,}784 & 4\,377 & 15\,121 & 794 & 3\,653 \\
\midrule
15 & 2 & 91 & 2\,578 & 4\,314 & 13\,834 & 5\,056 & 11\,712 \\
15 & 5 & 91 & 2\,578 & 4\,314 & 13\,834 & 160 & 1\,138 \\
\bottomrule
\end{tabular}
\end{center}
The \emph{number} of \og{}mass filing\fg{} clusters is sensitive to
$\param{T_{masse}}$ ($91$ down to $7$ between thresholds $15$ and $50$) but the
\emph{volume} of amendments affected is less so; the $\param{T_{ampl}}$ threshold
mostly shifts clusters between \og{}amplification\fg{} and \og{}reuse\fg{}. The
size distribution is in Appendix~\ref{sec:annexes-en}.

\paragraph{(d) Literature.} The typology is ad~hoc and descriptive; it belongs to
the field studying \emph{parliamentary obstruction} and mass amendment filing as a
debate tactic (cf.~Koger~\cite{koger2010} for the U.S. case; the Assembly's own
Rules of Procedure regulate the \og{}joint discussion\fg{} of identical
amendments).

\paragraph{(e) Limitations.} Four boxes, round thresholds: a cluster of $14$
identical amendments and one of $16$ fall on opposite sides of a boundary with no
theoretical basis. The \og{}convergence\fg{} box does not distinguish a technical
consensus from an opportunistic alliance.

\subsection{Vote discipline (individual cohesion)}\label{sec:discipline-en}

\paragraph{(a) Definition.} Let $G_d$ be the group of deputy $d$. Let $V_d$ be the
set of ballots where $d$ cast a vote $p_d(v)\in\{\text{for},\text{against},\text{abstention}\}$
\emph{and} where $p^\star(v,G_d)$ is defined. We set
\[
\mathrm{cast}(d)=|V_d|,\quad
\mathrm{aligned}(d)=\big|\{v\in V_d : p_d(v)=p^\star(v,G_d)\}\big|,\quad
\mathrm{discipline}(d)=\frac{\mathrm{aligned}(d)}{\mathrm{cast}(d)}.
\]
Shown in rankings only if $\mathrm{cast}(d)\ge 50$ (or $\ge 100$ for the
\og{}dissidents\fg{} list, sorted by ascending discipline); unaffiliated (NI)
deputies are excluded.

\paragraph{(b) Algorithm \& parameters.} Everything rests on the official field
$p^\star$: we do \emph{not} compute a group's majority position ourselves. On the
snapshot, $\mathrm{discipline}$ has a mean of $\approx 91.9~\%$, a median of
$\approx 92.5~\%$, a bottom decile of $\approx 85~\%$.

\paragraph{(c) Sensitivity.} Effect of the cast-vote floor:
\begin{center}\small
\begin{tabular}{rrrrrr}
\toprule
floor & eligible & mean & median & $p_{10}$ & $p_{90}$ \\
\midrule
0   & 639 & 91.92 & 92.56 & 84.92 & 98.30 \\
20  & 628 & 91.87 & 92.50 & 84.93 & 98.25 \\
\textbf{50}  & \textbf{623} & \textbf{91.85} & \textbf{92.50} & \textbf{84.99} & \textbf{98.21} \\
100 & 611 & 91.83 & 92.45 & 85.02 & 98.18 \\
200 & 600 & 91.87 & 92.46 & 85.06 & 98.18 \\
\bottomrule
\end{tabular}
\end{center}
The indicator is \emph{very robust} to the floor: the mean and quantiles move by
only hundredths of a point between floors $0$ and $200$. The choice $\ge 50$ is
therefore not a sensitive parameter — it only avoids displaying a \og{}100~\%\fg{}
based on three votes.

\paragraph{(d) The \og{}circular bias\fg{}.} The position $p^\star(v,G_d)$ is the
plurality of the votes of $G_d$'s members, and $d$ is one of them: $d$ is thus
compared, in part, to a position $d$ helped define. This is a \emph{real} bias but
(i)~it is not of our making — it is the Assembly's official convention;
(ii)~for a group of $50$ to $120$ members, the weight of one individual vote in
the plurality is bounded by $\le 2~\%$; (iii)~the case of an exact tie, which would
make the plurality ambiguous, is arbitrated upstream by the Assembly and never
appears in the data (the field $p^\star$ is always populated and always in
$\{\text{for},\text{against},\text{abstention}\}$, cf.~§\ref{sec:notations-en}).

\paragraph{(e) Literature.} The discipline so defined is the \emph{individual}
(per-deputy) version of party cohesion indices: Rice index~\cite{rice1925},
the Agreement Index of Hix-Noury-Roland~\cite{hnr2005,hnr2007}, Carey's
discussion~\cite{carey2007} of \og{}competing principals\fg{}. It does \emph{not}
estimate ideal points: for that one would need NOMINATE-type or \emph{Optimal
Classification} methods (Poole-Rosenthal~\cite{pr1997}, Poole~\cite{poole2005}),
out of scope.

\paragraph{(e') Limitations.} (i)~Perimeter restricted to recorded nominal
ballots: the \og{}discipline\fg{} measured is that of \emph{cast and recorded}
votes, not group unity in general. (ii)~An abstention consistent with the group's
position counts as \og{}aligned\fg{}. (iii)~An often-absent deputy has a low
$\mathrm{cast}$: their discipline is measured on few cases (hence the floor).

\subsection{Inter-group cohesion matrix}\label{sec:matrice-en}

\paragraph{(a) Definition.} For two groups $G_1,G_2$, let $V_{12}$ be the set of
ballots where both have a defined majority position. The cohesion is
\[
\mathrm{coh}(G_1,G_2)=\frac{\big|\{v\in V_{12}:p^\star(v,G_1)=p^\star(v,G_2)\}\big|}{|V_{12}|},
\qquad \mathrm{coh}(G,G)=1.
\]

\paragraph{(b) Algorithm \& parameters.} No free parameter; a self-join on the
majority-position table.

\paragraph{(c) Sensitivity.} A debatable choice: \og{}abstention $=$ abstention\fg{}
counts as agreement. Recomputing the matrix without the ballots where both groups
abstain, the mean absolute difference is $0.21$~point and the maximum $0.83$~point.
The effect is therefore \emph{negligible}.

\paragraph{(d) Literature.} It is an \emph{agreement} matrix between blocs in the
sense of Hix-Noury-Roland~\cite{hnr2007}; it provides no one-dimensional ordering
(for that, see the spatial methods cited in §\ref{sec:discipline-en}).

\paragraph{(e) Limitations.} Cohesion computed on \emph{group positions}, not
individual votes: a group split $51/49$ and a unanimous group count the same. Two
\og{}against\fg{} votes for opposite reasons are indistinguishable from genuine
agreement.

\subsection{Ternary diagram of blocs}\label{sec:ternaire-en}

\paragraph{(a) Definition.} Given three \emph{anchor} groups $A_1,A_2,A_3$ placed
at the vertices of an equilateral triangle at positions $\mathbf{x}_{A_1},\mathbf{x}_{A_2},\mathbf{x}_{A_3}$,
a group $G$ with cohesions $c_i=\mathrm{coh}(G,A_i)$ is placed at
\[
(\alpha,\beta,\gamma)=\frac{(c_1,c_2,c_3)}{c_1+c_2+c_3},\qquad
\mathbf{x}_G=\alpha\,\mathbf{x}_{A_1}+\beta\,\mathbf{x}_{A_2}+\gamma\,\mathbf{x}_{A_3}.
\]
The default anchors are LFI, EPR, RN; the user can change them.

\paragraph{(b)--(c) Algorithm \& sensitivity.} The only \og{}parameter\fg{} is the
choice of the three anchors: the diagram is merely a \emph{re-projection} of the
cohesion matrix §\ref{sec:matrice-en} onto three chosen axes, and changes with
them.

\paragraph{(d)--(e) Literature \& limitations.} A visualisation tool, with no claim
of measurement: a group's position \emph{depends on the frame of reference}. Do
not read the distances as ideological distances.

\subsection{Coalitions by ballot (by topic)}\label{sec:coalitions-en}

\paragraph{(a) Definition.} We select the \og{}highly attended and divisive\fg{}
ballots: (a)~$\text{cast}(v)\ge \param{T_v}=300$ ($\approx 52~\%$ of the 577 seats);
(b)~$|\mathrm{for}(v)-\mathrm{against}(v)|\le \param{T_e}=50$. We sort by (b)
ascending then date descending and keep the top $N=10$; if fewer than $N/2$
ballots pass, we relax criterion (b). For each retained ballot, we display each
group's majority position.

\paragraph{(b)--(c) Parameters \& sensitivity.} Number of eligible ballots by
$(\param{T_v},\param{T_e})$:
\begin{center}\small
\begin{tabular}{r|rrrrr}
\toprule
$\param{T_v}\backslash\param{T_e}$ & 10 & 25 & 50 & 75 & 100 \\
\midrule
200 & \phantom{0}95 & 252 & 565 & 913 & 1\,200 \\
250 & \phantom{0}25 & \phantom{0}81 & 193 & 328 & \phantom{0,}440 \\
\textbf{300} & \phantom{0}\textbf{11} & \phantom{0}\textbf{36} & \phantom{00}\textbf{79} & 118 & \phantom{0,}155 \\
350 & \phantom{00}3 & \phantom{0}14 & \phantom{00}24 & \phantom{0}47 & \phantom{0,0}55 \\
400 & \phantom{00}2 & \phantom{00}7 & \phantom{00}11 & \phantom{0}19 & \phantom{0,0}20 \\
\bottomrule
\end{tabular}
\end{center}
At the production setting $(300,50)$, $79$ ballots pass: the fallback (b) is not
triggered. The selection is sensitive to both thresholds, as expected (it is a
filter, not an estimator).

\paragraph{(d)--(e) Literature \& limitations.} A purely descriptive approach: we
show \emph{who votes like whom} on close ballots, with no model. Major limitation:
a ballot \og{}close in votes\fg{} may be so because of asymmetric low attendance as
much as because of a real cleavage.

\subsection{\og{}Template\fg{} ministerial answers}\label{sec:templates-en}

\paragraph{(a) Definition.} For each written answer $q$ whose raw text exceeds
$200$ characters, we normalise (strip HTML, collapse whitespace), optionally remove
the leading salutation, lowercase: $t_q$. The answer is discarded if $|t_q|<100$.
The \emph{prefix key} is $\pi_L(q)=t_q[1:L]$ with $L=250$. A \og{}template
answer\fg{} is a prefix shared by at least $m=3$ distinct answers. The
\emph{template rate} is
\[
\rho \;=\; \frac{1}{|Q_{\text{elig}}|}\sum_{\pi \,:\, |\{q:\pi_L(q)=\pi\}|\ge m} \big|\{q:\pi_L(q)=\pi\}\big|.
\]
On the snapshot, $|Q_{\text{elig}}|=10\,766$ and $\rho=27.9~\%$.

\paragraph{(b)--(c) Parameters \& sensitivity.} Grid of the rate $\rho$ (in \%) by
prefix length $L$ and occurrence threshold $m$:
\begin{center}\small
\begin{tabular}{r|rrrr}
\toprule
$L\backslash m$ & 2 & \textbf{3} & 5 & 10 \\
\midrule
100  & 42.1 & 31.6 & 23.1 & 13.4 \\
\textbf{250}  & 37.0 & \textbf{27.9} & 20.3 & 11.6 \\
500  & 33.7 & 24.6 & 17.2 & \phantom{0}9.9 \\
1000 & 30.2 & 22.1 & 15.3 & \phantom{0}9.0 \\
\bottomrule
\end{tabular}
\end{center}
The rate is \emph{strongly sensitive to the threshold $m$} (ratio $\approx 3$--$4$
between $m=2$ and $m=10$) and \emph{moderately sensitive to $L$} (from $\approx
32~\%$ to $\approx 22~\%$ at fixed $m$). On this corpus, removing the salutation has
\emph{no} effect (the \og{}with\fg{} and \og{}without\fg{} removal variants
coincide exactly, the stored text not starting with such formulas). The value
shown on the site should therefore be read as \emph{one} point of a family of
measures, not as a constant: the honest phrasing is \og{}$\approx 28~\%$ of answers
share a $250$-character prefix with at least two others\fg{}.

\paragraph{(d) Literature.} Parliamentary questions as an oversight instrument, and
their answers as an object of study, are treated by Wiberg~\cite{wiberg1995},
Russo-Wiberg~\cite{rw2010}, Martin~\cite{martin2011}. Spotting standardised answers
by prefix grouping is a \emph{home-grown} heuristic, akin to boilerplate detection;
it is not a validated measure of \og{}officialese\fg{}.

\paragraph{(e) Limitations.} (i)~A common prefix is a \emph{necessary} but not
\emph{sufficient} condition for an \og{}empty\fg{} answer: two answers may share a
framing header then diverge entirely. (ii)~The threshold $m=3$ and the length
$L=250$ are arbitrary (cf.~grid). (iii)~The corpus is limited to published answers
longer than $200$ characters; unanswered questions do not enter the denominator of
$\rho$.

\subsection{Strategic absenteeism}\label{sec:absence-en}

\paragraph{(a) Definition.} Among ballots with $\text{cast}(v)\ge \param{T_v}=200$,
$v$ is said \emph{divisive} if $|\mathrm{for}(v)-\mathrm{against}(v)|/\text{cast}(v)<\param{c}=0.10$,
\emph{consensual} if $>\param{C}=0.50$, otherwise unclassified. For a deputy $d$,
let $\mathrm{pres}_{cli}(d)$ (resp.\ $\mathrm{pres}_{con}(d)$) be the \emph{share}
of divisive (resp.\ consensual) ballots where $d$ cast a vote. We set
$\mathrm{gap}(d)=\mathrm{pres}_{con}(d)-\mathrm{pres}_{cli}(d)$ (in points) and
$\mathrm{ratio}(d)=\mathrm{pres}_{con}(d)/\mathrm{pres}_{cli}(d)$. Shown if $d$ is
currently sitting and $\mathrm{tot}_{cli}(d)\ge \param{T_{cli}}=30$.

\paragraph{(b)--(c) Parameters \& sensitivity.} Four thresholds
($\param{T_v},\param{c},\param{C},\param{T_{cli}}$). The exact values for the
$4^4$ explored combinations are in Appendix~\ref{sec:annexes-en}; one observes that
the median of the \emph{gap} is close to $0$ — most deputies have a similar
attendance on divisive and consensual ballots — and that only outliers have a
marked \emph{gap}, whose amplitude varies from $\approx 13$ to $\approx 31$ points
depending on the thresholds.

\paragraph{(d)--(e) Literature \& limitations.} The idea relates to the literature
on legislative \emph{shirking} and selective participation in votes. Strong
limitations: (i)~not voting $\ne$ being absent — a present deputy may decline to
take part; (ii)~\og{}divisive\fg{} is defined by the \emph{outcome} of the vote,
not by its political salience \emph{a priori}; (iii)~a positive \emph{gap} has many
non-strategic causes (agenda, illness, distant constituency). A signal to explore,
not a verdict.

\subsection{\og{}Phantom\fg{} amendments}\label{sec:fantomes-en}

\paragraph{(a) Definition.} For first signatory $d$, over all their amendments:
\[
\begin{aligned}
\mathrm{defended}(d) &= \big|\{a : \mathrm{sort}(a)\in\{\text{Adopted, Rejected, Withdrawn, Discussed}\}\}\big|,\\[2pt]
\mathrm{phantom}(d) &= \big|\{a : \mathrm{sort}(a)=\text{Not defended, or empty}\}\big|,\\[2pt]
\mathrm{pct\_phantom}(d) &= \mathrm{phantom}(d)\,/\,\mathrm{total}(d).
\end{aligned}
\]
Shown if $\mathrm{total}(d)\ge \param{T_a}=50$; sorted by descending
$\mathrm{pct\_phantom}$.

\paragraph{(b)--(c) Parameters \& sensitivity.} Census of \og{}sort\fg{} codes
(share of all amendments): Rejected 22.9~\%; Inadmissible 16.3~\%; Pending 14.1~\%;
Adopted 13.0~\%; Withdrawn 10.2~\%; Discussed 9.0~\%; Lapsed 8.7~\%; Not defended
5.9~\%; Erased $\approx$0~\%. (No amendment has an empty/null sort on the snapshot.)
So $\mathrm{phantom}$ reduces in practice to \og{}Not defended\fg{} ($5.9~\%$),
$\mathrm{defended}$ covers $\approx 55~\%$, and the rest (Inadmissible, Pending,
Lapsed) is counted \emph{neither} as defended \emph{nor} as phantom. Effect of the
floor $\param{T_a}$:
\begin{center}\small
\begin{tabular}{rrrrr}
\toprule
$\param{T_a}$ & eligible authors & of which sitting & $\overline{\mathrm{pct\_phant.}}$ & $\mathrm{pct\_phant.}^{\max}$ \\
\midrule
10  & 558 & 513 & 7.4 & 61.5 \\
20  & 530 & 491 & 7.2 & 52.4 \\
\textbf{50}  & \textbf{439} & \textbf{412} & \textbf{6.5} & \textbf{50.0} \\
100 & 313 & 300 & 6.6 & 50.0 \\
200 & 167 & 163 & 5.8 & 34.7 \\
\bottomrule
\end{tabular}
\end{center}
The mean rate is stable ($\approx 6$--$7.5~\%$); the \emph{maximum} decreases with
the floor (extreme values are carried by low-volume authors). A \emph{continuous
view} (scatter plot $\mathrm{total}\times\mathrm{pct\_phantom}$, without a floor) is
provided in the deposit so as not to hide this dependence.

\paragraph{(d)--(e) Literature \& limitations.} No direct literature; we quantify an
asymmetry between \emph{filing} and \emph{defending} amendments. Limitations:
(i)~the \og{}Not defended\fg{} category also covers time pressure in the chamber,
independent of the author's will; (ii)~the choice to include \og{}Discussed\fg{}
in \og{}defended\fg{} and to exclude \og{}Lapsed\fg{} is debatable (the sensitivity
grid on the definition of \og{}defended\fg{} is in the deposit); (iii)~\og{}first
signatory\fg{} $\ne$ \og{}actual author\fg{} for collective amendments.

\subsection{Ministerial response delays and rankings}\label{sec:delais-en}

\paragraph{(a) Definition.} For a question $q$, $\mathrm{delay}(q)$ is the number of
days between filing and answer (undefined if unanswered). For a ministry (in the
sense of the short \og{}questioned ministry\fg{}) $\mu$:
\[
\overline{\mathrm{delay}}(\mu)=\operatorname*{mean}_{q:\,\mathrm{delay}(q)\ \text{defined}}\mathrm{delay}(q),\qquad
\mathrm{response\_rate}(\mu)=\frac{|\{q:\text{status}=\text{answered}\}|}{|\{q\ \text{addressed to }\mu\}|}.
\]
Shown if $|\{q\}|\ge \param{T_q}=30$ \emph{and} $|\{q:\text{status}=\text{answered}\}|\ge \param{T_r}=5$.
The response rate is shown alongside the delay: a short delay does not mean the same
thing if only $1$ question in $10$ received an answer.

\paragraph{(b)--(c) Parameters \& sensitivity.} Number of ministries retained by
$(\param{T_q},\param{T_r})$: $(10,1)\to 73$; $(20,*)\to 62$; $(30,*)\to 57$;
$(50,*)\to 52$; $(100,*)\to 40$. The $\param{T_r}$ threshold is almost inoperative
(a ministry receiving $\ge 30$ questions has almost always answered $\ge 5$).

\paragraph{(d)--(e) Literature \& limitations.} Cf.~§\ref{sec:templates-en} for
questions as an object of study. Limitations: the mean is sensitive to extreme
values (one answer at $700$~days drags the mean up); the \og{}questioned
ministry\fg{} may differ from the \og{}assigned ministry\fg{} that actually
answers.

\subsection{Simple derived indicators}\label{sec:simples-en}

\paragraph{(a) Cumulative activity.} $\mathrm{act}(d)=n_Q(d)+n_A(d)+n_V(d)$
(questions asked $+$ amendments filed $+$ votes cast). \emph{Structural warning}:
$n_V$ is an order of magnitude larger than $n_Q$ and $n_A$ (a diligent deputy casts
thousands of votes but files dozens of amendments and asks dozens of questions); the
cumulative index is therefore \emph{mechanically dominated by attendance at
ballots}. The site systematically accompanies it with the isolated sub-rankings
(questions only, amendments only, attendance only).

\paragraph{(b) Amendment adoption rate.} $\mathrm{adopted}(d)/\mathrm{total}(d)$
(per deputy or per group).

\paragraph{(c) Question response rate.} cf.~§\ref{sec:delais-en}.

\paragraph{(d) Age, seniority.} Derived from the date of birth and the mandate
start date. Trivial.

\paragraph{(e) Parametrisable ranking builder.} The site offers some thirty
\og{}metrics\fg{} combinable with optional filters (by group, period, text type,
etc.). They are all \emph{re-expositions} of the primitives above (counts, ratios,
means); there is no additional method.

\paragraph{(f) Constituency statistics.} $\mathrm{turnout}=\mathrm{voters}/\mathrm{registered}$,
population (INSEE), registered voters (Ministry of the Interior), joined by
$(\text{department code},\ \text{constituency number})$. Coverage on the snapshot:
$566$ constituencies with INSEE population (the $11$ constituencies for French
citizens abroad have none), $577$ with registered voters.

\section{Reproducibility}\label{sec:repro-en}
\begin{itemize}
  \item \textbf{Application code}: public, MIT-licensed
        (\href{https://github.com/RogericBot/577deputes}{github.com/RogericBot/577deputes}).
        The database and the photographs are not redistributed; they rebuild with
        the command \code{anqp bootstrap} (download and ingest the sources).
  \item \textbf{Derived caches}: the tables \code{deputy\_discipline\_cache},
        \code{groupe\_discipline\_cache}, \code{deputy\_amd\_cache},
        \code{groupe\_amd\_cache}, \code{scrutin\_groupes}, \code{amendement\_clusters},
        \code{circo\_stats} are recomputed at ingestion time; they contain nothing
        that does not derive from the sources.
  \item \textbf{Sensitivity analyses}: the script \code{scripts/methodo\_sensitivity.py}
        (standard library only, read-only) reproduces all the tables in this
        document; it is attached to the Zenodo deposit with its outputs (CSV and
        Markdown). All random draws use a fixed seed.
  \item \textbf{Snapshot}: the figures in this document correspond to a copy of the
        database taken for its writing (volume in §1.3); they may differ slightly
        from the current state of the site, which updates daily.
\end{itemize}

\section{General limitations and known biases}\label{sec:limites-en}
\begin{enumerate}[label=L\arabic*.]
  \item \textbf{Ballot perimeter.} Only recorded nominal ballots held in the
        plenary chamber exist in the data. Show-of-hands votes and committee votes
        are not there. Every vote indicator (discipline, cohesion, absenteeism,
        coalitions) inherits this limitation. Moreover, the Assembly very often
        votes with only a fraction of the chamber present (median $\approx 150$
        voters out of $577$): a low count reflects that day's attendance, not a
        truncated vote, and there is no quorum required for most votes.
  \item \textbf{Amendment \og{}author\fg{}.} The field used is the \emph{first
        signatory}; for a collective amendment it does not necessarily reflect the
        drafting.
  \item \textbf{Straddling legislatures.} Some files opened under the
        16\textsuperscript{th} legislature and continued under the
        17\textsuperscript{th} are attached to their legislature of origin; hence
        small discrepancies between counters depending on the perimeter chosen.
  \item \textbf{Heuristic nature of the typologies.} \og{}Mass filing\fg{},
        \og{}convergence\fg{}, \og{}amplification\fg{}, \og{}reuse\fg{},
        \og{}template answer\fg{} are labels defined by round thresholds, not
        theoretical categories.
  \item \textbf{No intent inference.} The indicators describe measurable
        regularities. They do not say \emph{why}: a mass filing may enrich the
        debate or slow it down; an attendance \emph{gap} often has prosaic causes;
        an answer with a common prefix may be substantive after that prefix.
  \item \textbf{No ideal-point estimation.} No indicator produces a left-right scale
        or a spatial map; the cohesion matrix and the ternary diagram are only
        re-projected pairwise agreements.
  \item \textbf{Dependence on the data producers.} A gap, a delay or a format change
        at the Assembly, INSEE or the Ministry of the Interior propagates directly;
        the site explicitly signals when a source is unavailable rather than
        inventing.
  \item \textbf{Snapshot.} The site updates daily; any cited figure is dated by the
        day it was computed.
\end{enumerate}

\section{Appendices}\label{sec:annexes-en}
The full tables below are also provided in CSV/Markdown in the deposit (one file per
table, plus \code{SUMMARY.md}).
\begin{itemize}\small
  \item \textbf{A1} Jaccard threshold on the full set; \textbf{A1bis} histogram of
        candidate similarities; \textbf{A2} $k$ and $\theta$ on a sample (exact
        Jaccard); \textbf{A3} LSH curve $P(s)=1-(1-s^r)^b$; \textbf{A4} exposure to
        $L_{\max}$; \textbf{A4bis} floor $s_{\min}$.
  \item \textbf{B} typology thresholds; \textbf{B2} cluster size distribution.
  \item \textbf{C0} values of \code{positionMajoritaire}; \textbf{C1} discipline
        floor.
  \item \textbf{D} cohesion matrix, \og{}shared abstention\fg{} effect.
  \item \textbf{F} \og{}coalitions by ballot\fg{} thresholds.
  \item \textbf{G} \og{}template answers\fg{} grid ($L\times m\times$ salutation
        removal).
  \item \textbf{H} \og{}strategic absenteeism\fg{} grid ($\param{T_v}\times\param{c}\times\param{C}\times\param{T_{cli}}$).
  \item \textbf{I0} census of sorts; \textbf{I1} phantom floor; \textbf{I2}
        continuous view (scatter plot).
  \item \textbf{J} ministry-filter thresholds; \textbf{K} constituency-statistics
        coverage.
\end{itemize}
   % Partie II — English version (squelette à ce stade)
%==============================================================================
%  Parte III — Versión en español (traducción completa, refleja la Parte I).
%==============================================================================
\clearpage
\selectlanguage{spanish}
\part{Versión en español}
\setcounter{section}{0}   % cada parte lingüística se numera 1, 2, 3, … por su cuenta

\section{Introducción y alcance}

\subsection{Objeto}
\emph{577deputes.fr} es una aplicación web \emph{local-first}: descarga los
conjuntos de datos abiertos oficiales, los ingiere en una base de datos SQLite y
expone páginas de consulta e indicadores derivados para la XVII\textsuperscript{a}
legislatura de la Asamblea Nacional francesa (en curso desde julio de 2024). La
aplicación no produce ningún dato primario; todo lo que muestra es, o bien un dato
en bruto tomado tal cual de una fuente pública, o bien el resultado de un cálculo
documentado aquí.

\subsection{Fuentes}
\begin{itemize}
  \item \textbf{Asamblea Nacional} (\href{https://data.assemblee-nationale.fr}{data.assemblee-nationale.fr},
        \emph{Licence Ouverte~2.0}): actores, mandatos y órganos; expedientes
        legislativos, documentos y actos de procedimiento; enmiendas; votaciones
        públicas nominales y votos individuales; preguntas escritas, orales y al
        Gobierno; orden del día de las sesiones y actas literales.
  \item \textbf{INSEE}: población legal por circunscripción legislativa.
  \item \textbf{Ministerio del Interior} (data.gouv.fr): inscritos, votantes y
        abstención por circunscripción (elecciones legislativas de 2024).
\end{itemize}
Los conjuntos de datos de la Asamblea se actualizan a diario mediante peticiones
condicionales (ETag / Last-Modified); un fallo de descarga no altera las demás
fuentes.

\subsection{Volumen (instantánea utilizada para este documento)}
\begin{center}\small
\begin{tabular}{@{}lr@{}}
\toprule
Diputados indexados (de los cuales 577 en ejercicio) & $\approx 630$ \\
Preguntas parlamentarias & $\approx 17\,000$ \\
\quad de las cuales con respuesta publicada $>200$ caracteres & $10\,766$ \\
Enmiendas (XVII\textsuperscript{a} legislatura) & $108\,115$ \\
Votaciones públicas nominales & $6\,490$ \\
Votos individuales & $1\,043\,305$ \\
Pares (grupo, votación) con posición mayoritaria & $77\,748$ \\
\bottomrule
\end{tabular}
\end{center}

\subsection{Lo que queda fuera del alcance}
\begin{itemize}
  \item \textbf{Las votaciones a mano alzada} y \textbf{las votaciones en
        comisión}: el conjunto de datos \og{}Scrutins\fg{} solo contiene las
        votaciones públicas nominales celebradas en el pleno. Todas las cifras de
        cohesión, disciplina y absentismo se apoyan en ese perímetro
        (cf.~§\ref{sec:limites-es}).
  \item \textbf{Los expedientes abiertos bajo la XVI\textsuperscript{a}
        legislatura} y continuados bajo la XVII\textsuperscript{a}: están indexados
        pero asignados a su legislatura de origen, de ahí pequeñas discrepancias
        entre contadores.
  \item \textbf{La inferencia de intención}: ningún indicador de este documento
        pretende establecer un motivo. \og{}Presentación masiva\fg{}, \og{}enmienda
        fantasma\fg{}, \og{}absentismo estratégico\fg{} describen \emph{regularidades
        medibles}, no designios.
\end{itemize}

\section{Notación y convenciones}\label{sec:notations-es}
Escribimos $J(A,B) = |A\cap B| / |A\cup B|$ para el coeficiente de Jaccard de dos
conjuntos. Para una votación $v$, $\mathrm{a\_favor}(v)$, $\mathrm{en\_contra}(v)$,
$\mathrm{abst}(v)$ son los recuentos oficiales y
$\text{emitidos}(v)=\mathrm{a\_favor}(v)+\mathrm{en\_contra}(v)+\mathrm{abst}(v)$.
Para un grupo parlamentario $G$ y una votación $v$, $p^\star(v,G)$ es la
\emph{posición mayoritaria} del grupo tal como la publica la Asamblea (campo
\code{positionMajoritaire}), con valores en $\{\text{a favor},\text{en contra},\text{abstención}\}$.
En la instantánea utilizada, ese campo está cumplimentado para los $77\,748$ pares
$(G,v)$ y nunca toma otro valor (47,9~\% \og{}a favor\fg{}, 46,8~\% \og{}en
contra\fg{}, 5,3~\% \og{}abstención\fg{}): por lo tanto no existe, a nuestro nivel,
ningún empate exacto que arbitrar.

Cada subsección de indicador sigue el mismo esquema: \textbf{(a)~definición},
\textbf{(b)~algoritmo y parámetros}, \textbf{(c)~sensibilidad}, \textbf{(d)~literatura},
\textbf{(e)~límites}. Las cifras de sensibilidad provienen del script
\code{methodo\_sensitivity.py} ejecutado sobre la instantánea anterior; las tablas
completas están en el anexo~\ref{sec:annexes-es} y en el depósito.

\section{Los indicadores}

\subsection{Detección de enmiendas casi idénticas}\label{sec:clusters-es}

\paragraph{(a) Definición.} Se busca agrupar las enmiendas cuyo texto dispositivo
es \og{}esencialmente el mismo\fg{}, compartan o no autor. Para una enmienda $a$ se
construye $S_a$, el conjunto de los $k$-\emph{shingles} de palabras (secuencias de
$k$ palabras consecutivas) del texto normalizado, cada uno troceado a 64~bits. Dos
enmiendas $a,a'$ se consideran casi idénticas si $J(S_a,S_{a'})\ge\theta$. Un
\emph{clúster} es una componente conexa (en el sentido de la unión de esos pares)
de tamaño $\ge 2$.

\paragraph{(b) Algoritmo y parámetros.}
\begin{enumerate}[label=\arabic*.]
  \item \emph{Normalización}: supresión de etiquetas HTML; descomposición Unicode
        (NFKD) y supresión de diacríticos; paso a minúsculas; supresión de todo
        carácter fuera de \texttt{[a-z0-9 ]}; normalización de espacios; truncado a
        $L_{\max}=12\,000$ caracteres.
  \item \emph{Shingling}: $k = 5$; $S_a=\{h(w_i\cdots w_{i+k-1})\}$ con $h$ un
        hash BLAKE2b de 64~bits. Enmienda descartada si $|S_a| < s_{\min}=8$.
  \item \emph{Firma MinHash}: $K=64$ funciones hash $g_j(x)=x\oplus\sigma_j$
        ($\sigma_j$ semillas deterministas), firma $\mathrm{sig}_j(a)=\min_{x\in S_a} g_j(x)$.
  \item \emph{Bloqueo LSH}: la firma se divide en $b=16$ bandas de $r=4$
        componentes; $a$ y $a'$ son \emph{candidatos} si coinciden en al menos una
        banda. La probabilidad de que un par de similitud real $s$ sea candidato es
        $P(s)=1-(1-s^{\,r})^{b}$ (curva en S, umbral $\approx (1/b)^{1/r}\approx 0{,}55$).
  \item \emph{Verificación}: para cada par candidato, cálculo del Jaccard exacto
        sobre $S_a$, $S_{a'}$; par retenido si $J\ge\theta=0{,}80$.
  \item \emph{Agrupamiento}: union-find sobre los pares retenidos; clústeres de
        tamaño $\ge 2$.
\end{enumerate}
En la instantánea: $62\,476$ enmiendas pasan el umbral $s_{\min}$; se obtienen
$\approx 9\,460$ clústeres que agrupan $28\,124$ enmiendas (la base de producción
registra $28\,055$ — la diferencia de $0{,}25~\%$ proviene de una salvaguarda de
tamaño de \emph{bucket} introducida en el script de verificación, y constituye
nuestro control de reproducibilidad).

\paragraph{(c) Sensibilidad.}
\emph{Umbral $\theta$} (conjunto completo, LSH de producción):
\begin{center}\small
\begin{tabular}{rrrr}
\toprule
$\theta$ & pares retenidos & clústeres & enmiendas agrupadas (\%) \\
\midrule
0,70 & 180\,407 & 9\,638 & 32\,953 \;(52,8~\%) \\
0,75 & 147\,438 & 9\,334 & 30\,794 \;(49,3~\%) \\
0,78 & \phantom{0}88\,007 & 9\,649 & 29\,159 \;(46,7~\%) \\
\textbf{0,80} & \phantom{0}\textbf{60\,381} & \textbf{9\,461} & \textbf{28\,124 \;(45,0~\%)} \\
0,82 & \phantom{0}54\,380 & 9\,253 & 26\,834 \;(43,0~\%) \\
0,85 & \phantom{0}43\,701 & 9\,023 & 25\,212 \;(40,4~\%) \\
0,90 & \phantom{0}34\,480 & 8\,440 & 22\,782 \;(36,5~\%) \\
0,95 & \phantom{0}30\,679 & 7\,840 & 20\,934 \;(33,5~\%) \\
\bottomrule
\end{tabular}
\end{center}
La sensibilidad es \emph{moderada}: $\pm0{,}05$ alrededor de $0{,}80$ desplaza el
número de enmiendas agrupadas en $\approx\pm10$ a $\pm15~\%$. Pero la distribución
de las similitudes de los pares candidatos (anexo~\ref{sec:annexes-es}) muestra que
$0{,}80$ cae en una \emph{zona de transición}: hay una gran masa de pares en
$[0{,}75;0{,}85)$ (enmiendas que comparten mucho texto estándar sin ser copias),
mientras que las copias casi literales forman un pico distinto por encima de
$0{,}95$. Un umbral más estricto ($0{,}90$) solo retendría las copias casi
verbatim; uno más laxo ($0{,}70$--$0{,}75$) absorbería también las reescrituras
pesadas. La elección de $0{,}80$ es un compromiso asumido, no un punto
\og{}natural\fg{} de la nube. Para $\theta<0{,}80$, la cobertura de los pares
candidatos por el LSH 16$\times$4 disminuye (véase abajo): los recuentos mostrados
ahí son cotas inferiores.

\emph{Tamaño de los shingles $k$} (muestra de $\approx 3\,500$ enmiendas, Jaccard
exacto sobre todos los pares, $\theta=0{,}80$): $k=3 \Rightarrow 12{,}4~\%$
agrupadas; $k=5 \Rightarrow 8{,}3~\%$; $k=7 \Rightarrow 6{,}5~\%$. Shingles más
cortos son más permisivos; $k=5$ es el ajuste mediano.

\emph{Configuración LSH} ($b\times r$): $P(s)=1-(1-s^r)^b$. Para el ajuste de
producción 16$\times$4: $P(0{,}80)=0{,}9998$, $P(0{,}70)=0{,}988$, $P(0{,}50)=0{,}64$.
El LSH solo afecta a la \emph{cobertura} de los pares candidatos; la
\emph{precisión} final está garantizada por la verificación Jaccard exacta que
sigue. En el umbral de producción, la cobertura es por tanto casi perfecta.

\emph{Truncado $L_{\max}$ y umbral mínimo $s_{\min}$}: solo $192$ enmiendas
($0{,}18~\%$) superan los $12\,000$ caracteres tras la normalización; elevar
$s_{\min}$ de $8$ a $50$ excluiría el $38~\%$ de las enmiendas actualmente retenidas
(el valor $8$ es deliberadamente permisivo, para conservar los textos dispositivos
cortos).

\paragraph{(d) Literatura.} La representación por \emph{shingles} y la estimación de
la similitud de Jaccard mediante min-wise hashing se deben a
Broder~\cite{broder1997,broder2000}; el bloqueo por bandas (LSH) sigue a
Gionis-Indyk-Motwani~\cite{gionis1999} y la exposición de
Leskovec-Rajaraman-Ullman~\cite{lru2014}. La aplicación a la detección de
reutilización de texto legislativo está explorada notablemente por el
\emph{Legislative Influence Detector}~\cite{burgess2016} y los trabajos sobre la
difusión de \og{}model bills\fg{} (Hertel-Fernandez~\cite{hf2014},
Jansa et al.~\cite{jansa2019}, Linder et al.~\cite{linder2020}). Nuestro aporte no
es metodológico: es la \emph{aplicación} de piezas conocidas al flujo de enmiendas
de la Asamblea, con una tipología descriptiva.

\paragraph{(e) Límites.} (i)~La normalización \emph{trunca} más allá de $12\,000$
caracteres: dos enmiendas muy largas solo se comparan en su prefijo. (ii)~El
shingling de palabras ignora el orden más allá de la ventana $k$: una reordenación
de párrafos puede hacer caer el Jaccard. (iii)~El umbral $0{,}80$ está en la zona
ambigua de la nube: la frontera exacta de los clústeres es convencional. (iv)~Un
\og{}clúster\fg{} no dice nada de una coordinación intencional: un mismo modelo
puede circular vía una federación de asociaciones, un sindicato, un sitio
militante, sin concertación entre diputados.

\subsection{Tipología de los clústeres}\label{sec:typo-es}

\paragraph{(a) Definición.} Para un clúster $C$, sean $n(C)=|C|$ su tamaño y $g(C)$
el número de grupos parlamentarios distintos entre sus enmiendas. Se clasifica:
\[
\tau(C)=\begin{cases}
\text{presentación masiva} & \text{si } n(C)\ge \param{T_{masa}}=15,\\[2pt]
\text{convergencia intergrupos} & \text{si no, si } g(C)\ge 2,\\[2pt]
\text{amplificación intragrupo} & \text{si no, si } g(C)\le 1 \text{ y } n(C)\ge \param{T_{ampl}}=3,\\[2pt]
\text{reutilización simple} & \text{en otro caso.}
\end{cases}
\]
El tamaño siempre prevalece; no se infiere ninguna intención (la etiqueta interna
\og{}obstrucción\fg{} se neutralizó a \og{}presentación masiva\fg{}).

\paragraph{(b) Algoritmo y parámetros.} Cálculo directo sobre la tabla de clústeres
($\param{T_{masa}}=15$, $\param{T_{ampl}}=3$). En la instantánea (clústeres de
$\theta=0{,}80$): $91$ clústeres \og{}presentación masiva\fg{} ($2\,578$ enmiendas),
$4\,314$ \og{}convergencia\fg{} ($13\,834$), $773$ \og{}amplificación\fg{}
($3\,146$), $4\,283$ \og{}reutilización\fg{} ($8\,566$).

\paragraph{(c) Sensibilidad.}
\begin{center}\small
\begin{tabular}{rr rr rr rr}
\toprule
$\param{T_{masa}}$ & $\param{T_{ampl}}$ & \multicolumn{2}{c}{pres. masiva} & \multicolumn{2}{c}{convergencia} & \multicolumn{2}{c}{amplificación} \\
 & & cl. & enm. & cl. & enm. & cl. & enm. \\
\midrule
10 & 3 & 203 & 3\,875 & 4\,226 & 12\,812 & 749 & 2\,871 \\
\textbf{15} & \textbf{3} & \textbf{91} & \textbf{2\,578} & \textbf{4\,314} & \textbf{13\,834} & \textbf{773} & \textbf{3\,146} \\
20 & 3 & 43 & 1\,776 & 4\,352 & 14\,472 & 783 & 3\,310 \\
30 & 3 & 19 & 1\,218 & 4\,370 & 14\,880 & 789 & 3\,460 \\
50 & 3 & \phantom{0}7 & \phantom{0,}784 & 4\,377 & 15\,121 & 794 & 3\,653 \\
\midrule
15 & 2 & 91 & 2\,578 & 4\,314 & 13\,834 & 5\,056 & 11\,712 \\
15 & 5 & 91 & 2\,578 & 4\,314 & 13\,834 & 160 & 1\,138 \\
\bottomrule
\end{tabular}
\end{center}
El \emph{número} de clústeres \og{}presentación masiva\fg{} es sensible a
$\param{T_{masa}}$ ($91$ a $7$ entre los umbrales $15$ y $50$) pero el
\emph{volumen} de enmiendas afectado lo es menos; el umbral $\param{T_{ampl}}$
sobre todo hace deslizar clústeres entre \og{}amplificación\fg{} y
\og{}reutilización\fg{}. La distribución por tamaño está en el
anexo~\ref{sec:annexes-es}.

\paragraph{(d) Literatura.} La tipología es ad~hoc y descriptiva; se inscribe en el
campo del estudio de la \emph{obstrucción parlamentaria} y de la presentación
masiva de enmiendas como táctica de debate (cf.~Koger~\cite{koger2010} para el caso
estadounidense; el propio Reglamento de la Asamblea Nacional regula la
\og{}discusión común\fg{} de enmiendas idénticas).

\paragraph{(e) Límites.} Cuatro casillas, umbrales redondos: un clúster de $14$
enmiendas idénticas y uno de $16$ caen a uno y otro lado de una frontera sin
fundamento teórico. La casilla \og{}convergencia\fg{} no distingue un consenso
técnico de una alianza puntual.

\subsection{Disciplina de voto (cohesión individual)}\label{sec:discipline-es}

\paragraph{(a) Definición.} Sea $G_d$ el grupo del diputado $d$. Sea $V_d$ el
conjunto de las votaciones donde $d$ ha emitido un voto $p_d(v)\in\{\text{a favor},\text{en contra},\text{abstención}\}$
\emph{y} donde $p^\star(v,G_d)$ está definida. Se pone
\[
\mathrm{emitidos}(d)=|V_d|,\quad
\mathrm{alineados}(d)=\big|\{v\in V_d : p_d(v)=p^\star(v,G_d)\}\big|,\quad
\mathrm{disciplina}(d)=\frac{\mathrm{alineados}(d)}{\mathrm{emitidos}(d)}.
\]
Se muestra en las clasificaciones solo si $\mathrm{emitidos}(d)\ge 50$ (o $\ge 100$
para la lista de \og{}disidentes\fg{}, ordenada por disciplina creciente); los
diputados no inscritos (NI) quedan excluidos.

\paragraph{(b) Algoritmo y parámetros.} Todo se apoya en el campo oficial
$p^\star$: \emph{no} calculamos nosotros la posición mayoritaria de un grupo. En la
instantánea, $\mathrm{disciplina}$ tiene una media de $\approx 91{,}9~\%$, una
mediana de $\approx 92{,}5~\%$, un decil inferior de $\approx 85~\%$.

\paragraph{(c) Sensibilidad.} Efecto del umbral mínimo de votos emitidos:
\begin{center}\small
\begin{tabular}{rrrrrr}
\toprule
umbral & elegibles & media & mediana & $p_{10}$ & $p_{90}$ \\
\midrule
0   & 639 & 91,92 & 92,56 & 84,92 & 98,30 \\
20  & 628 & 91,87 & 92,50 & 84,93 & 98,25 \\
\textbf{50}  & \textbf{623} & \textbf{91,85} & \textbf{92,50} & \textbf{84,99} & \textbf{98,21} \\
100 & 611 & 91,83 & 92,45 & 85,02 & 98,18 \\
200 & 600 & 91,87 & 92,46 & 85,06 & 98,18 \\
\bottomrule
\end{tabular}
\end{center}
El indicador es \emph{muy robusto} al umbral: la media y los cuantiles solo varían
unas centésimas de punto entre los umbrales $0$ y $200$. La elección $\ge 50$ no es,
pues, un parámetro sensible — solo sirve para evitar mostrar un \og{}100~\%\fg{}
basado en tres votos.

\paragraph{(d) El \og{}sesgo circular\fg{}.} La posición $p^\star(v,G_d)$ es la
pluralidad de los votos de los miembros de $G_d$, y $d$ forma parte de él: $d$ es,
pues, comparado, en parte, con una posición que él ha contribuido a definir. Es un
sesgo \emph{real} pero (i)~no es obra nuestra — es la convención oficial de la
Asamblea; (ii)~para un grupo de $50$ a $120$ miembros, el peso de un voto individual
en la pluralidad está acotado por $\le 2~\%$; (iii)~el caso de un empate exacto, que
haría ambigua la pluralidad, lo arbitra aguas arriba la Asamblea y nunca aparece en
los datos (el campo $p^\star$ está siempre cumplimentado y siempre en
$\{\text{a favor},\text{en contra},\text{abstención}\}$, cf.~§\ref{sec:notations-es}).

\paragraph{(e) Literatura.} La disciplina así definida es la versión
\emph{individual} (por diputado) de los índices de cohesión partidaria: índice de
Rice~\cite{rice1925}, \emph{Agreement Index} de Hix-Noury-Roland~\cite{hnr2005,hnr2007},
discusión de Carey~\cite{carey2007} sobre los \og{}principales en competencia\fg{}.
\emph{No} estima puntos ideales: para ello harían falta métodos tipo NOMINATE o
\emph{Optimal Classification} (Poole-Rosenthal~\cite{pr1997}, Poole~\cite{poole2005}),
fuera de alcance.

\paragraph{(e') Límites.} (i)~Perímetro restringido a las votaciones públicas
nominales: la \og{}disciplina\fg{} medida es la de los votos \emph{emitidos y
registrados}, no la unidad de grupo en general. (ii)~Una abstención conforme con la
posición del grupo cuenta como \og{}alineada\fg{}. (iii)~Un diputado a menudo
ausente tiene un $\mathrm{emitidos}$ bajo: su disciplina se mide sobre pocos casos
(de ahí el umbral).

\subsection{Matriz de cohesión entre grupos}\label{sec:matrice-es}

\paragraph{(a) Definición.} Para dos grupos $G_1,G_2$, sea $V_{12}$ el conjunto de
las votaciones donde ambos tienen una posición mayoritaria definida. La cohesión es
\[
\mathrm{coh}(G_1,G_2)=\frac{\big|\{v\in V_{12}:p^\star(v,G_1)=p^\star(v,G_2)\}\big|}{|V_{12}|},
\qquad \mathrm{coh}(G,G)=1.
\]

\paragraph{(b) Algoritmo y parámetros.} Ningún parámetro libre; una auto-unión
sobre la tabla de posiciones mayoritarias.

\paragraph{(c) Sensibilidad.} Una elección discutible: una \og{}abstención $=$
abstención\fg{} cuenta como acuerdo. Recalculando la matriz sin las votaciones donde
ambos grupos se abstienen, la diferencia absoluta media es de $0{,}21$~puntos y el
máximo de $0{,}83$~puntos. El efecto es, pues, \emph{despreciable}.

\paragraph{(d) Literatura.} Es una matriz de \emph{acuerdo} entre bloques en el
sentido de Hix-Noury-Roland~\cite{hnr2007}; no proporciona ningún orden
unidimensional (para ello, véanse los métodos espaciales citados en
§\ref{sec:discipline-es}).

\paragraph{(e) Límites.} Cohesión calculada sobre \emph{posiciones de grupo}, no
sobre votos individuales: un grupo dividido $51/49$ y un grupo unánime cuentan
igual. Dos \og{}en contra\fg{} por motivos opuestos son indistinguibles de un
acuerdo de fondo.

\subsection{Diagrama ternario de los bloques}\label{sec:ternaire-es}

\paragraph{(a) Definición.} Dados tres grupos \emph{ancla} $A_1,A_2,A_3$ situados en
los vértices de un triángulo equilátero en posiciones $\mathbf{x}_{A_1},\mathbf{x}_{A_2},\mathbf{x}_{A_3}$,
un grupo $G$ de cohesiones $c_i=\mathrm{coh}(G,A_i)$ se sitúa en
\[
(\alpha,\beta,\gamma)=\frac{(c_1,c_2,c_3)}{c_1+c_2+c_3},\qquad
\mathbf{x}_G=\alpha\,\mathbf{x}_{A_1}+\beta\,\mathbf{x}_{A_2}+\gamma\,\mathbf{x}_{A_3}.
\]
Las anclas por defecto son LFI, EPR, RN; el usuario puede cambiarlas.

\paragraph{(b)--(c) Algoritmo y sensibilidad.} El único \og{}parámetro\fg{} es la
elección de las tres anclas: el diagrama no es más que una \emph{re-proyección} de
la matriz de cohesión §\ref{sec:matrice-es} sobre tres ejes elegidos, y cambia con
ellos.

\paragraph{(d)--(e) Literatura y límites.} Herramienta de visualización, sin
pretensión de medida: la posición de un grupo \emph{depende del marco de
referencia}. No leer las distancias como distancias ideológicas.

\subsection{Coaliciones por votación (por tema)}\label{sec:coalitions-es}

\paragraph{(a) Definición.} Se seleccionan las votaciones \og{}muy concurridas y
divisivas\fg{}: (a)~$\text{emitidos}(v)\ge \param{T_v}=300$ ($\approx 52~\%$ de los
577 escaños); (b)~$|\mathrm{a\_favor}(v)-\mathrm{en\_contra}(v)|\le \param{T_e}=50$.
Se ordena por (b) creciente y luego fecha decreciente y se conservan las $N=10$
primeras; si pasan menos de $N/2$ votaciones, se relaja el criterio (b). Para cada
votación retenida, se muestra la posición mayoritaria de cada grupo.

\paragraph{(b)--(c) Parámetros y sensibilidad.} Número de votaciones elegibles según
$(\param{T_v},\param{T_e})$:
\begin{center}\small
\begin{tabular}{r|rrrrr}
\toprule
$\param{T_v}\backslash\param{T_e}$ & 10 & 25 & 50 & 75 & 100 \\
\midrule
200 & \phantom{0}95 & 252 & 565 & 913 & 1\,200 \\
250 & \phantom{0}25 & \phantom{0}81 & 193 & 328 & \phantom{0,}440 \\
\textbf{300} & \phantom{0}\textbf{11} & \phantom{0}\textbf{36} & \phantom{00}\textbf{79} & 118 & \phantom{0,}155 \\
350 & \phantom{00}3 & \phantom{0}14 & \phantom{00}24 & \phantom{0}47 & \phantom{0,0}55 \\
400 & \phantom{00}2 & \phantom{00}7 & \phantom{00}11 & \phantom{0}19 & \phantom{0,0}20 \\
\bottomrule
\end{tabular}
\end{center}
En el ajuste de producción $(300,50)$, $79$ votaciones pasan: el repliegue (b) no se
activa. La selección es sensible a ambos umbrales, como cabía esperar (es un filtro,
no un estimador).

\paragraph{(d)--(e) Literatura y límites.} Enfoque puramente descriptivo: se muestra
\emph{quién vota como quién} en las votaciones ajustadas, sin modelo. Límite
principal: una votación \og{}ajustada en votos\fg{} puede serlo por baja
concurrencia asimétrica tanto como por una verdadera fractura.

\subsection{Respuestas ministeriales \og{}tipo\fg{}}\label{sec:templates-es}

\paragraph{(a) Definición.} Para cada respuesta escrita $q$ cuyo texto en bruto
supera los $200$ caracteres, se normaliza (supresión HTML, espacios), se retira
eventualmente la fórmula de cortesía inicial, se pasa a minúsculas: $t_q$. La
respuesta se descarta si $|t_q|<100$. La \emph{clave de prefijo} es
$\pi_L(q)=t_q[1:L]$ con $L=250$. Una \og{}respuesta tipo\fg{} es un prefijo
compartido por al menos $m=3$ respuestas distintas. La \emph{tasa de respuestas
tipo} es
\[
\rho \;=\; \frac{1}{|Q_{\text{eleg}}|}\sum_{\pi \,:\, |\{q:\pi_L(q)=\pi\}|\ge m} \big|\{q:\pi_L(q)=\pi\}\big|.
\]
En la instantánea, $|Q_{\text{eleg}}|=10\,766$ y $\rho=27{,}9~\%$.

\paragraph{(b)--(c) Parámetros y sensibilidad.} Rejilla de la tasa $\rho$ (en \%)
según la longitud de prefijo $L$ y el umbral de ocurrencias $m$:
\begin{center}\small
\begin{tabular}{r|rrrr}
\toprule
$L\backslash m$ & 2 & \textbf{3} & 5 & 10 \\
\midrule
100  & 42,1 & 31,6 & 23,1 & 13,4 \\
\textbf{250}  & 37,0 & \textbf{27,9} & 20,3 & 11,6 \\
500  & 33,7 & 24,6 & 17,2 & \phantom{0}9,9 \\
1000 & 30,2 & 22,1 & 15,3 & \phantom{0}9,0 \\
\bottomrule
\end{tabular}
\end{center}
La tasa es \emph{fuertemente sensible al umbral $m$} (cociente $\approx 3$--$4$
entre $m=2$ y $m=10$) y \emph{moderadamente sensible a $L$} (de $\approx 32~\%$ a
$\approx 22~\%$ con $m$ fijo). En este corpus, retirar la fórmula de cortesía no
tiene \emph{ningún} efecto (las variantes \og{}con\fg{} y \og{}sin\fg{} retirada
coinciden exactamente, el texto almacenado no empieza por esas fórmulas). El valor
mostrado en el sitio debe, pues, leerse como \emph{un} punto de una familia de
medidas, no como una constante: la formulación honesta es \og{}$\approx 28~\%$ de
las respuestas comparten un prefijo de $250$~caracteres con al menos otras dos\fg{}.

\paragraph{(d) Literatura.} Las preguntas parlamentarias como instrumento de
control, y sus respuestas como objeto de estudio, son tratadas por
Wiberg~\cite{wiberg1995}, Russo-Wiberg~\cite{rw2010}, Martin~\cite{martin2011}. La
detección de respuestas estandarizadas por agrupamiento de prefijos es una
heurística \emph{casera}, próxima a la detección de \og{}boilerplate\fg{}; no es una
medida validada de \og{}lenguaje de madera\fg{}.

\paragraph{(e) Límites.} (i)~Un prefijo común es una condición \emph{necesaria} pero
no \emph{suficiente} de una respuesta \og{}vacía\fg{}: dos respuestas pueden
compartir un encabezado de encuadre y luego divergir por completo. (ii)~El umbral
$m=3$ y la longitud $L=250$ son arbitrarios (cf.~rejilla). (iii)~El corpus se limita
a las respuestas publicadas de más de $200$ caracteres; las preguntas sin respuesta
no entran en el denominador de $\rho$.

\subsection{Absentismo estratégico}\label{sec:absence-es}

\paragraph{(a) Definición.} Entre las votaciones con $\text{emitidos}(v)\ge \param{T_v}=200$,
se dice que $v$ es \emph{divisiva} si
$|\mathrm{a\_favor}(v)-\mathrm{en\_contra}(v)|/\text{emitidos}(v)<\param{c}=0{,}10$,
\emph{consensual} si $>\param{C}=0{,}50$, en otro caso sin clasificar. Para un
diputado $d$, sea $\mathrm{pres}_{div}(d)$ (resp.\ $\mathrm{pres}_{con}(d)$) la
\emph{proporción} de votaciones divisivas (resp.\ consensuales) en las que $d$ ha
emitido un voto. Se pone $\mathrm{gap}(d)=\mathrm{pres}_{con}(d)-\mathrm{pres}_{div}(d)$
(en puntos) y $\mathrm{ratio}(d)=\mathrm{pres}_{con}(d)/\mathrm{pres}_{div}(d)$. Se
muestra si $d$ está en ejercicio y $\mathrm{tot}_{div}(d)\ge \param{T_{div}}=30$.

\paragraph{(b)--(c) Parámetros y sensibilidad.} Cuatro umbrales
($\param{T_v},\param{c},\param{C},\param{T_{div}}$). Los valores exactos de las
$4^4$ combinaciones exploradas figuran en el anexo~\ref{sec:annexes-es}; se observa
que la mediana del \emph{gap} está cerca de $0$ — la mayoría de los diputados tienen
una presencia similar en votaciones divisivas y consensuales — y que solo los
\emph{outliers} tienen un \emph{gap} marcado, cuya amplitud varía de $\approx 13$ a
$\approx 31$ puntos según los umbrales.

\paragraph{(d)--(e) Literatura y límites.} La idea se relaciona con la literatura
sobre el \emph{shirking} legislativo y la participación selectiva en los votos.
Límites fuertes: (i)~no votar $\ne$ estar ausente — un diputado presente puede
abstenerse de participar; (ii)~\og{}divisiva\fg{} se define por el \emph{resultado}
del voto, no por su relevancia política \emph{a priori}; (iii)~un \emph{gap}
positivo tiene muchas causas no estratégicas (agenda, enfermedad, circunscripción
lejana). Una señal a explorar, no un veredicto.

\subsection{Enmiendas \og{}fantasma\fg{}}\label{sec:fantomes-es}

\paragraph{(a) Definición.} Para el primer firmante $d$, sobre el conjunto de sus
enmiendas:
\[
\begin{aligned}
\mathrm{defendidas}(d) &= \big|\{a : \mathrm{sort}(a)\in\{\text{Aprobada, Rechazada, Retirada, Discutida}\}\}\big|,\\[2pt]
\mathrm{fantasma}(d) &= \big|\{a : \mathrm{sort}(a)=\text{No defendida, o vacío}\}\big|,\\[2pt]
\mathrm{pct\_fantasma}(d) &= \mathrm{fantasma}(d)\,/\,\mathrm{total}(d).
\end{aligned}
\]
Se muestra si $\mathrm{total}(d)\ge \param{T_a}=50$; ordenado por
$\mathrm{pct\_fantasma}$ decreciente.

\paragraph{(b)--(c) Parámetros y sensibilidad.} Censo de los códigos \og{}sort\fg{}
(porcentaje del total de enmiendas): Rechazada 22,9~\%; Inadmisible 16,3~\%; En
tramitación 14,1~\%; Aprobada 13,0~\%; Retirada 10,2~\%; Discutida 9,0~\%; Decaída
8,7~\%; No defendida 5,9~\%; Borrada $\approx$0~\%. (Ninguna enmienda tiene un sort
vacío o nulo en la instantánea.) Por tanto $\mathrm{fantasma}$ se reduce en la
práctica a \og{}No defendida\fg{} ($5{,}9~\%$), $\mathrm{defendidas}$ cubre
$\approx 55~\%$, y el resto (Inadmisible, En tramitación, Decaída) no se cuenta
\emph{ni} como defendida \emph{ni} como fantasma. Efecto del umbral $\param{T_a}$:
\begin{center}\small
\begin{tabular}{rrrrr}
\toprule
$\param{T_a}$ & autores elegibles & de los cuales en ejercicio & $\overline{\mathrm{pct\_fant.}}$ & $\mathrm{pct\_fant.}^{\max}$ \\
\midrule
10  & 558 & 513 & 7,4 & 61,5 \\
20  & 530 & 491 & 7,2 & 52,4 \\
\textbf{50}  & \textbf{439} & \textbf{412} & \textbf{6,5} & \textbf{50,0} \\
100 & 313 & 300 & 6,6 & 50,0 \\
200 & 167 & 163 & 5,8 & 34,7 \\
\bottomrule
\end{tabular}
\end{center}
La media de la tasa es estable ($\approx 6$--$7{,}5~\%$); el \emph{máximo} decrece
con el umbral (los valores extremos los llevan autores de bajo volumen). Una
\emph{vista continua} (diagrama de dispersión $\mathrm{total}\times\mathrm{pct\_fantasma}$,
sin umbral) se proporciona en el depósito para no ocultar esta dependencia.

\paragraph{(d)--(e) Literatura y límites.} Sin literatura directa; cuantificamos una
asimetría entre \emph{presentación} y \emph{defensa} de enmiendas. Límites: (i)~la
categoría \og{}No defendida\fg{} también cubre restricciones de tiempo en el pleno,
independientes de la voluntad del autor; (ii)~la elección de incluir
\og{}Discutida\fg{} en \og{}defendida\fg{} y de excluir \og{}Decaída\fg{} es
discutible (la rejilla de sensibilidad sobre la definición de \og{}defendida\fg{}
está en el depósito); (iii)~\og{}primer firmante\fg{} $\ne$ \og{}autor real\fg{} en
el caso de enmiendas colectivas.

\subsection{Plazos de respuesta ministeriales y clasificaciones}\label{sec:delais-es}

\paragraph{(a) Definición.} Para una pregunta $q$, $\mathrm{plazo}(q)$ es el número
de días entre la presentación y la respuesta (indefinido si no se responde). Para un
ministerio (en el sentido del \og{}ministerio interpelado\fg{} corto) $\mu$:
\[
\overline{\mathrm{plazo}}(\mu)=\operatorname*{media}_{q:\,\mathrm{plazo}(q)\ \text{definido}}\mathrm{plazo}(q),\qquad
\mathrm{tasa\_resp}(\mu)=\frac{|\{q:\text{estado}=\text{con respuesta}\}|}{|\{q\ \text{dirigidas a }\mu\}|}.
\]
Se muestra si $|\{q\}|\ge \param{T_q}=30$ \emph{y} $|\{q:\text{estado}=\text{con respuesta}\}|\ge \param{T_r}=5$.
La tasa de respuesta se muestra junto al plazo: un plazo corto no significa lo mismo
si solo $1$ pregunta de cada $10$ ha obtenido respuesta.

\paragraph{(b)--(c) Parámetros y sensibilidad.} Número de ministerios retenidos
según $(\param{T_q},\param{T_r})$: $(10,1)\to 73$; $(20,*)\to 62$; $(30,*)\to 57$;
$(50,*)\to 52$; $(100,*)\to 40$. El umbral $\param{T_r}$ es casi inoperante (un
ministerio que recibe $\ge 30$ preguntas casi siempre ha respondido $\ge 5$).

\paragraph{(d)--(e) Literatura y límites.} Cf.~§\ref{sec:templates-es} para las
preguntas como objeto de estudio. Límites: la media es sensible a los valores
extremos (una respuesta a $700$~días arrastra la media); el \og{}ministerio
interpelado\fg{} puede diferir del \og{}ministerio asignado\fg{} que efectivamente
responde.

\subsection{Indicadores derivados simples}\label{sec:simples-es}

\paragraph{(a) Actividad acumulada.} $\mathrm{act}(d)=n_Q(d)+n_A(d)+n_V(d)$
(preguntas formuladas $+$ enmiendas presentadas $+$ votos emitidos).
\emph{Advertencia estructural}: $n_V$ es de un orden de magnitud superior a $n_Q$ y
$n_A$ (un diputado asiduo emite miles de votos pero presenta decenas de enmiendas y
formula decenas de preguntas); el índice acumulado está, pues, \emph{mecánicamente
dominado por la presencia en las votaciones}. El sitio lo acompaña sistemáticamente
de las subclasificaciones aisladas (solo preguntas, solo enmiendas, solo presencia).

\paragraph{(b) Tasa de aprobación de las enmiendas.} $\mathrm{aprobadas}(d)/\mathrm{total}(d)$
(por diputado o por grupo).

\paragraph{(c) Tasa de respuesta a las preguntas.} cf.~§\ref{sec:delais-es}.

\paragraph{(d) Edad, antigüedad.} Derivadas de la fecha de nacimiento y de la fecha
de inicio del mandato. Triviales.

\paragraph{(e) Constructor de rankings parametrizables.} El sitio ofrece una
treintena de \og{}métricas\fg{} combinables con filtros opcionales (por grupo,
período, tipo de texto, etc.). Todas son \emph{re-exposiciones} de las primitivas
anteriores (recuentos, ratios, medias); no hay ningún método adicional.

\paragraph{(f) Estadísticas por circunscripción.} $\mathrm{participación}=\mathrm{votantes}/\mathrm{inscritos}$,
población (INSEE), inscritos (Ministerio del Interior), unidos por
$(\text{código de departamento},\ \text{número de circunscripción})$. Cobertura en
la instantánea: $566$ circunscripciones con población INSEE (las $11$
circunscripciones de los franceses en el extranjero no la tienen), $577$ con
inscritos.

\section{Reproducibilidad}\label{sec:repro-es}
\begin{itemize}
  \item \textbf{Código de la aplicación}: público, con licencia MIT
        (\href{https://github.com/RogericBot/577deputes}{github.com/RogericBot/577deputes}).
        La base de datos y las fotografías no se redistribuyen; se reconstruyen con
        el comando \code{anqp bootstrap} (descarga e ingestión de las fuentes).
  \item \textbf{Cachés derivados}: las tablas \code{deputy\_discipline\_cache},
        \code{groupe\_discipline\_cache}, \code{deputy\_amd\_cache},
        \code{groupe\_amd\_cache}, \code{scrutin\_groupes}, \code{amendement\_clusters},
        \code{circo\_stats} se recalculan en la ingestión; no contienen nada que no
        derive de las fuentes.
  \item \textbf{Análisis de sensibilidad}: el script \code{scripts/methodo\_sensitivity.py}
        (solo biblioteca estándar, solo lectura) reproduce todas las tablas de este
        documento; está adjunto al depósito Zenodo con sus salidas (CSV y Markdown).
        Todos los sorteos aleatorios usan una semilla fija.
  \item \textbf{Instantánea}: las cifras de este documento corresponden a una copia
        de la base tomada para su redacción (volumen en §1.3); pueden diferir
        ligeramente del estado actual del sitio, que se actualiza a diario.
\end{itemize}

\section{Límites generales y sesgos conocidos}\label{sec:limites-es}
\begin{enumerate}[label=L\arabic*.]
  \item \textbf{Perímetro de las votaciones.} Solo existen en los datos las
        votaciones públicas nominales celebradas en el pleno. Las votaciones a mano
        alzada y las votaciones en comisión no figuran. Todo indicador de voto
        (disciplina, cohesión, absentismo, coaliciones) hereda este límite. Además,
        la Asamblea vota muy a menudo con solo una fracción del hemiciclo presente
        (mediana $\approx 150$ votantes de $577$): un recuento bajo refleja la
        concurrencia de ese día, no un voto truncado, y no se exige quórum para la
        mayoría de los votos.
  \item \textbf{\og{}Autor\fg{} de una enmienda.} El campo usado es el \emph{primer
        firmante}; para una enmienda colectiva no refleja necesariamente la
        redacción.
  \item \textbf{Legislaturas a caballo.} Algunos expedientes abiertos bajo la
        XVI\textsuperscript{a} legislatura y continuados bajo la XVII\textsuperscript{a}
        están asignados a su legislatura de origen; de ahí pequeñas discrepancias
        entre contadores según el perímetro elegido.
  \item \textbf{Carácter heurístico de las tipologías.} \og{}Presentación masiva\fg{},
        \og{}convergencia\fg{}, \og{}amplificación\fg{}, \og{}reutilización\fg{},
        \og{}respuesta tipo\fg{} son etiquetas definidas por umbrales redondos, no
        categorías teóricas.
  \item \textbf{Ninguna inferencia de intención.} Los indicadores describen
        regularidades medibles. No dicen \emph{por qué}: una presentación masiva
        puede enriquecer el debate o ralentizarlo; un \emph{gap} de presencia tiene
        a menudo causas prosaicas; una respuesta con prefijo común puede ser
        sustancial tras ese prefijo.
  \item \textbf{Sin estimación de puntos ideales.} Ningún indicador produce una
        escala izquierda-derecha ni un mapa espacial; la matriz de cohesión y el
        diagrama ternario son solo acuerdos por pares re-proyectados.
  \item \textbf{Dependencia de los productores de datos.} Una laguna, un retraso o
        un cambio de formato en la Asamblea, el INSEE o el Ministerio del Interior se
        repercute directamente; el sitio señala explícitamente cuándo una fuente no
        está disponible en lugar de inventar.
  \item \textbf{Instantánea.} El sitio se actualiza a diario; toda cifra citada está
        fechada en el día de su cálculo.
\end{enumerate}

\section{Anexos}\label{sec:annexes-es}
Las tablas completas siguientes se proporcionan también en CSV/Markdown en el
depósito (un archivo por tabla, más \code{SUMMARY.md}).
\begin{itemize}\small
  \item \textbf{A1} umbral Jaccard sobre el conjunto completo; \textbf{A1bis}
        histograma de las similitudes candidatas; \textbf{A2} $k$ y $\theta$ sobre
        muestra (Jaccard exacto); \textbf{A3} curva LSH $P(s)=1-(1-s^r)^b$;
        \textbf{A4} exposición a $L_{\max}$; \textbf{A4bis} umbral $s_{\min}$.
  \item \textbf{B} umbrales de tipología; \textbf{B2} distribución de los tamaños de
        clúster.
  \item \textbf{C0} valores de \code{positionMajoritaire}; \textbf{C1} umbral de
        disciplina.
  \item \textbf{D} matriz de cohesión, efecto \og{}abstención compartida\fg{}.
  \item \textbf{F} umbrales \og{}coaliciones por votación\fg{}.
  \item \textbf{G} rejilla \og{}respuestas tipo\fg{} ($L\times m\times$ retirada de
        las fórmulas).
  \item \textbf{H} rejilla \og{}absentismo estratégico\fg{} ($\param{T_v}\times\param{c}\times\param{C}\times\param{T_{div}}$).
  \item \textbf{I0} censo de los sorts; \textbf{I1} umbral fantasma; \textbf{I2}
        vista continua (diagrama de dispersión).
  \item \textbf{J} umbrales del filtro de ministerios; \textbf{K} cobertura de las
        estadísticas de circunscripción.
\end{itemize}
   % Partie III — Versión en español (esqueleto en esta etapa)
%==============================================================================

%------------------------------------------------------------------------------
% Bibliographie (manuelle — pas de BibTeX requis sur Overleaf)
%------------------------------------------------------------------------------
\selectlanguage{french}
\begin{thebibliography}{99}\small
\bibitem{rice1925} S.~A. Rice (1925). \emph{The behavior of legislative groups: a method of measurement.} Political Science Quarterly 40(1), 60--72.
\bibitem{hnr2005} S.~Hix, A.~Noury, G.~Roland (2005). \emph{Power to the parties: cohesion and competition in the European Parliament, 1979--2001.} British Journal of Political Science 35(2), 209--234.
\bibitem{hnr2007} S.~Hix, A.~Noury, G.~Roland (2007). \emph{Democratic Politics in the European Parliament.} Cambridge University Press.
\bibitem{carey2007} J.~M. Carey (2007). \emph{Competing principals, political institutions, and party unity in legislative voting.} American Journal of Political Science 51(1), 92--107.
\bibitem{pr1997} K.~T. Poole, H.~Rosenthal (1997). \emph{Congress: A Political-Economic History of Roll Call Voting.} Oxford University Press.
\bibitem{poole2005} K.~T. Poole (2005). \emph{Spatial Models of Parliamentary Voting.} Cambridge University Press.
\bibitem{broder1997} A.~Z. Broder (1997). \emph{On the resemblance and containment of documents.} Proc. Compression and Complexity of Sequences, 21--29.
\bibitem{broder2000} A.~Z. Broder, M.~Charikar, A.~M. Frieze, M.~Mitzenmacher (2000). \emph{Min-wise independent permutations.} Journal of Computer and System Sciences 60(3), 630--659.
\bibitem{gionis1999} A.~Gionis, P.~Indyk, R.~Motwani (1999). \emph{Similarity search in high dimensions via hashing.} Proc. VLDB, 518--529.
\bibitem{lru2014} J.~Leskovec, A.~Rajaraman, J.~D. Ullman (2014/2020). \emph{Mining of Massive Datasets}, ch.~3 (Finding Similar Items). Cambridge University Press.
\bibitem{burgess2016} M.~Burgess, E.~Giraudy, J.~Katz-Samuels, J.~Walsh, D.~Willis, L.~Haynes, R.~Ghani (2016). \emph{The Legislative Influence Detector: finding text reuse in state legislation.} Proc. ACM KDD, 57--66.
\bibitem{hf2014} A.~Hertel-Fernandez (2014). \emph{Who passes business's ``model bills''? Policy capacity and corporate influence in U.S. state politics.} Perspectives on Politics 12(3), 582--602.
\bibitem{jansa2019} J.~M. Jansa, E.~R. Hansen, V.~H. Gray (2019). \emph{Copy and paste lawmaking: legislative professionalism and policy reinvention in the states.} American Politics Research 47(4), 739--767.
\bibitem{linder2020} F.~Linder, B.~Desmarais, M.~Burgess, E.~Giraudy (2020). \emph{Text as policy: measuring policy similarity through bill text reuse.} Policy Studies Journal 48(2), 546--574.
\bibitem{wiberg1995} M.~Wiberg (1995). \emph{Parliamentary questioning: control by communication?} In H.~Döring (ed.), Parliaments and Majority Rule in Western Europe, 179--222.
\bibitem{rw2010} F.~Russo, M.~Wiberg (2010). \emph{Parliamentary questioning in 17 European parliaments: some steps towards comparison.} Journal of Legislative Studies 16(2), 215--232.
\bibitem{martin2011} S.~Martin (2011). \emph{Parliamentary questions, the behaviour of legislators, and the function of legislatures: an introduction.} Journal of Legislative Studies 17(3), 259--270.
\bibitem{koger2010} G.~Koger (2010). \emph{Filibustering: A Political History of Obstruction in the House and Senate.} University of Chicago Press.
\end{thebibliography}

\end{document}
